% concatenated from Mazur-paper_revised.tex
\documentclass[a4paper,10pt]{amsart}
\usepackage[utf8]{inputenc}
\usepackage{amsmath}
\usepackage{amsfonts}
\usepackage{amssymb}
\usepackage{graphicx}
\usepackage[left=2.4cm,right=2.4cm,top=2.5cm,bottom=2.5cm]{geometry}
\numberwithin{equation}{section}
\graphicspath{ {images/} }
\usepackage{amsmath,amssymb, amsthm} 
\usepackage{lipsum}

\usepackage{amsmath}
\usepackage{amsfonts}
\usepackage{amssymb}
\usepackage{graphicx}
%\usepackage[left=2.4cm,right=2.4cm,top=2.5cm,bottom=2.5cm]{geometry}
\numberwithin{equation}{section}
\graphicspath{ {images/} }
\usepackage{amsmath,amssymb, amsthm} 
\usepackage{lipsum}
\usepackage{stmaryrd}
\usepackage{dsfont}
\usepackage[all]{xy}

\makeatletter
\def\@settitle{\begin{center}%
  \baselineskip14\p@\relax
  \bfseries
  \uppercasenonmath\@title
  \@title
  \ifx\@subtitle\@empty\else
     \\[1ex]\uppercasenonmath\@subtitle
     \footnotesize\mdseries\@subtitle
  \fi
  \end{center}%
}
\def\subtitle#1{\gdef\@subtitle{#1}}
\def\@subtitle{}
\makeatother

\newcommand\blfootnote[1]{%
  \begingroup
  \renewcommand\thefootnote{}\footnote{#1}%
  \addtocounter{footnote}{-1}%
  \endgroup
}


\usepackage{tikz-cd}
\usepackage{hyperref}
\usepackage{bm}
\hypersetup{colorlinks, linkcolor=blue, citecolor=black, urlcolor=black}

\makeatletter
\def\@settitle{\begin{center}%
  \baselineskip14\p@\relax
  \bfseries
  \uppercasenonmath\@title
  \@title
  \ifx\@subtitle\@empty\else
     \\[1ex]\uppercasenonmath\@subtitle
     \footnotesize\mdseries\@subtitle
  \fi
  \end{center}%
}
\def\subtitle#1{\gdef\@subtitle{#1}}
\def\@subtitle{}
\makeatother



\usepackage{tikz-cd}
\usepackage{hyperref}
\usepackage{bm}
\hypersetup{colorlinks, linkcolor=blue, citecolor=black, urlcolor=black}

\usepackage[english]{babel}
\usepackage[colorinlistoftodos]{todonotes}

\usepackage[english]{babel}
\usepackage[colorinlistoftodos]{todonotes}

\theoremstyle{plain}
\newtheorem{thm}{Theorem}[subsection] % reset theorem numbering for each chapter

\usepackage{mathrsfs}

\theoremstyle{definition}
\newtheorem{defi}[thm]{Definition}
\newtheorem{conjecture}[thm]{Conjecture}
\newtheorem{defiprop}[thm]{Definition/Proposition}
\newtheorem{rmk}[thm]{Remark}
\newtheorem*{claim}{Claim}
\newtheorem*{fact}{Fact}
\newtheorem*{hyp}{Heegner hypothesis}
\newtheorem{rem}[thm]{Remark}
\theoremstyle{definition}
\newtheorem{ex}[thm]{Example}
\theoremstyle{plain}
\newtheorem{prop}[thm]{Proposition}
\theoremstyle{plain}
\newtheorem{lemma}[thm]{Lemma}
\theoremstyle{plain}
\newtheorem{cor}[thm]{Corollary}
\theoremstyle{plain}
\newtheorem{thmintro}{Theorem}
\renewcommand{\thethmintro}{\Alph{thmintro}}
\newtheorem{conj}[thmintro]{Conjecture}


\makeatletter
\newenvironment{subtheorem}[1]{%
  \def\subtheoremcounter{#1}%
  \refstepcounter{#1}%
  \protected@edef\theparentnumber{\csname the#1\endcsname}%
  \setcounter{parentnumber}{\value{#1}}%
  \setcounter{#1}{0}%
  \expandafter\def\csname the#1\endcsname{\theparentnumber \alph{#1}}%
  \ignorespaces
}{%
  \setcounter{\subtheoremcounter}{\value{parentnumber}}%
  \ignorespacesafterend
}
\makeatother
\newcounter{parentnumber}

\DeclareMathOperator{\ord}{ord}

\newcommand{\Sp}{\operatorname{Spec}}
\newcommand{\proj}{\operatorname{Proj}}
\newcommand{\mor}{\operatorname{Mor}}
\newcommand{\sym}{\operatorname{Sym}}
\newcommand{\dg}{\operatorname{deg}}
\newcommand{\Hom}{\operatorname{Hom}}
\newcommand{\Frac}{\operatorname{Frac}}
%\newcommand{\M}{\operatorname{Max}}
\newcommand{\Der}{\operatorname{Der}}
\newcommand{\CH}{\operatorname{CH}}
%\newcommand{\AF}{\operatorname{AF}}
\newcommand{\trAF}{\widetilde{\AF}}
%\newcommand{\BF}{\operatorname{BF}}
\newcommand{\res}{\operatorname{Res}}
\newcommand{\ind}{\operatorname{ind}}
\newcommand{\Gal}{\operatorname{Gal}}
\newcommand{\msum}{\sum_{a \in (\Z/M\Z)^{\times}}}
\newcommand{\mC}{\C[(\Z/M\Z)^{\times}]}
\newcommand{\mZ}{(\Z/M\Z)^{\times}}
\newcommand{\dm}{\mathcal{DM}_{-}(k)}
\newcommand{\mof}{\Oo_F/(M\Oo_F+\Z)}
\newcommand{\iio}{\iota_{M,\N,b^{-1}a}}
\newcommand{\utau}{\underline{\tau}}

\newcommand{\bQ}{\mathbb{Q}}

\newcommand\SmallMatrix[1]{{%
  \tiny\arraycolsep=0.3\arraycolsep\ensuremath{\begin{pmatrix}#1\end{pmatrix}}}}
  

\makeatletter  
\newcommand{\colim@}[2]{%
  \vtop{\m@th\ialign{##\cr
    \hfil$#1\operator@font colim$\hfil\cr
    \noalign{\nointerlineskip\kern1.5\ex@}#2\cr
    \noalign{\nointerlineskip\kern-\ex@}\cr}}%
}
\newcommand{\colim}{%
  \mathop{\mathpalette\colim@{\rightarrowfill@\scriptscriptstyle}}\nmlimits@
}
\renewcommand{\varprojlim}{%
  \mathop{\mathpalette\varlim@{\leftarrowfill@\scriptscriptstyle}}\nmlimits@
}
\renewcommand{\varinjlim}{%
  \mathop{\mathpalette\varlim@{\rightarrowfill@\scriptscriptstyle}}\nmlimits@
}
\makeatother  

\newcommand{\limm}{\varprojlim}


\newcommand{\Set}{\left( \operatorname{Set} \right)}
\newcommand{\Top}{\left( \operatorname{Top} \right)}
\newcommand{\SchS}{\left( \operatorname{Sch} / S \right)}
\newcommand{\SchX}{\left( \operatorname{Sch} / X \right)}


\newcommand{\El}{E_0}


\newcommand{\N}{\mathfrak{N}}
\newcommand{\M}{\mathfrak{M}}
\newcommand{\MM}{\mathbb{M}}
\newcommand{\zita}{\mathfrak{z}}
\newcommand{\Zita}{\mathfrak{Z}}
\newcommand{\W}{\mathcal{W}}
\newcommand{\Z}{\mathbb{Z}}
\newcommand{\Q}{\mathbb{Q}}
\newcommand{\Qi}{\mathbb{Q}(i)}
\newcommand{\Qp}{\mathbb{Q}_{p}}
\newcommand{\Zp}{\mathbb{Z}_{p}}
\newcommand{\Fp}{\mathbb{F}_{p}}
\newcommand{\Ff}{\mathbb{F}}
\newcommand{\K}{\mathbb{K}}
\newcommand{\R}{\mathbb{R}}
\newcommand{\C}{\mathbb{C}}
\newcommand{\A}{\mathbb{A}}
\newcommand{\E}{\mathcal{E}}
\newcommand{\Oo}{\mathcal{O}}
\newcommand{\J}{\mathcal{J}}
\newcommand{\I}{\mathcal{I}}
\newcommand{\Hi}{\mathbb{H}}
\newcommand{\hh}{\mathcal{H}}
\newcommand{\B}{\mathcal{B}}
\newcommand{\D}{\mathcal{D}}
\newcommand{\Ss}{\mathcal{S}}
\newcommand{\F}{\mathcal{F}}
\newcommand{\T}{\mathcal{T}}
\newcommand{\p}{\mathfrak{p}}
\newcommand{\cc}{\mathfrak{c}}
\newcommand{\q}{\mathfrak{q}}
\newcommand{\tmap}{\mathbf{t}}
\newcommand{\umap}{\mathbf{u}}
\newcommand{\m}{\mathfrak{m}}
\newcommand{\X}{\mathfrak{X}}
\newcommand{\kk}{\underline{k}}
\newcommand{\ideal}{\mathfrak{a}}
\newcommand{\Af}{A_{f}}
\newcommand{\nAF}{\mathcal{AF}}
\newcommand{\tAF}{\widetilde{\nAF}}
\newcommand{\mat}{\SmallMatrix{a & b \\ c & d}}
\newcommand{\blam}{\bar{\lambda}}
\newcommand{\OG}{\Omega^{(\kappa_1,\kappa_2)}}
\newcommand{\et}{\rm \acute{e}t}

\newcommand{\bg}{\boldsymbol{g}}

\newcommand{\sk}{\vspace{0.1in}}

\newcommand{\isoarrow}{\stackrel{~}{\rightarrow}}

\font\wncyr=wncyr9.8
\newcommand{\sha}{\text{\wncyr{W}}}
\newcommand{\CE}{\mathcal{E}}
\newcommand{\CM}{\mathcal{M}}

\DeclareMathOperator{\Area}{Area}
\DeclareMathOperator{\std}{Std}
\DeclareMathOperator{\Ext}{Ext}
\DeclareMathOperator{\SL}{SL}
\DeclareMathOperator{\ch}{ch}
\DeclareMathOperator{\Las}{L_{As}}
\DeclareMathOperator{\Lasprim}{L^{prim}_{As}}
\DeclareMathOperator{\Lambdaas}{\Lambda_{As}}
\DeclareMathOperator{\as}{As}
\DeclareMathOperator{\GL}{GL}
\DeclareMathOperator{\frob}{Frob}
\DeclareMathOperator{\gal}{Gal}
\DeclareMathOperator{\Gr}{Gr}
\DeclareMathOperator{\mot}{mot}
\DeclareMathOperator{\RG}{R\Gamma}
\DeclareMathOperator{\tor}{tor}
\DeclareMathOperator{\eis}{Eis}
\DeclareMathOperator{\End}{End}
%\DeclareMathOperator{\et}{\text{\'{e}}t}
\DeclareMathOperator{\fil}{Fil}
\DeclareMathOperator{\dr}{dR}
\DeclareMathOperator{\sing}{sing}
\DeclareMathOperator{\betti}{B}
\DeclareMathOperator{\tr}{Tr}
\DeclareMathOperator{\divi}{div}
\DeclareMathOperator{\reg}{reg}
\DeclareMathOperator{\sgn}{sgn}
\DeclareMathOperator{\tsym}{TSym}
\DeclareMathOperator{\del}{\mathcal{D}}
\DeclareMathOperator{\rhoas}{\rho^{As}}
\DeclareMathOperator{\eN}{\mathcal{E}_N(\C)}
\DeclareMathOperator{\re}{Re}
\DeclareMathOperator{\loc}{loc}
\DeclareMathOperator{\id}{Id}
\DeclareMathOperator{\HT}{HT}
\DeclareMathOperator{\nm}{Nm}
\DeclareMathOperator{\cores}{cores}
\DeclareMathOperator{\stab}{Stab}
\DeclareMathOperator{\vol}{Vol}
\DeclareMathOperator{\im}{Im}
\DeclareMathOperator{\Aut}{Aut}
\DeclareMathOperator{\length}{length}


\newcommand{\Lcal}{\mathcal{L}}
\setcounter{tocdepth}{2}

\newcommand{\cF}{\mathcal{F}}
\newcommand{\CF}{\mathcal{F}}
\newcommand{\Fcal}{\mathcal{F}}
\newcommand{\Gcal}{\mathcal{G}}
\newcommand{\cL}{\mathscr{L}}
\newcommand{\cN}{\mathscr{N}}
\newcommand{\kap}{\kappa}
\newcommand{\rank}{\mathrm{rank}}
\newcommand{\fm}{\mathfrak{m}}
\newcommand{\rH}{\mathrm{H}}



\begin{document}
%\title{Iwasawa main conjectures for elliptic curves at Eisenstein primes}
\title{Mazur's main conjecture %for rational elliptic curves 
%and Perrin-Riou's main conjecture 
at Eisenstein primes}\blfootnote{\emph{2020 Mathematics Subject Classification}. 11R23, 11G05, 11G40.}
\author{Francesc Castella}
\address[F.~Castella]{University of California Santa Barbara, South Hall, Santa Barbara, CA 93106, USA}
\email{castella@ucsb.edu}

\author{Giada Grossi}
\address[G.~Grossi]{CNRS, Institut Galilée, Université Sorbonne Paris Nord, 93430 Villetaneuse, FRANCE}
\email{grossi@math.univ-paris13.fr}

 
\author{Christopher Skinner} 
\address[C.~Skinner]{Princeton University, Fine Hall, Washington Road, Princeton, NJ 08544-1000, USA}
\email{cmcls@princeton.edu}


\begin{abstract}
Let $E/\Q$ be an elliptic curve and let $p$ be an odd prime of good reduction for $E$. Assume that $E$ admits a rational $p$-isogeny $\varphi:E\rightarrow E'$, and let $\phi:G_{\bQ}\rightarrow\mathbb{F}_p^\times$ be the character by which $G_{\bQ}$ acts on ${\rm ker}(\varphi)$. In this paper, we prove the Iwasawa main conjecture for $E$, as formulated by B.\,Mazur in 1972, when $\phi\vert_{G_p}\neq 1,\omega$, where $G_p\subset G_{\bQ}$ is a decomposition group at $p$ and $\omega$ is the Teichm\"uller character. 

Two key innovations in our proof are a Kolyvagin system argument for the Selmer group of $E$ twisted by anticyclotomic Hecke characters arbitrarily close to the trivial character, and a congruence argument exploiting Beilinson--Flach classes and their explicit reciprocity laws. 
%Under the same condition on $\phi$, our results imply the $p$-part of the Birch--Swinnerton-Dyer formula when ${\rm ord}_{s=1}L(E,s)\leq 1$.
%
%Our proof is based on a study of the anticyclotomic Iwasawa theory of $E$ over an imaginary quadratic field $K$ in which $p$ splits, and a congruence argument exploiting the cyclotomic Euler system of Beilinson--Flach classes.
\end{abstract}


%\date{\today}
\maketitle
\tableofcontents


\section{Introduction}
\addtocontents{toc}{\protect\setcounter{tocdepth}{2}}
%\addtocontents{lof}{\protect\setcounter{tocdepth}{0}}
%\renewcommand{\thetheorem}{\Alph{thm}}


Let $E$ be an elliptic curve over $\Q$, and let $p$ be an odd prime of good reduction for $E$. In the early 1970s, motivated by Iwasawa's theory for the $p$-part of class groups of number fields in $\Z_p$-extensions, Mazur 
%\cite{mazur-towers} and Mazur--Swinnerton-Dyer \cite{mazur
initiated a parallel study for the arithmetic of $E$ over the cyclotomic $\Z_p$-extension $\Q_\infty/\Q$. In particular, in \cite{mazur-towers} he proved a foundational ``control theorem'' for the $p$-primary Selmer group ${\rm Sel}_{p^\infty}(E/\Q_\infty)$, viewed as a module over the Iwasawa algebra $\Lambda_\Q=\Z_p\llbracket{{\rm Gal}(\Q_\infty/\Q)\rrbracket}$,  
%\subset\rH^1(\Q_\infty,E[p^\infty])$, allowing him to establish the finite generatedness of $E(\Q_\infty)$ under some hypotheses, 
and formulated an analogue of Iwasawa's main conjecture, expressing the characteristic ideal of the Pontryagin dual of $\X_{\rm ord}(E/\Q_\infty)={\rm Sel}_{p^\infty}(E/\Q_\infty)^\vee$ in terms of the $p$-adic $L$-function attached to $E$ by the work Mazur--Swinnerton-Dyer \cite{M-SwD}:
\begin{equation}\label{MC}
\ch_{\Lambda_\Q}\bigl(\X_{\rm ord}(E/\Q_\infty)\bigr)\overset{}=\bigl(\Lcal_p^{\rm MSD}(E/\Q)\bigr)\tag{MC}
\end{equation}
(see [\emph{op.\,cit.}, Conj.~3]). By the Weierstrass preparation theorem, conjecture (\ref{MC}) can be viewed as the equality between two integral $p$-adic polynomials attached to $E$, one by means of the arithmetic of $E$ (i.e., its Mordell--Weil and Tate--Shafarevich groups) 
over $\Q_\infty$ and the other by 
means of the modular symbols associated with $E$ (available thanks to its modularity \cite{wiles,TW,BCDT}), encoding the special values of the Hasse--Weil $L$-function of $E$ twisted by finite order characters of ${\rm Gal}(\Q_\infty/\Q)$.  


The main result in this paper is the proof (under a mild hypothesis) of Mazur's main conjecture (\ref{MC}) when $p$ is an 
\emph{Eisenstein prime} for $E$, meaning that $E$ admits a rational $p$-isogeny. (In the non-Eisenstein case, the conjecture is largely known \cite{kato-euler-systems,skinner-urban,wan-hilbert}.) In other words, we consider the case where $E[p]$ is \emph{reducible} as a module over the absolute Galois group $G_\Q={\rm Gal}(\overline{\Q}/\Q)$, so that
\[
E[p]^{ss}=\mathbb{F}(\phi)\oplus\mathbb{F}(\psi)
\] 
as $G_\Q$-modules, where $\phi,\psi:G_\Q\rightarrow\mathbb{F}^\times$ are characters whose product $\phi\psi=\omega$ is the mod $p$ cyclotomic character. 
%For Eisenstein primes $p$, Kato proved \cite{kato-euler-systems} that $\X_{\rm ord}(E/\Q_\infty)$ is $\Lambda_\Q$-torsion and that its characteristic ideal contains $\Lcal_p^{\rm MSD}(E/\Q)$ after inverting $p$. This ambiguity of powers of $p$ was subsequently removed by W\"uthrich \cite{wuthrich-int}. On the other hand, in their foundational paper \cite{greenvats}, Greenberg--Vatsal proved conjecture (\ref{MC}) under the assumption that the \emph{isogeny character} 
%\[
%\phi:G_\Q\rightarrow\mathbb{F}_p^\times,
%\] 
%i.e., the character describing the action of $G_\Q$ on the kernel $C\subset E[p]$ of the underlying rational $p$-isogeny, satisfies:
%\begin{equation}\label{GV}
%\textrm{$\phi$ is either unramified at $p$ and odd, or ramified at $p$ and even.}\tag{GV}
%\end{equation}
%Under this hypothesis, they could show that both $\X_{\rm ord}(E/\Q_\infty)$ and $\Lcal_p^{\rm MSD}(E/\Q)$ have vanishing $\mu$-invariant by building on the work of Ferrero--Washington \cite{FW} and the proof of Iwasawa's main conjecture by Mazur--Wiles \cite{MW}, thereby reducing their result on (\ref{MC}) to a delicate comparison of Iwasawa $\lambda$-invariants.

%Without hypothesis (\ref{GV}), the situation is known to be more complicated. Indeed,  Greenberg showed that if $E[p^\infty]$ contains a $G_\Q$-invariant cyclic subgroup $\Phi$ of order $p^m$ which is ramified at $p$ and odd, then $\X_{\rm ord}(E/\Q_\infty)$ has $\mu$-invariant $\geq m$ (see \cite[Prop.~5.7]{greenberg-cetraro}). On the analytic side, a conjecture by Stevens \cite{stevens-invmath} predicts a similar phenomenon for $\Lcal_p^{\rm MSD}(E/\Q)$.
%(see \cite[Prop.~5.7]{greenberg-cetraro}). 
%The methods in this paper allow us to prove Mazur's main conjecture (\ref{MC}) for Eisenstein primes regardless of the value of the $\mu$-invariant. In particular, our results include giving a new proof for the case previously handled by Greenberg--Vatsal.


%Let $G_p\subset G_\Q$ be a decomposition group at $p$ and let $\omega:G_\Q\rightarrow\mathbb{F}_p^\times$ be the Teichm\"uller character. Our main result towards Mazur's main conjecture is the following.

\begin{thmintro}\label{thm:CYC}
Let $E/\Q$ be an elliptic curve, and let $p>2$ be a prime of good reduction  for $E$. Suppose that $p$ is Eisenstein with 
$\phi\vert_{G_p}\neq 1,\omega$, 
where $G_p\subset G_\Q$ is a decomposition group at $p$. Then $\X_{\rm ord}(E/\Q_\infty)$ is $\Lambda_\Q$-torsion with
%Suppose $\phi\vert_{G_p}\neq 1,\omega$. Then $\X(E/\Q_\infty)$ is $\Lambda_\Q$-torsion, with
\[
\ch_{\Lambda_\Q}\bigl(\X_{\rm ord}(E/\Q_\infty)\bigr)\overset{}=\bigl(\Lcal_p^{\rm MSD}(E/\Q)\bigr),
\]
and hence Mazur's main conjecture holds.
%as ideals in $\Lambda_\Q$.
\end{thmintro}


In the setting of Theorem~\ref{thm:CYC}, Kato proved \cite{kato-euler-systems} that $\X_{\rm ord}(E/\Q_\infty)$ is $\Lambda_\Q$-torsion and that its characteristic ideal contains $\Lcal_p^{\rm MSD}(E/\Q)$ after inverting $p$. This ambiguity of powers of $p$ was subsequently removed by W\"uthrich \cite{wuthrich-int}. On the other hand, in a foundational paper \cite{greenvats}, Greenberg--Vatsal proved conjecture (\ref{MC}) for Eisenstein primes $p$ under the assumption that
%under the assumption that the \emph{isogeny character} 
%\[
%\phi:G_\Q\rightarrow\mathbb{F}_p^\times,
%\] 
%i.e., the character describing the action of $G_\Q$ on the kernel $C\subset E[p]$ of the underlying rational $p$-isogeny, satisfies:
\begin{equation}\label{GV}
\textrm{$\phi$ is either}\,
\begin{cases}
\textrm{unramified at $p$ and odd, or} \\
\textrm{ramified at $p$ and even.}\tag{GV}
\end{cases}
\end{equation}
Under this hypothesis, they could show that both $\X_{\rm ord}(E/\Q_\infty)$ and $\Lcal_p^{\rm MSD}(E/\Q)$ have vanishing $\mu$-invariant by building on the work of Ferrero--Washington \cite{FW} and %the proof of Iwasawa's main conjecture by 
Mazur--Wiles \cite{MW}, thereby reducing their result on (\ref{MC}) to a delicate comparison of Iwasawa $\lambda$-invariants.

Without hypothesis (\ref{GV}), the situation is known to be more complicated. Indeed,  Greenberg showed that if $E[p^\infty]$ contains a $G_\Q$-invariant cyclic subgroup $\Phi$ of order $p^m$ which is ramified at $p$ and odd, then $\X_{\rm ord}(E/\Q_\infty)$ has $\mu$-invariant $\geq m$ (see \cite[Prop.~5.7]{greenberg-cetraro}). On the analytic side, a conjecture by Stevens \cite{stevens-invmath} predicts a similar phenomenon for $\Lcal_p^{\rm MSD}(E/\Q)$.
%(see \cite[Prop.~5.7]{greenberg-cetraro}). 
The methods in this paper allow us to prove Mazur's main conjecture (\ref{MC}) for Eisenstein primes regardless of the value of the $\mu$-invariant. In particular, our results include giving a new proof for the case previously handled by Greenberg--Vatsal.


\subsection{Some ideas from the proof} 

Our proof of Theorem~\ref{thm:CYC} goes through  \emph{anticyclotomic} Iwasawa theory for $E$ over an auxiliary imaginary quadratic field $K/\Q$ in which $p$ splits. Roughly speaking, this is used to show that $\mathfrak{X}_{\rm ord}(E/\Q_\infty)$ and $\Lcal_p^{\rm MSD}(E/\Q)$ have the same Iwasawa invariants\footnote{Strictly speaking, our results only show these equalities for \emph{the sum} of the Iwasawa invariants of $E$ and the quadratic twist $E^K$, but this suffices for the proof of Theorem~\ref{thm:CYC} thanks to Kato's work.}: 
\[
\mu(\mathfrak{X}_{\rm ord}(E/\Q_\infty))=\mu(\Lcal_p^{\rm MSD}(E/\Q))
\quad\textrm{and}\quad
\lambda(\mathfrak{X}_{\rm ord}(E/\Q_\infty))=\lambda(\Lcal_p^{\rm MSD}(E/\Q)),
\]
even in situations of positive $\mu$-invariant. 
%(Strictly speaking, our results only show these equalities for \emph{the sum} of the Iwasawa invariants of $E$ and the quadratic twist $E^K$, but this suffices for the proof of Theorem\,\ref{thm:CYC} thanks to Kato's work.) 


%From Kato's work--as refined in by W\"uthrich \cite{wuthrich-int}--applied to $E$ and $E^K$, it follows that a Selmer groups $\mathfrak{X}_{\rm ord}(E/K_\infty^+)$ 

More precisely, our argument rests on the proof of an Iwasawa main conjecture for $E$ over the anticyclotomic $\Z_p$-extension of $K$ (see Theorem~\ref{thm:AC} below), and a congruence argument building on the cyclotomic Euler system of Beilinson--Flach classes of Lei--Loeffler--Zerbes \cite{RSCMF} and Kings--Loeffler--Zerbes \cite{explicit}. 
%In the rest of this introduction we describe these ingredients in some detail.


%Two key ingredients in th of Theorem~\ref{thm:CYC} are the proof of an Iwasawa main conjecture for $E$ over the \emph{anticyclotomic} $\Z_p$-extension of an auxiliary imaginary quadratic field $K/\Q$ in which $p$ splits, and a special case of the cyclotomic Euler system of Beilinson--Flach classes of Lei--Loeffler--Zerbes \cite{RSCMF} and Kings--Loeffler--Zerbes \cite{explicit}. Roughly speaking, this is used to show that $\mathfrak{X}_{\rm ord}(E/\Q_\infty)$ and $\Lcal_p^{\rm MSD}(E/\Q)$ have \emph{the same} Iwasawa invariants, even in situations of positive $\mu$-invariant.  
 
 %serving as a bridge between the cyclotomic and anticyclotomic settings. 
%

\subsubsection*{Anticyclotomic main conjectures}

Denote by $N$ the conductor of $E$, and let $K$ be an imaginary quadratic field such that
\begin{equation}\label{eq:intro-disc}
\textrm{the discriminant $D_K$ is odd and $D_K\neq -3$,}\tag{disc}
\end{equation}
and such that the following \emph{Heegner hypothesis} holds:
\begin{equation}\label{eq:intro-Heeg}
\textrm{every prime $\ell\vert N$ splits in $K$.}\tag{Heeg}
\end{equation}
Let $\Gamma_K^-={\rm Gal}(K^{-}_{\infty}/K)$ be the Galois group of the anticyclotomic $\Z_p$-extension of $K$, and 
%let $\gamma^-\in\Gamma_K^-$ be a topological generator. 
for each $n$ denote by $K_n^-$ the subfield of $K^{-}_{\infty}$ with $[K_n^-:K]=p^n$. In sharp contrast with the case of $\Q_\infty/\Q$, one can show that ${\rm rank}_\Z E(K_n^-)$ is unbounded as $n\to\infty$, and therefore the Pontryagin dual $\X_{\rm ord}(E/K_\infty^-)$ 
of the Selmer group ${\rm Sel}_{p^\infty}(E/K_\infty^-)$ has positive rank as a module over the anticyclotomic Iwasawa algebra $\Lambda_K^-=\Z_p\llbracket{\Gamma_K^-}\rrbracket$. This unbounded growth is accounted for by the existence of Heegner points on $E$ associated with a given modular parametrization
\[
\pi:X_0(N)\rightarrow E.
\]
This system of points gives rise to a $\Lambda_K^-$-adic class $\kappa_1^{\rm Hg}$ which was shown to be non-torsion by Cornut \cite{cornut} and Vatsal \cite{vatsal} in their proof of ``Mazur's conjecture'' on higher Heegner points \cite{mazur-ICM83}. 
As recalled below, a formulation of the Iwasawa main conjecture in this context was given by Perrin-Riou \cite{perrinriou}. Write
\[
%\mathcal{S}:=\biggl(\varprojlim_n\varprojlim_n{\rm Sel}_{p^m}(E/K_n^-)\biggr)\otimes\Q_p,\quad\mathcal{X}:=\mathfrak{X}(E/K_\infty^-)\otimes\Q_p,
\mathfrak{S}_{\rm ord}(E/K_\infty^-):=\varprojlim_n\varprojlim_m{\rm Sel}_{p^m}(E/K_n^-),
\]
which is a compact $\Lambda_K^-$-module containing $\kappa_1^{\rm Hg}$. 

%DO WE NEED TO INVERT P?-------------------
%Let us denote by $\mathfrak{S}_{\rm ord}(E/K_\infty^-)_{\Q_p}$ (respectively $\X_{\rm ord}(E/K_\infty^-)_{\Q_p}$) the $\Lambda^-_K\otimes \Q_p$-module $\mathfrak{S}_{\rm ord}(E/K_\infty^-)\otimes \Q_p$ (resp. $\X_{\rm ord}(E/K_\infty^-)\otimes \Q_p$).
%\begin{conj}\label{conj:PR}
%Let $E/\Q$ be an elliptic curve, $p>2$ a prime of good ordinary reduction for $E$, and let $K$ be an imaginary quadratic field satisfying {\rm(\ref{eq:intro-Heeg})} and {\rm(\ref{eq:intro-disc})}. Then $\mathfrak{S}_{\rm ord}(E/K_\infty^-)_{\Q_p}$ and $\X_{\rm ord}(E/K_\infty^-)_{ \Q_p}$ both
%have $\Lambda^-_K\otimes \Q_p$-rank one, and
%\[
%{\rm char}_{\Lambda^-_{K}\otimes\Q_p}\bigl(\X_{\rm ord}(E/K_\infty^-)_{\Q_p,\rm tors}\bigr)={\rm char}_{\Lambda^-_{K}\otimes \Q_p}\bigl(\mathfrak{S}_{\rm ord}(E/K_\infty^-)_{\Q_p}/(\kappa_1^{\rm Hg})\bigr)^2 
%\] 
%where the subscript {\rm tors} denotes the $\Lambda^-_{K}\otimes \Q_p$-torsion submodule. 
%\end{conj}
%
%%by $\mathcal{H}\subset\mathcal{S}$ the $\Lambda_K^-\otimes\Q_p$-submodule generated by $\kappa_1^{\rm Hg}$.
\begin{conj}[Perrin-Riou]\label{conj:PR}
Let $E/\Q$ be an elliptic curve, $p>2$ a prime of good ordinary reduction for $E$, and let $K$ be an imaginary quadratic field satisfying {\rm(\ref{eq:intro-Heeg})} and {\rm(\ref{eq:intro-disc})}. Then $\mathfrak{S}_{\rm ord}(E/K_\infty^-)$ and $\X_{\rm ord}(E/K_\infty^-)$ both
have $\Lambda^-_K$-rank one, and
\[
{\rm char}_{\Lambda^-_{K}}\bigl(\X_{\rm ord}(E/K_\infty^-)_{\rm tors}\bigr)={\rm char}_{\Lambda^-_{K}}\bigl(\mathfrak{S}_{\rm ord}(E/K_\infty^-)/(\kappa_1^{\rm Hg})\bigr)^2, 
\] 
where the subscript {\rm tors} denotes the $\Lambda^-_{K}$-torsion submodule. 
\end{conj}

The first general result towards Conjecture~\ref{conj:PR} for Eisenstein primes $p$ was obtained in \cite{eisenstein}. Namely, it was shown 
%(see [\emph{op.\,cit.}, Cor.~D]) 
that Perrin-Riou's main conjecture holds for Eisenstein primes $p$ under the additional hypotheses that \begin{equation}\label{eq:intro-spl}
\textrm{$(p)=v\bar{v}$ splits in $K$,}\tag{spl} 
\end{equation} 
that $\phi\vert_{G_p}\neq 1,\omega$, and that 
\begin{equation}\label{eq:intro-Sel}
\textrm{the $\Z_p$-corank of ${\rm Sel}_{p^\infty}(E/K)$ is $1$}.\tag{Sel}
\end{equation}
%(see [\emph{op.\,cit.}, Cor.~D].) 


In particular, hypothesis (\ref{eq:intro-Sel}), %---which was harmless for the arithmetic applications in analytic rank $\leq 1$ studied in \cite{eisenstein}, but 
which obviously excludes elliptic curves with ${\rm rank}_\Z E(\Q)\geq 2$, was imposed to account for the inability of the methods in [\emph{op.\,cit.}, \S{3}] to control certain error terms appearing at height one primes of $\Lambda_K^-$ approaching the augmentation ideal $\mathfrak{P}_0\subset \Lambda_K^-$. 

A key technical innovation in this paper is the proof of a Kolyvagin system argument for the Selmer group of $E/K$ twisted by characters of $\Gamma_K^-$ arbitrarily close to $1$ yielding the ``upper bound''  divisibility in Conjecture~\ref{conj:PR} after localization at height one primes of $\Lambda_K^-$ approaching $\mathfrak{P}_0$. %including $\mathfrak{P}_0$ itself.
%The Kolyvagin system arguments in this paper 
%(contained in Sect.~\ref{sec:anticyc}) 
%complement those developed in 
Together with complementary results obtained in \cite{eisenstein}, this argument yields the following.

\begin{thmintro}\label{thm:AC}
Let $E/\mathbb{Q}$ be an elliptic curve, let $p\nmid 2N$ be an Eisenstein prime for $E$, and let $K$ be an imaginary quadratic field satisfying {\rm (\ref{eq:intro-Heeg})}, {\rm (\ref{eq:intro-disc})}, and {\rm (\ref{eq:intro-spl})}. Suppose that   
$\phi\vert_{G_p}\neq 1,\omega$. Then Conjecture~\ref{conj:PR} holds.
\end{thmintro}

To go from Theorem~\ref{thm:AC} to Theorem~\ref{thm:CYC}, we use a reformulation of the former in terms of $p$-adic $L$-functions. Let $\Z_p^{\rm ur}$ be the completion of the ring of integers of the maximal unramified extension of $\Q_p$, and put
\[
\Lambda^{-,\rm ur}_K:=\Lambda_K^-\widehat\otimes_{\Z_p}\Z_p^{\rm ur}.
\]
%By the well-known equivalence between 
It follows from the explicit reciprocity law for $\kappa_1^{\rm Hg}$, that Conjecture~\ref{conj:PR} is equivalent to the Iwasawa--Greenberg main conjecture for the $p$-adic $L$-function $\Lcal_p^{\rm BDP}(f/K)$ introduced in \cite{BDP}. %(see e.g. \cite[Thm.~5.2]{BCK}), 
Under the same hypotheses of Theorem~\ref{thm:AC}, we thus deduce that a different Greenberg Selmer group denoted $\X_{\rm Gr}(E/K_\infty^-)$ 
%(with the relaxed and the strict local condition at the primes above $v$ and $\bar{v}$, respectively) 
is $\Lambda_K^-$-torsion, with
\[
{\rm char}_{\Lambda_K^-}\bigl(\mathfrak{X}_{\rm Gr}(E/K_{\infty}^-)\bigr)\Lambda_K^{-,\rm ur}=\bigl(\Lcal_p^{\rm BDP}(f/K)\bigr)
\]
as ideals in $\Lambda_K^{-,\rm ur}$. This equality of characterisitc ideals is the first key ingredient in the proof of Theorem~\ref{thm:CYC}.

\subsubsection*{Comparing Iwasawa invariants}


In \cite{RSCMF} and \cite{explicit}, 
%(extending earlier construction of Bertolini--Darmon--Rotger \cite{BDR1,BDR2}, and Lei--Loeffler--Zerbes \cite{LLZ}), 
Lei--Loeffler--Zerbes and Kings--Loeffler--Zerbes constructed a cyclotomic Euler system %of Beilinson--Flach classes 
(over $\Q$) for the Rankin--Selberg convolution of two modular forms moving in Hida families. To be able to use these classes as a bridge between the anticyclotomic $\Z_p$-extension $K_\infty^-/K$ and the cyclotomic $\Z_p$-extension $\Q_\infty/\Q$ (or rather its translate by $K$),  %(crossing at the trivial character), 
here we are led to consider a variant of their construction attached to the pair $(f,\boldsymbol{g})$, where $f$ %\in S_2(\Gamma_0(N))$
is the weight $2$ newform attached to $E$, and $\boldsymbol{g}$ is a suitable CM Hida family. Because our $\boldsymbol{g}$ specializes in weight $1$ to the $p$-irregular Eisenstein series ${\rm Eis}_{1,\eta}$, 
%(since we assume (\ref{eq:spl})), 
where $\eta=\eta_{K/\Q}$ is the quadratic character associated to $K/\Q$, in fact we use a refinement of the construction in \cite{explicit} studied in \cite{BST}.

Let $\Lambda_K$ (resp. $\Lambda_K^+$) be the Iwasawa algebra for the $\Z_p^2$-extension of $K$ (resp. the cyclotomic $\Z_p$-extension $K_\infty^+/K$). In particular, from these works we obtain a two-variable Iwasawa cohomology class 
\[
BF\in\rH^1_{\rm Iw}(K_\infty,T_pE_\bullet),
\]
where $E_\bullet/\Q$ is the distinguished elliptic curve in the isogeny class of $E$ constructed by W\"uthrich \cite{wuthrich-int}. Combined with the  relations between different $p$-adic $L$-functions established in Sect.~\ref{sec:Lp}, we also deduce two explicit reciprocity laws:
\begin{enumerate}
\item One relating ${\rm loc}_{\bar{v}}(BF)$ to a $p$-adic $L$-function $\Lcal_p^{\rm PR}(E_\bullet/K)\in\Lambda_K$ whose image $\Lcal_p^{\rm PR}(E_\bullet/K)^+$ under the natural projection $\Lambda_K\rightarrow\Lambda_K^+$ satisfies
\[
\Lcal_p^{\rm PR}(E_\bullet/K)^+=\Lcal_p^{\rm MSD}(E_\bullet/\Q)\cdot\Lcal_p^{\rm MSD}(E_\bullet^K/\Q)
\]
up to a unit, where $E_\bullet^K$ is the twist of $E_\bullet$ by the quadratic field $K$.
%character $\eta$; 
\item A second one relating ${\rm loc}_v(BF)$ to a $p$-adic $L$-function $\Lcal_p^{\rm Gr}(f/K)\in\Lambda_K^{\rm ur}=\Lambda_K\widehat\otimes\Z_p^{\rm ur}$ whose image $\Lcal_p^{\rm Gr}(f/K)^-$ under the natural projection $\Lambda_K^{\rm ur}\rightarrow\Lambda_K^{-,\rm ur}$ satisfies
\[
\Lcal_p^{\rm Gr}(f/K)^-=\Lcal_p^{\rm BDP}(f/K)
\]
up to a unit. 
\end{enumerate}

%The above relations between $p$-adic $L$-functions and their counterparts for characteristic ideals of Selmer groups are shown in Sect.~\ref{sec:Lp} and Sect.~\ref{sec:Sel}, respectively.
%With these results at hand, 
%(established in Sect.~\ref{BF}, following the above relations between $p$-adic $L$-functions and their counterparts for characteristic ideals of Selmer groups contained in Sect.~\ref{sec:Lp} and Sect.~\ref{sec:Sel}, respectively), 
%one can give a simplified proof of Theorem~\ref{thm:CYC} under some additional hypotheses, including hypothesis (\ref{eq:intro-Sel}) for a suitably chosen $K$ (of course, this condition cannot be arranged as long as ${\rm rank}_\Z E(\Q)\geq 2$). 
%As it helps clarify our overall strategy and motivates the more technical arguments that are needed in general, in Sect.\,\ref{sec:one} we explain the proof of Theorem~\ref{thm:CYC} in this rank one case. 
%before tackling the general case.

The next key idea to have Theorem~\ref{thm:AC} to bring to bear on the proof of Theorem~\ref{thm:CYC} is to consider the restriction of $BF$ to $\rH^1_{\rm Iw}(K_\infty^+,T_pE_\bullet)$, and exploit the fact that the anticyclotomic and the cyclotomic $\Z_p$-extensions meet at the trivial character (in other words, $K_\infty^-\cap K_\infty^+=K$). When $\Lcal_p^{\rm BDP}(f/K)(1)\neq 0$ (which by the main result of \cite{BDP} can only happen when ${\rm rank}_\Z E(\Q)\leq 1$), the argument for the implication Theorem~\ref{thm:AC}$\Longrightarrow$Theorem~\ref{thm:CYC} is relatively simple, and to help orient the reader, this simpler case is explained in detail in Sect.~\ref{sec:one}.  
%simplified proof of Theorem\,\ref{thm:CYC}.  As an illustration and motivation of the more technical arguments that are needed in general, in Sect.\,\ref{sec:one} we explain the proof in this rank one case. %before tackling the general case.

To make the argument work in arbitrary rank, we take a character $\alpha$ of $\Gamma_K^-$ with
\[
\alpha\equiv 1\;({\rm mod}\,p^m)
%\alpha\equiv 1\pmod{p^m}
\]
for some $m\gg 0$ such that $\Lcal_p^{\rm BDP}(f/K)(\alpha)\neq 0$ (as always possible by the nonvanishing of $\Lcal_p^{\rm BDP}(f/K)$), and consider $\alpha$-twisted versions of the above Selmer groups and $p$-adic $L$-functions projected to $\Lambda_K^+$. With the aid of a cyclotomic Euler system extending the $\alpha$-twist of the class $BF$ projected to $\rH^1_{\rm Iw}(K_\infty^+,T_pE_\bullet)$, building on Theorem~\ref{thm:AC} we deduce that the twisted Selmer group $\mathfrak{X}_{\rm ord}(E_\bullet(\alpha)/K_\infty^+)$ is $\Lambda_K^+$-torsion, with
\begin{equation}\label{eq:alpha-MC}
\ch_{\Lambda_K^+}\bigl(\X_{\rm ord}(E_\bullet(\alpha)/K_\infty^+)\bigr)=\bigl(\Lcal_p^{\rm PR}(E_\bullet(\alpha)/K)^+\bigr)
\end{equation} 
as ideals in $\Lambda_K^+$. On the other hand, Kato's divisibility \cite{kato-euler-systems} (as refined by W\"uthrich \cite{wuthrich-int}) applied to $E_\bullet$ and $E_\bullet^K$ implies that the untwisted Selmer group $\mathfrak{X}_{\rm ord}(E_\bullet/K_\infty^+)$ is $\Lambda_K^+$-torsion, with the divisibility
\begin{equation}\label{eq:one-MC}
\ch_{\Lambda_K^+}\bigl(\X_{\rm ord}(E_\bullet/K_\infty^+)\bigr)\supset\bigl(\Lcal_p^{\rm PR}(E_\bullet/K)^+\bigr)
\end{equation} 
as ideals in $\Lambda_K^+$. By a new congruence argument using the study of the variation of both sides of \eqref{eq:alpha-MC} as $\alpha$ varies carried out in the earlier parts of the paper, %building on the earlier results in this paper,  
%with $\alpha$ sufficiently close to $1$, 
we deduce from this equality (for $\alpha$ sufficiently close to $1$) that $\X_{\rm ord}(E_\bullet/K_\infty^+)$ and $\Lcal_p^{\rm PR}(E_\bullet/K)^+$ have the same Iwasawa invariants, and so equality holds in (\ref{eq:one-MC}). The proof of Theorem~\ref{thm:CYC} for both $E_\bullet$ and the original elliptic curve $E$, can then be deduced from Kato's work and the invariance of Mazur's main conjecture under isogenies.  

%When $\Lcal_p^{\rm BDP}(f/K)(1)\neq 0$ (which by the main result of \cite{BDP} can only happen when ${\rm rank}_\Z E(\Q)\leq 1$), one can give a simplified proof of Theorem\,\ref{thm:CYC}.  As an illustration and motivation of the more technical arguments that are needed in general, in Sect.\,\ref{sec:one} we explain the proof in this rank one case. %before tackling the general case.


\subsection{Application to the $p$-part of the Birch--Swinnerton Dyer formula}

As a standard consequence of Theorem~\ref{thm:CYC}, 
 we deduce most cases of the $p$-part of the Birch--Swinnerton Dyer formula for elliptic curves $E/\Q$ of analytic rank $\leq 1$ at Eisenstein primes $p$. 

\begin{thmintro}\label{thm:pBSD}
Let $E/\Q$ and $p>2$ be as in Theorem~\ref{thm:CYC}, and let $r\in\{0,1\}$. If ${\rm ord}_{s=1}L(E,s)=r$, then
%If $E/\Q$ has analytic rank $r\in\{0,1\}$, $p>2$ is a prime of good reduction, and $E$ has a rational $p$-isogeny with kernel $\mathbb{F}_p(\phi)$ with $\phi\vert_{G_p}\neq 1, \omega$, then
\begin{displaymath}
\ord_{p}\biggl(\frac{L^{(r)}(E, 1)}{\operatorname{Reg}(E / \Q) \cdot \Omega_{E}}\biggr)=\ord_{p}\biggl(\# \sha(E / \Q) \prod_{\ell \nmid \infty} c_{\ell}(E / \Q)\biggr),
\end{displaymath}
where 
\begin{itemize}
\item ${\rm Reg}(E/\Q)$ is the regulator of $E(\Q)$, 
\item $\Omega_E=\int_{E(\mathbb{R})}\vert\omega_E\vert$ is the N\'{e}ron  period associated to the N\'{e}ron differential $\omega_E$, and 
\item $c_\ell(E/\Q)$ is the Tamagawa number of $E$ at the prime $\ell$.
\end{itemize}
In other words, the $p$-part of the Birch--Swinnerton-Dyer formula for $E$ holds.
\end{thmintro}

\begin{proof}
In the case $r=0$, the argument is the same as in  \cite[Thm.~5.1.4]{eisenstein}, replacing the appeal to \cite[Thm.~1.3]{greenvats} by an appeal to our Theorem~\ref{thm:CYC}. In the case $r=1$, we argue as in \cite[Thm.~5.3.1]{eisenstein}, choosing a suitable $K$ with $L(E^K,1)\neq 0$ and applying our result in the rank $0$ case  to $E^K$.
\end{proof}

\begin{rem}
Previously, by the work of Greenberg--Vatsal \cite{greenvats} in analytic rank $0$, and of the authors with Lee \cite{eisenstein} in analytic rank $1$, Theorem~\ref{thm:pBSD} was only known for roughly ``half'' of the $p$-Eisenstein cases, as both results required a certain parity hypothesis on $\phi$.
\end{rem}

%Note that %as in \cite[$\S$5.3]{eisenstein} 
%the previously known cases of Theorem~\ref{thm:pBSD} for analytic rank $r=0$ (resp. $r=1$) required $E$ (resp. an odd quadratic twist of $E$) to satisfy (\ref{GV}), accounting for roughly only ``half'' of the possible cases.

It is easy to produce examples of elliptic curves to which Theorem \ref{thm:pBSD} can be applied. Let $p=5$ and $J$ be any of the three non-isomorphic elliptic curves of conductor 11; one could take for example 
the elliptic curve of Weierstrass equation
\[
y^2+y=x^3-x^2-10x-20,
\]
which is the strong Weil curve of conductor $11$.  
It is known that $J[5]$ is reducible and its semi-simplification is isomorphic to $\mu_5\oplus \Z/5\Z$. Moreover, the rank of $J$ is equal to zero. We can then consider $\psi$ any quadratic character such that $\psi(5)=-1$ and take $E=J_\psi$ the quadratic twist of $J$ by $\psi$. Since the root number of $J$ is 1, applying \cite[Thm.~B.1]{twist}, we can produce infinitely many $\psi$ such that $E=J_\psi$ has analytic (and hence algebraic) rank zero and for which the hypothesis of Theorem \ref{thm:pBSD} are satisfied for $p=5$. Note that prior to this paper, the only elliptic curves in this family for which the theorem was known are of the form $E=J_\psi$ where $\psi$ is odd.
%Note that the elliptic curves in this family for which the theorem was already known are the ones satisfying (\ref{GV}) (using \cite[Thm. 1.3]{greenvats}), i.e. the elliptic curves $E=J_\psi$ where $\psi$ is odd. 

Similarly, take $p=5$ and consider the elliptic curve
\begin{displaymath}
J: y^2+y=x^3+x^2-10x+10.
\end{displaymath} 
The curve $J$ has conductor $123$ and analytic rank $1$, and satisfies $J[5]^{ss}=\mu_5\oplus \Z/5\Z$ as $G_\Q$-modules. Let $\psi$ any quadratic character such that $\psi(5)=-1$ and take $E=J_\psi$ the quadratic twist of $J$ by $\psi$. Since the root number of $J$ is $-1$ (being of analytic rank one), by \cite[Thm.~B.2]{twist} we can find infinitely many $\psi$ as above for which $E=J_\psi$ also has analytic rank one and hence to which Theorem~\ref{thm:pBSD} can be applied for $p=5$. 
%Note that prior to this paper, the only elliptic curves in this family for which the theorem was known are of the form 
%Because of (\ref{GV}), 
Earlier result in this direction would only apply to the elliptic curve in this family of the form $E=J_\psi$ with $\psi$ even. 
%Note that the elliptic curves in this family for which the theorem was already known are the ones satisfying the negation of (\ref{GV}) (using \cite[Thm. F]{eisenstein}), i.e. the elliptic curves $E=J_\psi$ where $\psi$ is even. 

One can proceed similarly for $p=3,7,13,37$, taking quadratic twists of elliptic curves of rank zero or one and with a rational $p$-isogeny. If the character describing the kernel of the isogeny is not trivial (which has to be the case for $p=13$ or $p=37$), one might have to impose further conditions to the quadratic character at some bad primes in order to apply \cite[Thm.~B]{twist}.


\subsection{Further applications and relation to previous works}

In addition to being a key ingredient in the proof of our main results, the Kolyvagin system argument with error terms for the $p$-adic Tate module of $E$ twisted by characters close to the trivial one 
contained in Sect.~\ref{sec:anticyc} (see Theorem~\ref{thm:Zp-twisted}) is used in a forthcoming work of the authors with A.\,Burungale \cite{BCGS} to establish Kolyvagin's conjecture on the nonvanishing of derived Heegner classes under  mild hypotheses on $E[p]$, including the first cases for  Eisenstein primes $p$ (many cases in the non-Eisenstein case were first proved by W.\,Zhang \cite{wei-zhang}).

Finally, we note that the idea of using Beilinson--Flach classes to relate different Iwasawa main conjectures has appeared in earlier works, notably in \cite{wan-ss,CW-adv}, but the congruence argument mechanism introduced in this paper to deduce Theorem~\ref{thm:CYC} from Theorem~\ref{thm:AC} is new, and should be useful in other settings, 
%where vanishing of the relevant $\mu$-invariants (as is the case in the aforementioned references) is not known.
%, 
as we plan to exploit in forthcoming work. %In particular, this part of our argument also applies in the non-Eisenstein case.

%of bridging between the anticyclotomic and cyclotomic settings have appeared in earlier works, notably in \cite{wan-ss} on Kobayashi's main conjecture, and of Castella--Wan \cite{CW-adv}


%As further application of Theorem \ref{thm:AC} (and Theorem \ref{thm:CYC} respectively), we also point out that in \cite{kolyconj} we prove Kolyvagin's conjecture about the non-vanishing of the Kolyvagin system of Heegner points (and the analogous statement for the Kolyvagin system obtained from Kato's classes respectively) at Eisenstein primes. Crucial to the proof are the equalities in the main conjectures proved in this paper alongside the upper bound on the size of the Selmer group of certain twists of the Tate module of the elliptic curve in terms of the index of a Kolyvagin system class, which in the anticyclotomic setting we prove in Sect.~\ref{sec:anticyc} (Theorem~\ref{thm:Zp-twisted}).


%Recall that as in \cite[$\S$5.3]{eisenstein}, the case $r=1$ is proved choosing a suitable quadratic imaginary field $K$ so that Conjecture \ref{conj:PR} holds for $E/K$ and conjecture \ref{MC} holds for $E^K/\Q$. Therefore previous known cases of Theorem~\ref{thm:pBSD} for $p$ Eisenstein relied on the main result of \cite{greenvats} and one had to assume \eqref{GV} for $E/\Q$ if $r=0$ and its negation (i.e. \eqref{GV} for $E^K/\Q$) if $r=1$.

\subsection{Outline of the paper}

We begin in Sect.~\ref{sec:Lp} by introducing the different one- and two-variable 
$p$-adic $L$-functions, and the various relations between them that will be needed for our arguments. Then in Sect.~\ref{sec:Sel} we introduce corresponding Selmer groups, both primitive and imprimitive, and prove some ancillary results for the latter ones. In Sect.~\ref{BF} we discuss Beilinson--Flach classes and some direct applications of their relations with $p$-adic $L$-functions. As a guide to the general case, 
in Sect.~\ref{sec:one} we give a simplified proof of Theorem~\ref{thm:CYC} in the case of rank one.  
%illustrate our general strategy by outlining a proof of Theorem~\ref{thm:CYC} under some additional hypotheses (including (\ref{eq:intro-Sel}) for a suitably chosen $K$), for which the earlier results in this paper and those in \cite{eisenstein} are sufficient. 
In Sect.~\ref{sec:anticyc}
%, which is the most technical section of the paper, 
we prove our new Kolyvagin system bound for twists by characters arbitrarily close to $1$, resulting in Theorem~\ref{thm:AC}.
%the key Theorem~\ref{thm:Zp-twisted} for the $p$-adic Tate module of $E$ twisted by characters $\alpha$ of $\Gamma_K^-$ that are close to $1$. The result is an upper bound on the size of the Selmer group in terms of the index of a Kolyvagin system class with an `error term' that remains bounded as $\alpha$ approaches $1$. 
%The proof of Theorem~\ref{thm:AC} then follows easily. 
Finally, in Sect.~\ref{sec:MC} we run our congruence argument and conclude the proof of Theorem~\ref{thm:CYC}.

%Finally, similarly as in \cite{eisenstein}, we note that the methods and result of this paper should easily extend to the case of cuspidal newforms of weight two and trivial character that a congruent to an Eisenstein series at a prime above $p$. We leave this, and the further extension of our results to Hilbert modular forms, to future work.


%\subsection{Strategy and relation to previous works} 

%Let us know explain the main ideas that go into the proofs of these results. In \cite[$\S\S$1-2]{eisenstein}, we have proved, assuming $\phi_{G_{\Q_p}}\neq 1,\omega$, that the $\lambda$ and $\mu$-invariants in Conjecture \ref{conj:BDP} coincide and that the latter are zero. We then proved in $\S$3 of \emph{op. cit.} one of the divisibilities in Conjecture~\ref{conj:PR}: 
%{${\rm char}_{\Lambda^-}(\mathcal{X}_{\rm tors})$ divides ${\rm char}_{\Lambda^-}\bigl(\mathcal{S}/\mathcal{H}\bigr)^2$ in 
%$\Lambda^-[\frac{1}{p},\frac{1}{T}]$,
%which yields a corresponding divisibility in Conjecture \ref{conj:BDP}. The proof of this result relies on a Kolyvagin system argument for twists by anticyclotomic characters sufficiently close $p$-adically to the ``zeroes of $f_{\Lambda^-}$'', where $f_{\Lambda^-}$ is a generator of ${\rm char}_{\Lambda^-}\bigl(\mathcal{S}/\mathcal{H}\bigr)$. This allows to prove that for $m\gg 0$ and for every height one prime ideal $\mathfrak{P}=(g(T))$ of $\Lambda^-$ dividing $f_{\Lambda^-}$ we have
%\begin{equation}\label{eq:oldbound}
%m\cdot{\rm ord}_{\mathfrak{P}}\bigl({\rm char}_{\Lambda^-}(\mathcal{X}_{\rm tors}\bigr)\bigr)\leqslant 2m\cdot{\rm ord}_{\mathfrak{P}}(f_{\Lambda^-})+ \mathcal{E}_{\mathfrak{P},m}
%\end{equation}
%for an `error term' $\mathcal{E}_{\mathfrak{P},m}$ depending on $T_p(E)$, ${\rm rk}_{\Z_p}(\Lambda^-/\mathfrak{P})$ and linearly on the $p$-adic valuation of $g(0)+p^m$. In particular, if $\mathfrak{P}=(T)$ divides $f_{\Lambda^-}$, since $\mathcal{E}_{T,m}$ depends linearly on $m$, we cannot let $m$ go to infinity and deduce ${\rm ord}_{T}\bigl({\rm char}_{\Lambda^-}(\mathcal{X}_{\rm tors}\bigr)\bigr)\leqslant 2\;{\rm ord}_{T}(f_{\Lambda^-})$. Therefore, in this work, we have proved in $\S$\ref{sec:anticyc} (see Theorem \ref{thm:Zp-twisted}) a variation of \cite[Theorem 3.2.1]{eisenstein} which applies only to characters $\alpha$ $p$-adically close to the trivial character (say $\alpha \equiv 1 \mod p^m$) and involves an error term not depending on $m$. The strategy is similar to the one employed in \cite[$\S$3.3.3]{eisenstein} but technically more delicate. In order to remove the dependence on $m$, we tweak the \v{C}ebotarev argument to produce the Kolyvagin primes at the cost of having to work with the $p^m$-torsion of the Selmer groups involved (see Remark \ref{rmkkolyvaginargument} for more details). We then obtain \eqref{eq:oldbound} for $\mathfrak{P}=(T)$ and an error term $\mathcal{E}_{\mathfrak{P}}$ such that $\mathcal{E}_{\mathfrak{P}}/m$ goes to zero as $m$ goes to infinity, therefore proving Theorem \ref{thm:AC}.

%Having proved the anticyclotomic main conjecture, we then exploit the equality ${\rm char}_\Lambda^-\bigl(\mathfrak{X}_E\bigr)\Lambda^{\rm ur}=\bigl(\Lcal_p^{\rm BDP}(E/K)\bigr)$ at some choice of anticyclotomic character $\alpha$ congruent to one modulo a sufficiently high power of $p$ and such that $\Lcal_p^{\rm BDP}(E/K)(\alpha)\neq 0$ to prove Theorem \ref{thm:CYC}. More precisely, 


\subsection{Acknowledgements}
We would like to thank Ashay Burungale, Shinichi Kobayashi, and Romyar Sharifi for useful exchanges related to this work. We also thank the referees for their useful comments and suggestions.  
During the preparation of this paper, F.C. was supported by the NSF grants DMS-2101458 and DMS-2401321; G.G. was partially supported by the postdoctoral fellowship of the Fondation Sciences Math\'{e}matiques de Paris; C.S. was partially supported by the Simons Investigator Grant \#376203 from the Simons Foundation and by the NSF grant DMS-1901985. 



\section{$p$-adic $L$-functions}
\addtocontents{toc}{\protect\setcounter{tocdepth}{2}}
\label{sec:Lp}

Let $E/\Q$ be an elliptic curve of conductor $N$, and let $p\nmid 2N$ be a prime of good ordinary reduction for $E$. Fix embeddings $\iota_p:\overline{\Q}\hookrightarrow\overline{\Q}_p$ and $\iota_\infty:\overline{\Q}\hookrightarrow\C$. Let $K$ be an imaginary quadratic field of discriminant $D_K<0$ prime to $N$, and assume that
\begin{equation}\label{eq:spl}
\textrm{$(p)=v\bar{v}$ splits in $K$,}\tag{spl}
\end{equation}
with $v$ the prime above $p$ induced by $\iota_p$. 
We denote by $\Gamma_{\Q}$  (resp. $\Gamma_K$) the Galois group of the cyclotomic $\Z_p$-extension $\Q_{\infty}/\Q$ (resp. the $\Z_p^2$-extension $K_{\infty}/K$). Since $p>2$ the action of complex conjugation gives an eigenspace decomposition 
\[
\Gamma_K=\Gamma_K^+\times \Gamma_K^-.
\] 
Note that $\Gamma_K^+$ (resp. $\Gamma_K^-$) is identified with the Galois group of the cyclotomic $\Z_p$-extension $K_\infty^+/K$ (resp. the anticyclotomic $\Z_p$-extension of $K^-_{\infty}/K$), and 
%we have $\Gamma^+_K\simeq \Gamma_\Q$. Let
hence $\Gamma^+_K$ is naturally identified with $\Gamma_\Q$. Let 
\[
\Lambda_\Q=\Z_p\llbracket{\Gamma_\Q}\rrbracket,\quad
\Lambda_K=\Z_p\llbracket{\Gamma_K}\rrbracket,\quad\Lambda_K^\pm=\Z_p\llbracket\Gamma_K^\pm\rrbracket
\]
be the corresponding Iwasawa algebras, so in particular $\Lambda_K^+$ is naturally identified with $\Lambda_\Q$.
%We write $\gamma^{\pm}$ for a fixed topological generator of $\Gamma^\pm$.


\subsection{Cyclotomic $p$-adic $L$-function}
\label{subsec:L-cyc} 

Fix a modular parametrization $\pi_E:X_0(N)\rightarrow E$, and denote by $c_E\in\Z$ the corresponding Manin constant, so that 
\[
\pi_E^*(\omega_E)=c_E\cdot 2\pi i f(\tau)d\tau,
\] 
where $\omega_E$ is a minimal differential on $E$ and $f=\sum_{n=1}^\infty a_nq^n\in S_2(\Gamma_0(N))$ is the  newform associated with $E$.  Pick generators $\delta^\pm$ of $H_1(E,\Z)^\pm$, and define the N\'{e}ron periods $\Omega_E^\pm$ by 
\[
\Omega_E^\pm=\int_{\delta^\pm}\omega_E.
\]
We normalize the $\delta^\pm$
%periods 
so that $\Omega_E^+\in\R_{>0}$ and $\Omega_E^-\in i\R_{>0}$.

The Fourier coefficient $a_p$ is a $p$-adic unit by hypothesis, and we let $\alpha_p$ be the $p$-adic unit root of $x^2-a_px+p$. 

\begin{thm}
\label{thm:MSD}
There exists an element $\Lcal^{\rm MSD}_p(E/\Q)\in\Lambda_\Q$ such that for any finite order character $\chi$ of $\Gamma_\Q$ of conductor $p^r$ with $r>0$, we have
\[
\Lcal^{\rm MSD}_p(E/\Q)(\chi)=\frac{p^r}{\tau(\overline{\chi})\alpha_p^r}\cdot\frac{L(E,\overline{\chi},1)}{\Omega_E^+},
\]
where $\tau(\overline{\chi})=\sum_{a\;{\rm mod}\;p^r}\overline{\chi}(a)e^{2\pi ia/p^r}$ is the Gauss sum, and
\[
\Lcal_p^{\rm MSD}(E/\Q)(1)=\bigl(1-\alpha_p^{-1}\bigr)^2\cdot\frac{ L(E,1)}{\Omega_E^+}.
\]
\end{thm}

\begin{proof}
The construction of $\Lcal_p^{\rm MSD}(E/\Q)$ as an element in $\Lambda_{\Q}\otimes\Q_p$ with the stated interpolation property is given in \cite[\S{9}]{M-SwD}. 
%Its integrality properties are studied in \cite{greenvats} and \cite{wuthrich-int}, among others. In particular, the 
The integrality of $\Lcal_p^{\rm MSD}(E/\Q)$ is shown in \cite[Prop.~3.7]{greenvats} in the case where $E[p]$ is irreducible as a $G_\Q$-module, and in \cite[Cor.~18]{wuthrich-int} in the reducible case. Since the latter integrality result and some of the ingredients in the proof will be important later, we briefly outline the argument. First, W\"uthrich shows the existence of an elliptic curve $E_\bullet/\Q$ in the isogeny class of $E$ with $\mathcal{L}_p^{\rm MSD}(E_\bullet/\Q)\in\Lambda_{\Q}$ and whose $p$-adic Tate module satisfies 
\[
T_pE_{\bullet}\simeq V_{\Z_p}(f)(1),
\] 
where $V_{\Z_p}(f)$ is the geometric lattice in the $p$-adic Galois representation $V_{\Q_p}(f)$ associated to $f$ considered in \cite[\S{8.3}]{kato-euler-systems}. Building on the theorem of Ferrero--Washington \cite{FW}, Kato's divisibility in the Iwasawa main conjecture for $E_\bullet$ ``without zeta functions''  \cite[Thm.~12.5(4)]{kato-euler-systems} is shown to be integral, from which the integrality of $\mathcal{L}_p^{\rm MSD}(E_\bullet/\Q)$ and of $\mathcal{L}_p^{\rm MSD}(E/\Q)$ itself follows by global duality and the invariance of Mazur's main conjecture under isogenies (see Prop. \ref{prop:isog-inv}).
\end{proof}

%\begin{rmk}
%The integrality properties of $\Lcal_p^{\rm MSD}(E/\Q)$ are studied in \cite{greenvats}
%\footnote{\textcolor{blue}{More precisely, for any $\eta\in(\Z/Mp\Z)^\times$ with $p\nmid M$ one can define $\Lcal_p^{\rm MSD}(E/\Q,\eta)$, and by \cite[Remark~3.9]{greenvats} under hypothesis (\ref{assumption}) doesn't hold one knows integrality of $\Lcal_p^{\rm MSD}(E/\Q,\eta)$ for odd $\eta$.}} 
%and \cite{wuthrich-int}.  In particular, by \cite[Corollary~18]{wuthrich-int} one knows the inclusion  
%\[
%\Lcal_p^{\rm MSD}(E/\Q)\in\Lambda_\Q.
%\] 
%This is shown in two steps. First, one shows the existence of $E_\bullet$ isogenous to $E$ with $\mathcal{L}_p^{\rm MSD}(E_\bullet/\Q)\in\Lambda_{\Q}$ and such that 
%\[
%T_pE_{\bullet}\simeq V_{\Z_p}(f)(1)
%\] 
%is the lattice considered by Kato \cite{kato-euler-systems}. Then, building on the theorem of Ferrero--Washington \cite{FW}, Kato's divisibility in the Iwasawa main conjecture for $E_\bullet$ ``without zeta functions'' is shown to be integral, from where the integrality of $\mathcal{L}_p^{\rm MSD}(E/\Q)$ (rather than just $E_\bullet$) follows from global duality and the invariance of the main conjecture under isogenies.
%\end{rmk}

\subsection{Two-variable $p$-adic $L$-function, I}
\label{subsec:L-2var-I} 

%({\color{blue} under revision})
In this section we recall the $p$-adic $L$-function constructed in \cite{PR-JLMS} following Hida's $p$-adic Rankin method \cite{hida-measure-I}, and explain the relation with the $p$-adic $L$-function of Mazur--Swinnerton-Dyer. 

Let $\Sigma$ be the set of infinity types of Hecke characters $\psi$ of $K$ for which $s=1$ is a critical value of the Rankin $L$-function $L(f/K,\psi,s)$. Following the notations and conventions in \cite[\S{6.1}]{LLZ-K}, the set $\Sigma$ decomposes as 
\[
\Sigma=\Sigma^{(1)}\cup\Sigma^{(2)}\cup\Sigma^{(2')},
\]
where $\Sigma^{(1)}=\{(0,0)\}$ (i.e., corresponding to characters $\psi$ of finite order), $\Sigma^{(2)}=\{(a,b)\colon a\leq -1, b\geq 1\}$, and $\Sigma^{(2')}=\{(b,a)\colon a\leq -1, b\geq 1\}$. Note that the regions $\Sigma^{(2)}$ and $\Sigma^{(2')}$ are interchanged by the involution $\psi\mapsto\psi^\tau$, where $\psi^\tau$ denotes the composition of $\psi$ with the non-trivial automorphism $\tau\in{\rm Gal}(K/\Q)$.


For $\psi$ a finite order Hecke character of $K$ of conductor $\mathfrak{c}$, denote by $\theta(\psi)$ the associated theta series, which is an eigenform of weight one and level $M=D_K\mathbf{N}(\mathfrak{c})$. As in \cite[p.\,457]{PR-GZ}, define the ``Artin root number'' of $\psi$ to be the complex number $W(\psi)$ with $\vert W(\psi)\vert=1$ given by
\[
\theta(\psi)\vert_1\left(\begin{smallmatrix}&-1\\M&\end{smallmatrix}\right)=-iW(\psi)\cdot\theta(\psi);
\]
then the completed $L$-function $\Lambda(\psi,s)=M^{s/2}(2\pi)^{-s}\Gamma(s)\sum_{\mathfrak{a}}\psi(\mathfrak{a})\mathbf{N}(\mathfrak{a})^{-s}$ satisfies the functional equation $\Lambda(\psi,s)=W(\psi)\Lambda(\overline{\psi},1-s)$. As before, let $f\in S_2(\Gamma_0(N))$ be the newform associated with $E$, and put
\[
H_p(f)=\Bigl(1-\frac{p}{\alpha_p^2}\Bigr)\Bigl(1-\frac{1}{\alpha_p^2}\Bigr).
\] 
%where $\alpha_p$ is the $p$-adic unit root of $x^2-a_px+p$.
Denote by $c_f\in\Z_p$ the \emph{congruence number} of $f$, as defined in \cite[\S{7}]{hida-congruences,ribet-modp}; this is such that $c_f^{-1}$ is a $\Z_p$-generator of the congruence ideal $I_f\subset\Q_p$ defined in $\S\ref{subsec:ERL}$.  
%(thus $c_f$ is a $p$-adic unit if and only if there are no non-trivial congruences between $f$ and other normalized cusp forms in $S_2(\Gamma_0(N))$.

\begin{thm}
\label{thm:PR}
There exists an element $\Lcal_p(f/K,\Sigma^{(1)})\in c_f^{-1}\Lambda_K$ such that for every finite order character $\psi$ of $\Gamma_K$ of conductor $v^m\bar{v}^n$ with $m+n>0$, we have
\[
\Lcal_p(f/K,\Sigma^{(1)})(\psi)=\frac{W(\psi)p^{(m+n)/2}}{\alpha_p^{m+n}H_p(f)}\cdot\frac{L(f/K,\overline{\psi},1)}{8\pi^2\langle f,f\rangle_N},
\]
where $W(\psi)$ is the Artin root number of $\psi$, and
\[
\Lcal_p(f/K,\Sigma^{(1)})(1)=
%\prod_{w\vert p}(1-\alpha_p^{-1}\psi^{-1}(w))(1-\alpha_p^{-1\psi(w))
\frac{(1-\alpha_p^{-1})^{4}}{H_p(f)}
\cdot\frac{L(f/K,1)}{8\pi^2\langle f,f\rangle_N},
\]
where $\langle f,g\rangle_N=\int_{\Gamma_0(N)\backslash\mathcal{H}}\overline{f(\tau)}g(\tau)dxdy$ is the Petersson inner product on $S_2(\Gamma_0(N))$.
\end{thm}

\begin{proof}
This is a reformulation of \cite[Thm.~1.1]{PR-GZ}. For our later use, we note that this is the same as the two-variable $p$-adic $L$-function 
\[
L_p(f,\boldsymbol{g})\in c_f^{-1}\Z_p\widehat\otimes\Lambda_{\boldsymbol{g}}\llbracket\Gamma_\Q\rrbracket
\]
obtained by specializing to $f$ the three-variable $p$-adic Rankin $L$-series in \cite[Thm.~7.7.2]{explicit}, where $\boldsymbol{g}$ is the CM Hida family introduced in the proof of Lemma~\ref{lem:integral} below. 
%in  \cite[Theorem~2.7.4]{explicit}  
%(see also \cite[Theorem~B.2.1]{BL-PR}, whose term $\tau(\psi)$ is given by $iW(\psi)p^{-(m+n)/2}$ in our conventions). 
\end{proof}

\begin{defi}[``Perrin-Riou's $p$-adic $L$-function'']\label{def:PR}
Let $f\in S_2(\Gamma_0(N))$ be the newform associated to $E$, and let $\pi_E:X_0(N)\rightarrow E$ be a modular parametrization. Then we
put
\[
\Lcal_p^{\rm PR}(E/K):=\frac{\deg(\pi_E)}{c_E^2}\cdot H_p(f)\cdot \Lcal_p(f/K,\Sigma^{(1)}),\nonumber
\]
where $c_E$ is the Manin constant associated to $\pi_E$.
\end{defi}


%{\color{blue} \begin{rem} (To be removed later) If $E$ is the strong Weil curve (i.e., optimal w.r.t. $X_0(N)$), then (since $p\nmid N$) we know ${\rm deg}(\pi)\sim_p c_f$, where $c_f$ is the congruence number of $f$ (\cite[Thm.~2.1]{ARS}); and in general since $p\nmid 2N$ we know $p\nmid c_E$ by \cite[Cor.~4.1]{mazur}. By construction, the maps ${\rm Col}^f$ in Theorem~\ref{thm:KLZ} is such that $c_f\cdot{\rm Col}^f$ takes values in $\Lambda_K$. \end{rem}}

\begin{rem}\label{rem:periods}
Note that the factor $H_p(f)$ can be interpreted as an Euler factor coming from the adjoint $L$-function of $f$ (see \cite[Rem.~6.5.10]{KLZ-AJM}). Moreover, setting 
\[
\Omega_{E/K}:=\frac{1}{\sqrt{\vert D_K\vert}}\int_{E(\C)}\omega_E\wedge i\,\overline{\omega_E}
\]
we find
\[
8\pi^2\langle f,f\rangle_N=\int_{X_0(N)}\omega_f\wedge i\,\overline{\omega_f}
=\frac{\deg(\pi_E)}{c_E^2}\cdot\Omega_{E/K},
\]
and so Theorem~\ref{thm:PR} says that $\Lcal_p^{\rm PR}(E/K)$ interpolates the finite order twists of $L(f/K,1)$ normalized by the complex period $\Omega_{E/K}$.
\end{rem}


Denote by $\Lcal^{\rm PR}_p(E/K)^+\in\Lambda_\Q$ the image of $\Lcal^{\rm PR}_p(E/K)$ under the map induced by the projection $\Gamma_K\twoheadrightarrow\Gamma_K^+\simeq\Gamma_\Q$.

\begin{prop}
\label{prop:comp-Lcyc}
Up to a unit in $\Lambda_\Q^\times$, we have 
\[
\Lcal_p^{\rm PR}(E/K)^+=\Lcal_p^{\rm MSD}(E/\Q)\cdot\Lcal^{\rm MSD}_p(E^K/\Q),
\] 
where $E^K$ is the twist of $E$ by the quadratic character corresponding to $K/\Q$.
\end{prop}

\begin{proof}
%Suppose first that $E/\Q$ is any elliptic curve in the isogeny class associated to $f$. 
Let $\chi$ be a Dirichlet character of conductor $p^r$. Using the standard formula 
\[
W(\chi\circ\mathbf{N})=\chi(\vert D_K\vert)\frac{\tau(\chi)^2}{p^r}=\chi(\vert D_K\vert)\frac{p^r}{\tau(\overline{\chi})^2}
\]
(see e.g. \cite[Lem.~4.8.1]{miyake}) and the factorization
\[
L(f/K,\overline{\chi}\circ\mathbf{N},1)=L(E,\overline{\chi},1)\cdot L(E^K,\overline{\chi},1),
\]
from Theorem~\ref{thm:MSD} and Theorem~\ref{thm:PR} we find
\[
\begin{aligned}
\Lcal_p(f/K,\Sigma^{(1)})(\chi\circ\mathbf{N})\cdot H_p(f)&=W(\chi\circ\mathbf{N})\cdot\frac{p^r}{\alpha_p^{2r}}\cdot L(f/K,\overline{\chi}\circ\mathbf{N},1)\\
&=\chi(\vert D_K\vert)\cdot\frac{p^{2r}}{\tau(\overline{\chi})\alpha_p^{2r}}\cdot L(E,\overline{\chi},1)\cdot L(E^K,\overline{\chi},1)\\
&=\chi(\vert D_K\vert)\cdot\Lcal_p^{\rm MSD}(E/\Q)(\chi)\cdot\Lcal_p^{\rm MSD}(E^K/\Q)(\chi)\cdot\frac{\Omega_E^+\cdot\Omega_{E^K}^+}{8\pi^2\langle f,f\rangle_N}.
\end{aligned}
\]
Since the complex period $\Omega_{E/K}$ in Remark~\ref{rem:periods} satisfies
\[
\Omega_{E/K}=[E(\R):E^0(\R)]\cdot\Omega_E^+\cdot\Omega_{E^K}^+
\]
%On the other hand, from the period relation
%\[
%8\pi\langle f,f\rangle_N=\frac{{\rm deg}(\pi_E)}{c_E^2}\cdot\Omega_E^+\cdot\Omega_{E^K}^+\cdot[E(\R):E^0(\R)]\cdot\sqrt{\vert D_K\vert}
%\]
(see \cite[\S{V.2}]{grosszagier}), the result in now clear from the definition of $\Lcal_p^{\rm PR}(E/K)^+$.
%we obtain 
%\begin{equation}\label{eq:period-rel}
%\frac{\Omega_E^+\cdot\Omega_{E^K}^+}{8\pi^2\langle f,f\rangle_N}\;\sim_p\;\frac{c_E^2}{{\rm deg}(\pi_E)}.
%\end{equation}
%
%Now, if $E=A$ is the strong Weil curve in the isogeny class associated to $f$ and $\pi_A:X_0(N)\rightarrow A$ is the minimal parametrization, then one knows that ${\rm deg}(\pi_A)\sim_p c_f$ (see \cite[Thm.~2.10]{ARS}, which in this case follows from \cite[Prop.~3.3, Prop.~3.4]{abbes-ullmo}); and the corresponding Manin constant $c_A$ is a $p$-adic unit by \cite[Cor.~4.1]{mazur}. Hence the result follows from Definition~\ref{def:PR}, (\ref{eq:interp}) and (\ref{eq:period-rel}).
\end{proof}

\subsection{Anticyclotomic $p$-adic $L$-function}
\label{subsec:L-ac} 

We next recall the ``square-root'' $p$-adic $L$-function of Bertolini--Darmon--Prasanna \cite{BDP}, explicitly shown to be a $p$-adic measure in \cite{cas-hsieh1}. Assume that 
\begin{equation}\label{eq:Heeg}
\textrm{every prime $\ell\vert N$ splits in $K$,}\tag{Heeg}
\end{equation}
and fix an integral ideal $\mathfrak{N}\subset\mathcal{O}_K$ with $\mathcal{O}_K/\mathfrak{N}=\Z/N\Z$.  For simplicity, we also assume that
\begin{equation}\label{eq:disc}
\textrm{the discriminant $D_K<0$ is odd and $D_K\neq -3$}.\tag{disc}
\end{equation}

In the following we let $\Omega_p$ and $\Omega_K$ be CM periods attached to $K$ as in \cite[\S{2.5}]{cas-hsieh1} and put  $\Lambda_{K}^{-,\rm ur}=\Lambda_K\widehat\otimes\Z_p^{\rm ur}$, where $\Z_p^{\rm ur}$ is the completion of the ring of integers of the maximal unramified extension of $\Q_p$.

\begin{thm}
\label{thm:BDP}
There exists an element $\Lcal_p^{\rm BDP}(f/K)\in\Lambda_K^{-,{\rm ur}}$ characterized by the following interpolation property: for every character $\xi$ of $\Gamma_K^-$ crystalline at both $v$ and $\bar{v}$ and corresponding to a Hecke character of $K$ of infinity type $(n,-n)$ with $n\in\Z_{>0}$ and $n\equiv 0\pmod{p-1}$, we have
\[
\Lcal_p^{\rm BDP}(f/K)(\xi)=\frac{\Omega_p^{4n}}{\Omega_K^{4n}}\cdot\frac{\Gamma(n)\Gamma(n+1)\xi(\mathfrak{N}^{-1})}{4(2\pi)^{2n+1}\sqrt{D_K}^{2n-1}}\cdot\bigl(1-a_p\xi(\bar{v})p^{-1}+\xi(\bar{v})^2p^{-1}\bigr)^2\cdot L(f/K,\xi,1).
\]  
Moreover, $\Lcal_p^{\rm BDP}(f/K)$ is a nonzero element of $\Lambda_K^{-,{\rm ur}}$.
\end{thm}

\begin{proof}
See \cite[Thm.~2.1.1]{eisenstein} for the construction, 
which is deduced from \cite[\S{3}]{cas-hsieh1}. Since (\ref{eq:Heeg}) implies that $f$ does not have CM by $K$, the nonvanishing of $\Lcal_p^{\rm BDP}(f/K)$ follows from \cite[Thm.~3.9]{cas-hsieh1}.
\end{proof}

\begin{rmk}\label{rem:periods}
The CM period $\Omega_{K}\in\C^\times$ in Theorem~\ref{thm:BDP} agrees with that in \cite[(5.1.16)]{BDP}, but is \emph{different} from the period $\Omega_\infty$ defined in \cite[p.\,66]{deshalit} and \cite[(4.4b)]{hidatilouineI}. In fact, one has 
\[
\Omega_{\infty}=2\pi i\cdot\Omega_K.
\]
In terms of $\Omega_\infty$, the interpolation formula in Theorem~\ref{thm:BDP} reads
\[
\Lcal_{p}^{\rm BDP}(f/K)(\xi)=\frac{\Omega_p^{4n}}{\Omega_\infty^{4n}}\cdot\frac{\Gamma(n)\Gamma(n+1)\xi(\mathfrak{N}^{-1})}{4(2\pi)^{1-2n}\sqrt{D_K}^{2n-1}}\cdot\bigl(1-a_p\xi(\overline{v})p^{-1}+\xi(\bar{v})^2p^{-1}\bigr)^2\cdot L(f/K,\xi,1).
\]
This is the form of the interpolation that we shall use later.
\end{rmk}



\subsection{Two-variable $p$-adic $L$-function, II}
\label{subsec:L-2var-II} 

The main result of this section is Proposition~\ref{prop:comp-Lac}, relating the %anticyclotomic 
$p$-adic $L$-function $\Lcal_p^{\rm BDP}(f/K)$ of Theorem~\ref{thm:BDP} to the anticyclotomic projection of the following two-variable $p$-adic Rankin $L$-series. 


\begin{thm}\label{thm:Hida-2}
\label{thm:Gr}
There exists an element $\Lcal_p(f/K,\Sigma^{(2')})\in{\rm Frac}\,\Lambda_K$ such that for every character $\xi$ of $\Gamma_K$ 
%of conductor dividing $\mathfrak{f}$, 
crystalline at both $v$ and $\bar{v}$, and of infinity type $(b,a)$ with $a\leq -1$ and $b\geq 1$, we have
\[
\Lcal_p(f/K,\Sigma^{(2')})(\psi)=\frac{2^{a-b}i^{b-a-1}\Gamma(b+1)\Gamma(b)N^{a+b+1}}{(2\pi)^{2b+1}\langle\theta_{\psi_b},\theta_{\psi_b}\rangle_N}\cdot\frac{\mathcal{E}(\psi,f,1)}{(1-\psi^{1-\tau}(\bar{v}))(1-p^{-1}\psi^{1-\tau}(\bar{v}))}\cdot L(f/K,\psi,1),
\]
where $\theta_{\psi_b}$ is the theta series of weight $b-a+1\geq 3$ associated to the Hecke character $\psi_b=\psi\vert\;\vert^{-b}$ of $\infty$-type $(0,a-b)$, and 
\[
\mathcal{E}(\psi,f,1)=(1-p^{-1}\psi(\bar{v})\alpha_p)(1-\psi(\bar{v})\alpha^{-1}_p)(1-\psi^{-1}({v})\alpha_p^{-1})(1-p^{-1}\psi^{-1}({v})\alpha_p).
\]
\end{thm}

\begin{proof}
This is another instance of Hida's $p$-adic Rankin $L$-series, as explained in \cite[Thm.~6.1.3]{LLZ-K} (note, however, that we have reversed the roles of $v$ and $\bar{v}$ with respect to \emph{loc.\,cit.}).
\end{proof}

We also need to recall the interpolation property of the Katz $p$-adic $L$-functions \cite{katz}, following the formulation in \cite{deshalit}. Put $\Lambda_K^{\rm ur}=\Lambda_K\widehat\otimes_{\Z_p}\Z_p^{\rm ur}$.

\begin{thm}\label{thm:katz}
There exists an element $\Lcal_{\bar{v}}(K)\in\Lambda_K^{\rm ur}$ such that for every character $\xi$ of $\Gamma_K$ of infinity type $(k,j)$ with $0\leq-j<k$ satisfies 
\[
\Lcal_{\bar{v}}(K)(\xi)=\frac{\Omega_p^{k-j}}{\Omega_\infty^{k-j}}\cdot\Gamma(k)\cdot\biggl(\frac{\sqrt{D_K}}{2\pi}\biggr)^j\cdot(1-\xi^{-1}(\bar{v})p^{-1})(1-\xi({v}))\cdot L(\xi,0).
\]
Similarly, there exists an element $\Lcal_v(K)\in\Lambda_K^{\rm ur}$ such that for every character $\xi$ of $\Gamma_K$ of infinity type $(j,k)$ with $0\leq-j<k$, we have 
\[
\Lcal_{v}(K)(\xi)=\frac{\Omega_p^{k-j}}{\Omega_\infty^{k-j}}\cdot\Gamma(k)\cdot\biggl(\frac{\sqrt{D_K}}{2\pi}\biggr)^j\cdot(1-\xi^{-1}({v})p^{-1})(1-\xi(\bar{v}))\cdot L(\xi,0).
\]
Moreover, we have the functional equation 
\[
\Lcal_{\bar{v}}(\xi)=\Lcal_{v}(\xi^{-1}\mathbf{N}^{-1}),
\] 
where the equality is up to a $p$-adic unit.
\end{thm}

\begin{proof}
This is \cite[Thm.~II.4.14]{deshalit}, with $\Lcal_{\bar{v}}(K)$ (resp. $\Lcal_{{v}}(K)$) corresponding to the measure $\mu(\bar{v}^\infty)$ (resp. $\mu(v^\infty)$) in \emph{loc.\,cit.}. On the other hand, as explained in \cite[Lem.~3.3.2(b)]{7author}, the stated functional equation is a reformulation of \cite[Thm.~II.6.4]{deshalit}.
\end{proof}


\begin{defi}[``Greenberg's $p$-adic $L$-function'']\label{def:Gr}
Put
\[
\Lcal_p^{\rm Gr}(f/K):=h_K\cdot\Lcal_v(K)^-\cdot\Lcal_p(f/K,\Sigma^{(2')}),
\]
where $h_K$ is the class number of $K$ and $\Lcal_v(K)^-$ the image of $\Lcal_v(K)$ under the map $\Lambda_K^{\rm ur}\rightarrow\Lambda_K^{\rm ur}$ given by $\gamma\mapsto\gamma^{1-\tau}$ for $\gamma\in\Gamma_K$. 
%and $w_K=\vert\mathcal{O}_K^\times\vert$. 
\end{defi}

Note that \emph{a priori} we have $\Lcal_p^{\rm Gr}(f/K)\in{\rm Frac}\,\Lambda_K^{\rm ur}$. 

\begin{lemma}\label{lem:integral}
The $p$-adic $L$-function $\Lcal_p^{\rm Gr}(f/K)$ is integral, i.e., $\Lcal_p^{\rm Gr}(f/K)\in\Lambda_K^{\rm ur}$.
\end{lemma}

\begin{proof}
Denote by $\eta=\eta_{K/\Q}$ the quadratic character corresponding to $K/\Q$, and let ${\rm Eis}_{1,\eta}(q)$ be the weight one Eisenstein series
\[
{\rm Eis}_{1,\eta}(q)=\sum_{n\geq 1}q^n\sum_{d\vert n}\eta_{}(d).
\]
Since $p$ splits in $K$, ${\rm Eis}_{1,\eta}(q)$ is $p$-irregular. Letting $g$ denote the unique $p$-stabilization of ${\rm Eis}_{1,\eta}$, by \cite[Thm.~A(i)]{Betina-Dimitrov-Pozzi} there is a unique cuspidal Hida family $\boldsymbol{g}$ passing through $g$. Moreover, $\boldsymbol{g}$ is of CM type, given explicitly as the $q$-series
\[
\boldsymbol{g}=\sum_{(\mathfrak{a},\bar{v})=1}[\mathfrak{a}]q^{\mathbf{N}(\mathfrak{a})}\in\Lambda_{\boldsymbol{g}}\llbracket q\rrbracket,
\]
where $\Lambda_{\boldsymbol{g}}=\Z_p\llbracket\Gamma_v\rrbracket$, for $\Gamma_v$ the Galois group of the maximal $\Z_p$-extension inside $K_\infty/K$ unramified at $v$, and where $[\mathfrak{a}]$ denotes the natural image of the ideal $\mathfrak{a}\subset\mathcal{O}_K$ in $\Gamma_v$ under the Artin reciprocity map. 

Now, the $p$-adic $L$-function $\mathcal{L}_p(f/K,\Sigma^{(2')})$ in Theorem~\ref{thm:Hida-2} arises from Hida's $p$-adic Rankin $L$-series 
\[
L_p(\boldsymbol{g},f)\in I_{\boldsymbol{g}}^{\rm cusp}\widehat\otimes_{\Z_p}\Z_p\llbracket\Gamma_\Q\rrbracket,
\] 
where $I_{\boldsymbol{g}}^{\rm cusp}\subset{\rm Frac}\,\Lambda_{\boldsymbol{g}}$ is the cuspidal congruence ideal of $\boldsymbol{g}$ (see $\S\ref{subsec:ERL}$ for the definition).  
Thus if $H_{\bg}^{\rm cusp}\in\Lambda_{\boldsymbol{g}}$ denotes a characteristic power series for the denominator of  $I_{\boldsymbol{g}}$, then the product $H_{\bg}^{\rm cusp}\cdot\Lcal_p(f/K,\Sigma^{(2)})$ is integral, so it suffices to show that $h_K\cdot\Lcal_v(K)^-$ is divisible by $H_{\bg}^{\rm cusp}$. 

By \cite[Thm.~0.3]{hidatilouineII} and Rubin's proof of the Iwasawa main conjecture for $K$, \cite{rubinmainconj},  one has that such divisibility holds up to powers of the augmentation ideal $(\gamma_v-1)\subset\Z_p\llbracket\Gamma_v\rrbracket$; since by \cite[Thm.~A(i)]{Betina-Dimitrov-Pozzi} one knows that $H_{\bg}^{\rm cusp}$ is not divisible by $\gamma_v-1$, the result follows.
\end{proof}

Denote by $\Lcal_p^{\rm Gr}(f/K)^-$ the image of $\Lcal_p^{\rm Gr}(f/K)$ under the natural projection $\Lambda_K^{\rm ur}\rightarrow\Lambda_K^{-,{\rm ur}}$.

\begin{prop}\label{prop:factor-L-Gr}
\label{prop:comp-Lac}
We have the equality
%Up to a unit, we have the equality
\[
\Lcal_p^{\rm Gr}(f/K)^-\cdot \Lambda_K^{-,\rm{ur}}=\Lcal_p^{\rm BDP}(f/K)\cdot \Lambda_K^{-,\rm{ur}}.
\]
\end{prop}

\begin{proof}
This follows from a direct comparison of the interpolation formulas in Theorem~\ref{thm:Gr}, Theorem~\ref{thm:BDP} and Theorem~\ref{thm:katz}, together with an application of Dirichlet's class number formula (cf. \cite[Thm.~1.7]{castellaheights}). 

Indeed, let $\xi$ be a Hecke character of infinity type $(n,-n)$, $n\in\Z_{\geq 0}$, as in the statement of Theorem~\ref{thm:BDP}. Then the character $\xi^{1-\tau}\mathbf{N}^{-1}$, of infinity type $(2n+1,1-2n)$, is in the range of interpolation of $\Lcal_{\bar{v}}(K)$, and noting that $L(\xi^{1-\tau}\mathbf{N}^{-1},0)=L(\xi^{1-\tau},1)$, by Theorem~\ref{thm:katz} we have
\begin{equation}\label{eq:ev-Katz}
\Lcal_{\bar{v}}(K)(\xi^{1-\tau}\mathbf{N}^{-1})=\frac{\Omega_p^{4n}}{\Omega_\infty^{4n}}\cdot\Gamma(2n+1)\cdot\biggl(\frac{2\pi}{\sqrt{D_K}}\biggr)^{2n-1}\cdot(1-\xi^{1-\tau}(\bar{v}))(1-\xi^{1-\tau}({v})p^{-1})\cdot L(\xi^{1-\tau},1).
\end{equation}
On the other hand, from Hida's formula for the adjoint $L$-value (see \cite[Thm~.7.1]{hidatilouineI}) and Dirichlet's class number formula we obtain %(cf. \cite[p.\,414]{JSW})
\begin{equation}\label{eq:classnumber}
\langle \theta_{\xi_n},\theta_{\xi_n}\rangle_N\;\sim_p\;
\Gamma(2n+1)\cdot\frac{1}{2^{4n-1}\pi^{2n+1}}\cdot h_K\cdot L(\xi^{1-\tau},1),
\end{equation}
where $\sim_p$ denotes equality up to $p$-adic unit independent of $n$, and similarly as in Theorem~\ref{thm:Hida-2}, $\xi_n$ is the theta series of weight $2n+1\geq 3$ associated to the Hecke character $\xi_n=\xi\vert\;\vert^{-n}$ of infinity type $(0,-2n)$. Combining (\ref{eq:ev-Katz}), (\ref{eq:classnumber}) and the functional equation in Theorem~\ref{thm:katz} this gives
\[
h_K\cdot\Lcal_{v}(K)(\xi^{1-\tau})\;\sim_p\;\frac{\Omega_p^{4n}}{\Omega_\infty^{4n}}\cdot\biggl(\frac{2\pi}{\sqrt{D_K}}\biggr)^{2n-1}\cdot(1-\xi^{1-\tau}(\bar{v}))(1-\xi^{1-\tau}({v})p^{-1})\cdot 2^{4n-1}\pi^{2n+1}\cdot \langle\theta_{\xi_n},\theta_{\xi_n}\rangle_N.
\]
Noting that the $p$-Euler factor $\mathcal{E}(\psi,f,1)$ in Theorem~\ref{thm:Hida-2} satisfies
\[
\mathcal{E}(\xi,f,1)=\bigl(1-a_p\xi(\overline{v})p^{-1}+\xi(\bar{v})^2p^{-1}\bigr)^2,
\]
from Definition~\ref{def:Gr}, Theorem~\ref{thm:Hida-2}, and Theorem~\ref{thm:BDP} we thus find that
\[
\Lcal_p^{\rm Gr}(f/K)(\xi)\;\sim_p\;\xi(\mathfrak{N})\cdot 2^{3n-2}i^{2n-1}\cdot\Lcal_p^{\rm BDP}(f/K)(\xi).
\]
Since $\xi(\mathfrak{N})\cdot 2^{3n-2}i^{2n-1}$ is interpolated by a unit in $\Lambda_K^{-,\rm{ur}}$, this completes the proof.
\end{proof}

\subsection{Twists and imprimitive $p$-adic $L$-functions}\label{subsec:tw-L}

%We conclude this section by introducing twisted %and imprimitive  
%variants of some of the above $p$-adic $L$-functions. 

Let $\alpha:\Gamma_K\rightarrow R^\times$ be a character with values in the ring of integers $R$ of a finite extension $\Phi/\Q_p$ with uniformiser $\varpi\in R$. Let $\Lambda_{K,R} = R\widehat\otimes_{\Z_p}\Lambda_K = R[\![\Gamma_K]\!]$, and define
\[
{\rm Tw}_\alpha:\Lambda_{K,R}\rightarrow\Lambda_{K,R}
\]
to be the $R$-linear isomorphism given by $\gamma\mapsto\alpha(\gamma)\gamma$ for $\gamma\in\Gamma_K$. Denote by 
$\Lcal_p^{\rm PR}(E(\alpha)/K), \Lcal_p^{\rm Gr}(f(\alpha)/K)$ the image of $\Lcal_p^{\rm PR}(E/K), \Lcal_p^{\rm Gr}(f/K)$, respectively, under ${\rm Tw}_\alpha$. %Let $\varpi\in R$ be a uniformizer. 

\begin{lemma}\label{lem:cong-L}
Suppose $\alpha\equiv 1\pmod{\varpi^m}$. Then 
\[
\Lcal_p^{\rm PR}(E(\alpha)/K)^+\equiv\Lcal_p^{\rm PR}(E/K)^+\;({\rm mod}\,\varpi^m).
\]
\end{lemma}

\begin{proof}
This is clear from the definitions.
%This is clear, since letting $\gamma^+\in\Gamma^+_K$ be a topological generator, the hypothesis implies $(\alpha(\gamma^+)\gamma^_-1)^n\equiv(\gamma-1)^n\pmod{\varpi^m}$ for all $n>0$, 
%
%This is clear from Lemma~\ref{lem:integral} and the definition of $, since composing $\alpha$ with the projection $\Gamma_K\twoheadrightarrow\Gamma^-$, we may view $\Lcal_{\rm Gr}(E(\alpha)/K_\infty)$ as the image of $\Lcal_{\rm Gr}(E/K_\infty)$ under the map $\Lambda_K^{\rm ur}\to\Lambda_K^{\rm ur}$ given by $\gamma\mapsto\alpha(\gamma)\gamma$ for $\gamma\in\Gamma_K$.
\end{proof}


%\subsection{Imprimitive $p$-adic $L$-functions} 

For a prime $w$ split in $K$, lying over the rational prime $\ell\neq p$, we let $\Gamma_w^\pm$ be the corresponding decomposition group in $\Gamma_K^\pm$, and $\gamma_w^\pm\in\Gamma_w^\pm$ be the image of an arithmetic Frobenius ${\rm Frob}_w$ under the projection $G_K\rightarrow\Gamma_K^\pm$. 

\begin{defi}\label{def:w-Euler}
Put 
\[
\mathcal{P}^\pm_w(\alpha):=P_w(\ell^{-1}\gamma_w^\pm)\in\Lambda_{K,R}^{\pm},
\]
where $P_w(X)=\det(1-{\rm Frob}_wX\,\vert\,V(\alpha)_{I_w})$ is the Euler factor at $w$ of the $L$-function of $V(\alpha)=T_pE(\alpha)\otimes\Q_p$. For $S'$ a finite set of primes $w$ as above, define
\begin{align*}
\Lcal_{p}^{\rm PR}(E(\alpha)/K)^{+,S'}&:=\Lcal_{p}^{\rm PR}(E(\alpha)/K)^+\cdot\prod_{w\in S'}\mathcal{P}_w^+(\alpha),\\\Lcal_p^{\rm Gr}(f(\alpha)/K)^{\pm,S'}&:=\Lcal_p^{\rm Gr}(f(\alpha)/K)^\pm\cdot\prod_{w\in S'}\mathcal{P}^\pm_w(\alpha),
\end{align*}
and similarly $\Lcal_p^{\rm BDP}(f(\alpha)/K)^{S'}:=\Lcal_p^{\rm BDP}(f(\alpha)/K)\cdot\prod_{w\in S'}\mathcal{P}^-_w(\alpha)$.
\end{defi}

Of course, the results of Proposition~\ref{prop:comp-Lcyc} and Proposition~\ref{prop:comp-Lac} directly extend to their analogues for these $S'$-imprimitive $p$-adic $L$-functions.

\section{Selmer groups}\label{sec:Sel}

In this section, we let $E/\Q$ be an elliptic curve of conductor $N$, $p$ an odd  prime of good ordinary reduction for $E$, and $K$ an imaginary quadratic field satisfying (\ref{eq:Heeg}) and (\ref{eq:spl}). 
%Let $\Lambda_\Q$, $\Lambda_K$, $\Lambda_K^\pm$ be the Iwasawa algebras as introduced in $\S\ref{sec:Lp}$.


\subsection{Selmer structures}\label{subsec:Sel-str}

Let $\Sigma$ be a finite set of places of $\Q$ containing the prime $p$, $\infty$, and the prime factors of $N$. We assume throughout that 
\[
\textrm{all finite primes in $\Sigma$ split in $K$.}
\]
With a slight abuse of notation, we 
%shall 
also write $\Sigma$ 
%(resp. $S'$ for any subset $S'\subset\Sigma$) 
for the set of places of $K$ lying above the places in $\Sigma$. 
%and similarly for any subset $S'\subset\Sigma$.


\subsubsection{Discrete coefficients}\label{subsec:discrete-Sel}

For a discrete $\Z_p$-module $M$, we let
\[
M^\vee={\rm Hom}_{\rm cts}(M,\Q_p/\Z_p)
\]
be the Pontryagin dual. The module $\Lambda_{\Q}^{\vee}$ is equipped with a $G_\Q$-action via $\Psi^{-1}$, where $\Psi:G_\Q\rightarrow\Lambda_{\Q}^\times$ is the character arising from the projection $G_\Q\twoheadrightarrow\Gamma_\Q$. Similarly, $\Lambda_K^{\vee}$ and $(\Lambda^\pm_K)^{\vee}$ are equipped with $G_K$-actions.

\begin{defi}\label{defilocalconds}
Let $F$ be $\Q$ or $K$ and $w$ a prime above $p$. For $\Lambda$ any of the Iwasawa algebras $\Lambda_\Q$, $\Lambda_K$, or $\Lambda_K^\pm$, 
we put
%over $\Q$ or $K$, respectively,  we let 
\begin{align*}
\rH^1_{\rm rel}(F_w,T_pE\otimes \Lambda^{\vee})&=\rH^1(F_w,T_pE\otimes \Lambda^{\vee}),\\
\rH^1_{\ord}(F_w,T_pE\otimes \Lambda^{\vee})&= {\rm im}\bigl\{\rH^1(F_w,{\rm Fil}_w^+(T_pE)\otimes \Lambda^{\vee})\rightarrow\rH^1(F_w,T_pE\otimes \Lambda^{\vee})\bigr\},\\
\rH^1_{\rm str}(F_w,T_pE\otimes \Lambda^{\vee})&=\{0\}, 
\end{align*}
where ${\rm Fil}_w^+(T_pE):={\rm ker}\bigl\{T_pE\rightarrow T_p\tilde{E}\bigr\}$ with $\tilde{E}$ the reduction of $E$ at $w$.
\end{defi}
Let $G_{\Q,\Sigma}$ and $G_{K,\Sigma}$ denote the Galois group of the maximal extension of $\Q$ and $K$ respectively unramified outside $\Sigma$. %We consider the following Selmer structures: 
For $\bullet \in \{\ord, \rm str, \rm rel \}$ and $M=T_pE\otimes \Lambda_\Q^{\vee}$, we define the Selmer group
\begin{equation}
\rH^1_{\Fcal_{\bullet}}(\Q, M)=\ker \biggl(\rH^1(G_{\Q,\Sigma},M) \to \prod_{w\in\Sigma,w\nmid p}\rH^1(\Q_w, M)\times\frac{\rH^1(\Q_{p},M)}{\rH^1_{\bullet}(\Q_{p},M)}\biggr).
\end{equation}
Similarly, for $\star,\bullet \in \{\ord, \rm str, \rm rel \}$ and $M=T_pE\otimes\Lambda^{\vee}$, where $\Lambda$ is any of the Iwasawa algebras $\Lambda_K$ or $\Lambda_K^\pm$, we let
\begin{equation}\label{eq:defSelK}
\rH^1_{\Fcal_{\star,\bullet}}(K, M)=\ker \biggl(\rH^1(G_{K,\Sigma},M) \to \prod_{w\in\Sigma,w\nmid p}\rH^1(K_w, M)\times\frac{\rH^1(K_v,M)}{\rH^1_{\star}(K_v,M)}\times\frac{\rH^1(K_{\bar{v}},M)}{\rH^1_{\bullet}(K_{\bar{v}},M)}\biggr).
\end{equation}
To ease notation, we write $\rH^1_{\Fcal_{\ord}}(K, M)=\rH^1_{\Fcal_{\ord,\ord}}(K, M)$ and $\rH^1_{\Fcal_{\rm Gr}}(K, M)=\rH^1_{\Fcal_{\rm rel, str}}(K, M)$, and put
\begin{align*}
\X_{\rm ord}(E/\Q_{\infty})&=\rH^1_{\Fcal_{\ord}}(\Q, T_pE\otimes\Lambda_\Q^{\vee})^{\vee},\\
\X_{\rm ord}(E/K_\infty^\pm)&=\rH^1_{\Fcal_{\ord}}(K, T_pE\otimes (\Lambda^\pm_K)^{\vee})^{\vee},\\
\X_{\rm Gr}(E/K_\infty^\pm)&=\rH^1_{\Fcal_{\rm Gr}}(K, T_pE\otimes (\Lambda^\pm_K)^{\vee})^{\vee}.
\end{align*}
It is a standard fact that these are finitely generated modules over the corresponding Iwasawa algebras. 

\subsubsection{Compact coefficients} 

For $\Lambda$ any of the Iwasawa algebras $\Lambda_\Q$, $\Lambda_K$, or $\Lambda_K^\pm$, consider the compact module
\[
T_pE\widehat\otimes_{\Z_p}\Lambda,
\]
where $\Lambda$ is equipped with a $G_K$-action via $\Psi:G_K\rightarrow\Lambda^\times$. For $\bullet\in\{\rm ord, \rm str, \rm rel\}$ and $w$ a prime of $K$ above $p$, we define the local conditions $\rH^1_{\bullet}(K_w,T_pE\widehat\otimes_{\Z_p}\Lambda)\subset\rH^1(K_w,T_pE\widehat\otimes_{\Z_p}\Lambda)$ similarly as in Definition~\ref{defilocalconds}, and for $\star,\bullet\in\{\rm ord,\rm str,\rm rel\}$ we define the Selmer group $\rH^1_{\Fcal_{*,\bullet}}(K,T_pE\widehat\otimes_{\Z_p}\Lambda)$ by the same recipe as in (\ref{eq:defSelK}). Put
\[
\mathfrak{S}_{\rm ord,\rm rel}(E/K_\infty)=\rH^1_{\Fcal_{\rm ord,\rm rel}}(K,T_pE\widehat\otimes_{\Z_p}\Lambda_K),\quad
\mathfrak{S}_{\rm ord}(E/K_\infty)=\rH^1_{\Fcal_{\rm ord,\rm ord}}(K,T_pE\widehat\otimes_{\Z_p}\Lambda_K),
\]
etc., and similarly for $E/K_\infty^\pm$ with $\Lambda_K$ in place of $\Lambda_K^\pm$, respectively.



\subsection{Iwasawa main conjectures}

The Iwasawa main conjecture for $E$, as formulated by Mazur \cite{mazur-towers}, \cite[\S{9.5}, Conj.~3]{M-SwD}, is the following. 

\begin{conjecture}[Mazur]\label{conj:IMC}
The module $\X_{\rm ord}(E/\Q_{\infty})$ is $\Lambda_\Q$-torsion, with 
\[
\ch_{\Lambda_\Q}\bigl(\X_{\rm ord}(E/\Q_{\infty})\bigr)=\bigl(\Lcal_p^{\rm MSD}(E/\Q)\bigr).
\]
\end{conjecture}

More generally, a vast generalization of Mazur's main conjecture to $p$-adic deformations of motives formulated by Greenberg \cite{greenberg-reps,greenberg-motives},
predicts the following.

\begin{conjecture}[Greenberg]\label{conj:IMC-K}
The modules $\X_{\rm ord}(E/K^+_\infty)$ and $\X_{\rm Gr}(E/K_\infty^-)$ are torsion over $\Lambda_K^+$ and $\Lambda_K^-$ respectively, with
\begin{align*}
\ch_{\Lambda_K^+}\bigl(\X_{\rm ord}(E/K_{\infty}^+)\bigr)&=\bigl(\Lcal_p^{\rm PR}(E/K)^+\bigr),\\
\ch_{\Lambda_K^-}\bigl(\X_{\rm Gr}(E/K_{\infty}^-)\bigr)\Lambda_K^{-,\rm ur}&=\bigl(\Lcal_p^{\rm Gr}(f/K)^-\bigr).
\end{align*}
\end{conjecture}

%Note that under hypothesis (\ref{eq:Heeg}) the anticyclotomic restriction $\Lcal_p^{\rm PR}(E/K)^-$ vanishes identically, and $\X_{\rm ord}(E/K_\infty^-)$ can be shown to have positive $\Lambda_K^-$-rank. 

In this paper we shall prove Mazur's Main Conjecture~\ref{conj:IMC} (in the case where $p$ is a good Eisenstein prime for $E$) by first proving  Conjecture~\ref{conj:IMC-K} for a suitable $K$. 
%and then using the following result to descend from $K$ to $\Q$. 
%As done here, we shall often identify the Iwasawa algebras $\Lambda_K^+$ and $\Lambda_\Q$ (via the natural projection $\Lambda_K^+\stackrel{\sim}{\rightarrow}\Gamma_\Q$).

%\begin{prop}\label{prop:comp-Selcyc}
%The restriction map from $G_\Q$ to $G_K$ induces a $\Lambda_\Q$-module isomorphism 
%\[
%\X_{\rm ord}(E/K_{\infty}^+)\simeq \X_{\rm ord}(E/\Q_\infty)\oplus\X_{\rm ord}(E^K/\Q_\infty),
%\]
%where $E^K$ is the twist of $E$ by the quadratic character corresponding to $K/\Q$. In particular, 
%\[
%\ch_{\Lambda^+_K}\bigl(\X_{\rm ord}(E/K_{\infty}^+)\bigr)=\ch_{\Lambda_{\Q}}\bigl(\X_{\rm ord}(E/\Q_\infty)\bigr)\cdot \ch_{\Lambda_{\Q}}\bigl(\X_{\rm ord}(E^K/\Q_\infty)\bigr).
%\]
%\end{prop}

%\begin{proof}
%This follows readily from the inflation-restriction exact sequence and %Shapiro's lemma (see e.g. \cite[Lem.\,3.6]{skinner-urban}).
%
%By \cite{kato-euler-systems}, the $\Lambda_{\Q}$-modules $\X(E/\Q)$, $\X(E^K/\Q)$ are finitely generated and torsion.
%By Shapiro's lemma together with $\Lambda^+_K\simeq \Lambda_{\Q}$ (see for example \cite[Lemma 3.6]{skinner-urban}), we find
%$$\rH^1_{\Fcal_{\ord}}(K, T_p(E)\otimes (\Lambda^+_K)^{\vee})\simeq \rH^1_{\Fcal_{\ord}}(\Q, T_p(E)\otimes (\Lambda_{\Q})^{\vee})\oplus \rH^1_{\Fcal_{\ord}}(\Q, T_p(E^K)\otimes (\Lambda_{\Q})^{\vee}).$$
%From this we deduce the second claim and the statement that $\X(E/K_{\infty}^+)$ is a torsion $\Lambda^+_K$-module. The non-existence of non-zero finite $\Lambda^+_K$-modules follows from the non-existence of non-zero finite $\Lambda_{\Q}$-module for $\X(E/\Q)^S$ and $\X(E^K/\Q)^S$, which is proved in \cite[Proposition 2.5]{greenvats}.
%\end{proof}



\subsubsection{Isogeny invariance}

Conjectures~\ref{conj:IMC} and \ref{conj:IMC-K} are known to be invariant under isogenies. This follows from a computation in global duality due to Schneider and Perrin-Riou \cite{schneider-isogenies,perrinriou}. 


%For the latter one, this invariance is immediate from the following result, where we let $\Fcal_{\rm ord}(E/K_\infty^+)\in\Lambda_K^+$ denote a characteristic power series for $\X_{\rm ord}(E/K_{\infty}^+)$ (taken to be zero if the module is non-torsion).

\begin{prop}\label{prop:isog-inv}
Suppose $E_1/\Q$ and $E_2/\Q$ are isogenous elliptic curves with good ordinary reduction at $p$. Assume that $\mathfrak{X}_{\rm ord}(E_i/K_\infty^\pm)$ is $\Lambda_K^+$-torsion ($i=1,2$), and let $\mathcal{F}_{\rm ord}(E_i/K_\infty^+)$ and $\mathcal{F}_{\rm ord}(E_i/\Q_\infty)$ be characteristic power series for $\mathfrak{X}_{\rm ord}(E_i/K_\infty^+)$ and $\mathfrak{X}_{\rm ord}(E/\Q_\infty)$, respectively. Then we have the equalities up to a $p$-adic unit:
\begin{align*}
\Omega_{E_1}\cdot\Fcal_{\rm ord}(E_1/\Q_\infty)&\;\sim_p\;\Omega_{E_2}\cdot\Fcal_{\rm ord}(E_2/\Q_\infty),\\
\Omega_{E_1/K}\cdot\Fcal_{\rm ord}(E_1/K_\infty^+)&\;\sim_p\;\Omega_{E_2/K}\cdot\Fcal_{\rm ord}(E_2/K_\infty^+).
\end{align*} 
In particular, the main conjectures for $\mathfrak{X}_{\rm ord}(E/\Q_\infty)$ and $\mathfrak{X}_{\rm ord}(E/K_{\infty}^+)$ are both invariant under isogenies.
\end{prop}

\begin{proof}
See \cite[Appendice]{perrinriou}.
\end{proof}

Although not directly needed for our arguments, we also note that the isogeny invariance of the main conjecture for $\mathfrak{X}_{\rm Gr}(E/K_\infty^-)$ similarly follows from the main result of \cite{PR-isogenies} (see \cite[Prop.\,2.9]{KO-iwasawa}).


\subsection{Imprimitive Selmer groups}\label{subsec:imp}

For any subset $S'\subset\Sigma$ consisting of primes away from $p$, we define $S'$-imprimitive versions of the  Selmer groups of $\S\ref{subsec:discrete-Sel}$ by relaxing the local conditions at the primes $w\in S'$, e.g. for $M=T_pE\otimes\Lambda_\Q^\vee$:
\[
\rH^1_{\Fcal_{\rm ord}^{S'}}(\Q,M)=\ker \biggl(\rH^1(G_{\Q,\Sigma},M) \to \prod_{w\in\Sigma\smallsetminus S',w\nmid p}\rH^1(\Q_w, M)\times\frac{\rH^1(\Q_{p},M)}{\rH^1_{\rm ord}(\Q_{p},M)}\biggr).
\] 
%In particular, for $S:=\Sigma\smallsetminus\{w\mid p\infty\}$, $M$ a $G_{\Q}$- or $G_{K}$-module as above, and $\bullet, \star\in \{\ord, \rm str, \rm rel \}$, we have 
%\[
%\rH^1_{\Fcal_{\bullet}^S}(\Q, M)=\ker \left(\rH^1(G_{\Q,\Sigma},M) \to \rH^1_{\bullet}(\Q_{p},M)\right),
%\]
%\[
%\rH^1_{\Fcal_{\star,\bullet}^S}(K, M)=\ker\biggl(\rH^1(G_{K,\Sigma},M) \to \frac{\rH^1(K_v,M)}{\rH^1_{\star}(K_v,M)}\times \frac{\rH^1(K_{\bar{v}},M)}{\rH^1_{\bullet}(K_{\bar{v}},M)}\biggr).
%\]
%
We denote with a superscript $S'$ the Pontryagin duals of these modules:
\begin{align*}
\X_{\rm ord}^{S'}(E/\Q_{\infty})&=\rH^1_{\Fcal_{\ord}^{S'}}(\Q, T_pE\otimes \Lambda_\Q^{\vee})^{\vee},\\
\X_{\rm ord}^{S'}(E/K^\pm_\infty)&=\rH^1_{\Fcal_{\ord}^{S'}}(K, T_pE\otimes (\Lambda^\pm_K)^{\vee})^{\vee},\\
\X_{\rm Gr}^{S'}(E/K_\infty^\pm)&=\rH^1_{\Fcal_{\rm Gr}^{S'}}(K, T_pE\otimes (\Lambda^\pm_K)^{\vee})^{\vee}.
\end{align*}

The next result will be used to descend from $K$ to $\Q$ (\emph{cf.} Proposition~\ref{prop:comp-Lcyc}). %we shall use the following basis result. 
As done here, in the following we shall often identify the Iwasawa algebras $\Lambda_K^+$ and $\Lambda_\Q$ (via the natural projection $\Lambda_K^+\stackrel{\sim}{\rightarrow}\Gamma_\Q$).


\begin{prop}\label{prop:comp-Selcyc}
Let $S'\subset\Sigma$ be any subset of primes not lying above $p$. Then the restriction map from $G_\Q$ to $G_K$ induces a $\Lambda_\Q$-module isomorphism 
\[
\X_{\rm ord}^{S'}(E/K_{\infty}^+)\simeq \X_{\rm ord}^{S'}(E/\Q_\infty)\oplus\X_{\rm ord}^{S'}(E^K/\Q_\infty).
\]
%where $E^K$ is the twist of $E$ by the quadratic character corresponding to $K/\Q$. 
In particular, 
\[
\ch_{\Lambda^+_K}\bigl(\X_{\rm ord}^{S'}(E/K_{\infty}^+)\bigr)=\ch_{\Lambda_{\Q}}\bigl(\X_{\rm ord}^{S'}(E/\Q_\infty)\bigr)\cdot \ch_{\Lambda_{\Q}}\bigl(\X_{\rm ord}^{S'}(E^K/\Q_\infty)\bigr).
\]
%Moreover, if $\phi_{G_{\Q_p}}\neq 1,\omega$, then
%$\X(E/K_{\infty}^+)$ is a finitely generated torsion $\Lambda^+_K$-module and $\X^S(E/K_{\infty}^+)$ has no non-zero finite $\Lambda^+_K$-submodule.
\end{prop}

\begin{proof}
This follows readily from the inflation-restriction exact sequence and Shapiro's lemma (see e.g. \cite[Lem.\,3.6]{skinner-urban}).
%
%By \cite{kato-euler-systems}, the $\Lambda_{\Q}$-modules $\X(E/\Q)$, $\X(E^K/\Q)$ are finitely generated and torsion.
%By Shapiro's lemma together with $\Lambda^+_K\simeq \Lambda_{\Q}$ (see for example \cite[Lemma 3.6]{skinner-urban}), we find
%$$\rH^1_{\Fcal_{\ord}}(K, T_p(E)\otimes (\Lambda^+_K)^{\vee})\simeq \rH^1_{\Fcal_{\ord}}(\Q, T_p(E)\otimes (\Lambda_{\Q})^{\vee})\oplus \rH^1_{\Fcal_{\ord}}(\Q, T_p(E^K)\otimes (\Lambda_{\Q})^{\vee}).$$
%From this we deduce the second claim and the statement that $\X(E/K_{\infty}^+)$ is a torsion $\Lambda^+_K$-module. The non-existence of non-zero finite $\Lambda^+_K$-modules follows from the non-existence of non-zero finite $\Lambda_{\Q}$-module for $\X(E/\Q)^S$ and $\X(E^K/\Q)^S$, which is proved in \cite[Proposition 2.5]{greenvats}.
\end{proof}





%\begin{prop}\label{propSelGrpm}
%Assume $\phi\vert_{G_p}\neq 1,\omega$. Then $\X_{\rm Gr}^S(E/K_{\infty}^-)$ and $\X_{\rm Gr}(E/K_{\infty}^-)$ are finitely generated torsion $\Lambda^-_K$-modules and $\X_{\rm Gr}^S(E/K_{\infty}^-)$ has no non-zero finite $\Lambda^-_K$-submodules.
%\end{prop}

%\begin{proof}
%The statement is proved in \cite[Proposition 1.4.2, Corollary 1.4.3]{eisenstein}, characterising the Selmer group $\rH^1_{\Fcal_{\rm Gr}}(K, T_p(E)\otimes (\Lambda^+_K)^{\vee})$ in terms of Selmer groups of the characters $\phi, \phi^{-1}\omega$ and then applying the results of \cite{rubinmainconj} relating the latter with the anticyclotomic projection of a Katz $p$-adic $L$-function. The assumptions on the character $\phi$ are used to ensure that $E[p](K_{\bar{v}})=0$ and therefore $H^0(K_{\bar{v}},T_p(E)\otimes (\Lambda^+_K)^{\vee})=0$. This allows to reduce the proof of the torsionness statement to studying the finiteness of the Selmer group of the $p$-torsion of $T_p(E)\otimes (\Lambda^+_K)^{\vee}$, which up to semi-simplification is given by a direct sum of the two characters. This then implies the triviality of finite $\Lambda^-_K$-submodules (\cite[Corollary 1.4.3]{eisenstein}).
%\end{proof}

\subsubsection{From imprimitive to primitive}

As observed by Greenberg, imprimitive Selmer groups as above tend to have better properties with respect to congruences than their primitive counterparts. For our arguments, we shall also find it convenient to work first with imprimitive Selmer group, and so the next results will be useful.


\begin{prop}\label{prop:prim-imprim-ord}
Assume that $E(K)[p]=0$ and that $\X_{\rm ord}(E/K_\infty^+)$ is $\Lambda_K^+$-torsion. Then for any $S'\subset\Sigma$ consisting of primes away from $p$, the Selmer group $\X_{\rm ord}^{S'}(E/K_\infty^+)$ is also $\Lambda_K^+$-torsion, with
\[
\ch_{\Lambda_K^+}\bigl(\X_{\rm ord}^{S'}(E/K_\infty^+)\bigr)=
\ch_{\Lambda_K^+}\bigl(\X_{\rm ord}(E/K_\infty^+)\bigr)\cdot\prod_{w\in S'}\bigl(\mathcal{P}_w^+(1)\bigr),
\]
where $\mathcal{P}_w^+(1)\in\Lambda_K^+$ is as in Definition~\ref{def:w-Euler}, with $\alpha=1$.
\end{prop}

\begin{proof}
From the assumption that $E(K)[p]=0$, we see that the $G_{K_\infty^+}$-invariants of ${\rm Hom}(T_pE,\mu_{p^\infty})$ are trivial, and so by \cite[Prop.~A.2]{pollack2011anticyclotomic} the global-to-local map defining $\rH^1_{\Fcal_{\rm ord}}(K,T_pE\otimes(\Lambda_K^+)^\vee)$ as in (\ref{eq:defSelK}) is surjective. We therefore find an exact sequence
\begin{equation}\label{eq:seq-imp-ord}
0\to\prod_{w\in S'}\rH^1(K_w,T_pE\otimes(\Lambda_K^+)^\vee)^\vee\to\X_{\rm ord}^{S'}(E/K_\infty^+)\to\X_{\rm ord}(E/K_\infty^+)\to 0.
\end{equation}
Since the primes $w\in S'$ split in $K$, by  \cite[Prop.~2.4]{greenvats} the module $\rH^1(K_w,T_pE\otimes(\Lambda_K^+)^\vee)^\vee$ is $\Lambda_K^+$-torsion, with
\[
\ch_{\Lambda_K^+}\bigl(\rH^1(K_w,T_pE\otimes(\Lambda_K^+)^\vee)^\vee\bigr)=\bigl(\mathcal{P}_w^+(1)\bigr).
\]
%(note that since $w$ split in $K$, it is finitely decomposed in $K_\infty/K$), 
The result now follows by taking characteristic ideals in (\ref{eq:seq-imp-ord}).
\end{proof}




\begin{cor}\label{cor:imp-IMC} 
Assume that $E(K)[p]=0$ and let $S'\subset\Sigma$ be any subset consisting of primes away from $p$. Then Conjecture~\ref{conj:IMC-K} for $\mathfrak{X}_{\rm ord}(E/K_\infty^+)$ holds if and only if $\mathfrak{X}_{\rm ord}^{S'}(E/K_\infty^+)$ is $\Lambda_K^+$-torsion, with
\[
\ch_{\Lambda_K^+}\bigl(\X_{\rm ord}^{S'}(E/K_{\infty}^+)\bigr)=\bigl(\Lcal_p^{\rm PR}(E/K)^{+,S'}\bigr).
\]
\end{cor}

\begin{proof}
Since for any prime $w\in S'$, the element $\mathcal{P}_w^+(1)\in\Lambda_K^+$ is nonzero, this is clear from Definition~\ref{def:w-Euler} and Proposition~\ref{prop:prim-imprim-ord}.
\end{proof}

%Let $S=\Sigma\smallsetminus\{w\mid p\infty\}$ and 
%Assume $E(K)[p]=0$ and that $\X_{\rm ord}(E/K_\infty^+)$ is $\Lambda_K^+$-torsion. Then
%$$\ch_{\Lambda^+_K}\bigl(\X_{\rm ord}(E/K_{\infty}^+)\bigr)=\ch_{\Lambda_{\Q}}\bigl(\X_{\rm ord}(E/\Q_\infty)\bigr)\cdot \ch_{\Lambda_{\Q}}\bigl(\X_{\rm ord}(E^K/\Q_\infty)\bigr).$$
%\end{cor}

%\begin{proof}
%Immediate from Propositions~\ref{prop:comp-Selcyc} and \ref{prop:prim-imprim-ord}, since $\mathcal{P}_w^+(1)$ ($w\in S$) is a nonzero element of $\Lambda_K^+$. Note that we are also implicitly using the analogue of Proposition~\ref{prop:comp-Selcyc} for $\X^S_{\rm ord}(E/\Q_\infty)$ and $\X^S_{\rm ord}(E^K/\Q_\infty)$ (see \cite[Cor. 2.3]{greenvats}) and the fact that every rational prime in $S$ splits in $K$.
%\end{proof}


\subsection{Anticyclotomic twists and congruences}\label{sec:seltwist}

We now introduce the twisted variants of the Selmer groups from the preceding sections that we shall need. The results of Proposition~\ref{prop:euler-char} and Proposition~\ref{prop:cong-Sel} will play an important role later.


%From now one we assume that $p$ is an Eisenstein prime for $E$, i.e. $E$ admits a rational $p$-isogeny. Let us write the semi-simplification of the $G_{\Q}$-module $E[p]$ as 
%\[
%E[p]^{ss}=\mathbb{F}_p(\phi)\oplus\mathbb{F}(\phi^{-1}\omega)
%\] 
%where $\phi:G_\Q\rightarrow\mathbb{F}_p^\times$ is a character, and $\omega$ is the Teichm\"uller character.

%We shall also need to consider twisted versions of the Selmer groups introduced in $\S\ref{subsec:Sel-str}$. Thus 
As in $\S\ref{subsec:tw-L}$, let $\Phi$ be a finite extension of $\Q_p$, and let $R$ be the ring of integers of $\Phi$ with uniformizer $\varpi$. We consider a character $\alpha:\Gamma_K^-\to R^{\times}$ which satisfies
\[
\alpha\equiv 1\;({\rm mod}\,\varpi^m)
\]
for some $m>0$. Let $S'\subset\Sigma$ be any subset consisting of primes away from $p$. Replacing $T_pE$ by the twist 
\[
T_\alpha:=T_pE\otimes_{\Z_p}R(\alpha),
\]  
we define (imprimitive) Selmer groups $\rH^1_{\Fcal_{\rm ord}^{S'}}(K,T_\alpha\otimes(\Lambda_K^\pm)^\vee)$ and $\rH^1_{\Fcal_{\rm Gr}^{S'}}(K_,T_\alpha\otimes(\Lambda_K^\pm)^\vee)$ (with Pontryagin duals $\X^{S'}_{\rm ord}(E(\alpha)/K_\infty^\pm)$ and $\X_{\Gr}^{S'}(E(\alpha)/K_\infty^\pm)$, respectively)  
in the same way as before, replacing ${\rm Fil}_w^+(T_pE)$ by ${\rm Fil}_w^+(T_pE)\otimes_{\Z_p}R(\alpha)$ in Definition \ref{defilocalconds}. We also consider their counterparts with $K_\infty^\pm$ and $\Lambda_K^\pm$ replaced by $K_\infty$ and $\Lambda_K$, respectively. 
Each of these Selmer groups is a module for the Iwasawa algebra $\Lambda_{R}= R\widehat\otimes_{\Z_p}\Lambda$ with $\Lambda$ either $\Lambda_\Q$, $\Lambda_K$, or $\Lambda_K^\pm$ (per the definitions).

Put $V_\alpha=T_\alpha\otimes\Q_p$ and $W_\alpha=T_\alpha\otimes\Q_p/\Z_p=V_\alpha/T_\alpha$. We shall also need to consider the following variant of the above Selmer groups for the module $W_\alpha:=T_\alpha\otimes\Q_p/\Z_p$:
\[
\rH^1_{\Fcal_{\rm Gr}^{S'}}(K,V_\alpha):=\ker\Biggl(\rH^1(G_{K,\Sigma},V_\alpha)\rightarrow
%\frac{\rH^1(K_v,V_\alpha)}{\rH^1(K_v,W_\alpha)_{\rm div}}
\prod_{w\in\Sigma\smallsetminus S', w\nmid p}\rH^1(K_w,V_\alpha)\times\rH^1(K_{\bar{v}},V_\alpha)\Biggr),
\]
%where $\rH^1(K_v,W_\alpha)_{\rm div}$ denotes the maximal divisible submodule of $\rH^1(K_v,W_\alpha)$. 
%(Note that this corresponds to the 
and the resulting $\rH^1_{\Fcal_{\rm Gr}^{S'}}(K,T_\alpha)$ and 
$\rH^1_{\Fcal_{\rm Gr}^{S'}}(K,W_\alpha)$ obtained by propagation via $0\rightarrow T_\alpha\rightarrow V_\alpha\rightarrow W_\alpha\rightarrow 0$. 
%Note that these are the same as the Selmer group $\rH^1_{\Fcal_{\rm ac}^{S'}}(K,V_\alpha)$, etc. in \cite[\S{2.3}]{JSW}, but with the roles of $v$ and $\bar{v}$ reversed. 
We also consider the %Bloch--Kato 
Selmer group $\rH^1_{\Fcal_{\rm ord}}(K,V_\alpha)$ consisting of classes that land in 
\[
%\rH^1_{\rm BK}(K_w,V_\alpha):={\rm ker}(\rH^1(K_w,V_\alpha)%\rightarrow\rH^1(K_w,V_\alpha\otimes\mathbf{B}_{\rm cris}))
%
\rH^1_{\rm ord}(K_w,V_\alpha):={\rm im}(\rH^1(K_w,{\rm Fil}_w^+V_\alpha)\rightarrow\rH^1(K_w,V_\alpha))
\] 
at the primes $w\mid p$ and are trivial at the primes $w\nmid p$, and its counterparts $\rH^1_{\Fcal_{\rm ord}}(K,T_\alpha)$ and $\rH^1_{\Fcal_{\rm ord}}(K,W_\alpha)$ defined by propagating the local conditions.

\subsubsection{Ancillary results for $\X_{\rm Gr}(E(\alpha)/K_\infty^\pm)$}
The main result of this section is a relation between the specializations  of the cyclotomic Selmer group $\X_{\rm Gr}(E(\alpha)/K_\infty^+)$ and the anticyclotomic Selmer group $\X_{\rm Gr}(E(\alpha)/K_\infty^-)$ at the trivial character. 
%The way we prove this is adapting the arguments of \cite[$\S$3]{jsw} to prove control theorems for both these Selmer groups and realise they yield the same result.

Define
\[
M_\alpha^\pm=T_\alpha\widehat\otimes_{\Z_p}(\Lambda_K^\pm)^\vee,
\]
and for any $S'\subset\Sigma$ consisting of primes not dividing $p$, put
\[
\mathcal{P}_{\rm Gr}(M_\alpha;S'):=\rH^1(K_{\bar{v}},M_\alpha)\times
\prod_{w\in\Sigma\smallsetminus S', w\nmid p}\rH^1(K_w,M_\alpha),
\]
and similarly,
\[
\mathcal{P}_{\Gr}(W_\alpha;S'):=\frac{\rH^1(K_v,W_\alpha)}{\rH^1(K_v,W_\alpha)_{\rm div}}\times\rH^1(K_{\bar{v}},W_\alpha)\times
\prod_{w\in\Sigma\smallsetminus S', w\nmid p}\rH^1(K_w,W_\alpha),
\]
so in particular $\rH^1_{\Fcal_{\Gr}^{S'}}(K,W_\alpha)$ is the kernel of the global-to-local map $\rH^1(G_{K,\Sigma},W_\alpha)\rightarrow\mathcal{P}_{\Gr}(W_\alpha;S')$.


%Let $G_p\subset G_\Q$ be a decomposition group at $p$.

%\begin{prop}\label{prop:prim-imprim}
%Assume that $\phi\vert_{G_{p}}\neq 1,\omega$ and that $\X_{\rm Gr}(E(\alpha)/K_\infty^\pm)$ is $\Lambda_{K,R}^\pm$-torsion. Let $S'\subset\Sigma$ be any subset of $\Sigma$ consisting of primes away from $p$. Then $\X_{\rm Gr}^{S'}(E(\alpha)/K_\infty^\pm)$ is also $\Lambda_{K,R}^\pm$-torsion, and 
%\[
%\ch_{\Lambda_K^\pm}\bigl(\X_{\rm Gr}^{S'}(E(\alpha)/K_\infty^\pm)\bigr)=
%\ch_{\Lambda_K^\pm}\bigl(\X_{\rm Gr}(E(\alpha)/K_\infty^\pm)\bigr)\cdot\prod_{w\in S'}\bigl(\mathcal{P}_w^\pm(\alpha)\bigr),
%\]
%where $\mathcal{P}_w^\pm(\alpha)\in\Lambda_{K,R}^\pm$ is as in Definition~\ref{def:w-Euler}.
%\end{prop}

%\begin{proof}
%The same argument as in the proof of Proposition~\ref{prop:prim-imprim-ord} applies, noting that $\rH^0(K,T_{\alpha}[\varpi])=0$ since $E(K)[p]=0$ and $\alpha\equiv 1\,({\rm mod}\,\varpi)$.
%
%The proof follows the same lines of the one for Proposition \ref{prop:prim-imprim-ord}. Since $\phi\vert_{G_{\Q_p}}\neq 1,\omega$, the $G_{K_\infty}$-invariants of ${\rm Hom}(T_\alpha,\mu_{p^\infty})$ are trivial, and so by \cite[Prop.~A.2]{pollack2011anticyclotomic} the global-to-local map cutting out $\rH^1_{\Fcal_{\rm Gr}}(K,T_\alpha\otimes(\Lambda_K^+)^\vee)$ as in (\ref{eq:defSelK}) is surjective. From this, we immediately obtain the exact sequence
%\begin{equation}\label{eq:seq-imp}
%0\to\prod_{w\in S'}\rH^1(K_w,T_\alpha\otimes(\Lambda_K^+)^\vee)^\vee\to\X_{\rm Gr}^{S'}(E(\alpha)/K_\infty^+)\to\X_{\rm Gr}(E(\alpha)/K_\infty^+)\to 0.
%\end{equation}
%Since by \cite[Prop.~2.4]{greenvats} the module $\rH^1(K_w,T_\alpha\otimes(\Lambda_K^+)^\vee)^\vee$ is $\Lambda_K^+$-torsion with
%\[
%\ch_{\Lambda_K^+}\bigl(\rH^1(K_w,T_\alpha\otimes(\Lambda_K^+)^\vee)^\vee\bigr)=\bigl(\mathcal{P}_w^+(\alpha)\bigr),
%\]
%taking characteristic ideals in (\ref{eq:seq-imp}) the result follows.
%\end{proof}

%Recall that we put $S=\Sigma\smallsetminus\{p,\infty\}$.

%\begin{prop}\label{prop:ps-null}
%Assume that $\phi\vert_{G_{p}}\neq 1,\omega$ and that 
%$\alpha:\Gamma^-\rightarrow R^\times$ is a character with $\alpha\equiv 1\pmod{\varpi^N}$ for some $N>0$ and such that 
%$\X_{\rm Gr}^S(E(\alpha)/K_\infty^\pm)$ is $\Lambda^\pm_K$-torsion. Then $\X_{\rm Gr}^S(E(\alpha)/K_\infty^\pm)$ has no non-zero finite $\Lambda^\pm_K$-submodules.
%\end{prop}

%\begin{proof}
%As in the proof of Proposition~\ref{prop:prim-imprim-ord}, the assumptions imply that the global-to-local map in $(\ref{eq:defSelK})$ is surjective, yielding the isomorphism
%\begin{equation}\label{eq:lem26}
%\X_{\rm Gr}^S(E(\alpha)/K_\infty^\pm)\simeq\frac{\rH^1(G_{K,\Sigma},M_{\alpha}^\pm)^\vee}{\rH^1(K_{\bar{v}},M_{\alpha}^\pm)^\vee},
%\end{equation}
%where $M_\alpha=T_\alpha\otimes(\Lambda_K^\pm)^\vee$. Since the same argument as in the proof of Proposition~1.2.5 and Lemma~1.3.1 in \cite{eisenstein} shows that $\rH^1(G_{K,\Sigma},M_{\alpha}^\pm)^\vee$ has no non-zero finite $\Lambda^\pm_K$-submodules and $\rH^1(K_{\bar{v}},M_{\alpha})$ is $\Lambda^\pm_K$-cofree, respectively, the result follows from (\ref{eq:lem26}) and \cite[Lem.~2.6]{greenvats}.
%\end{proof}

%The next auxiliary result 
%, which will play a  particularly important role later, shows that the order of the Greenberg Selmer groups for the cyclotomic and the anticyclotomic $\Z_p$-extensions of $K$ agree when specialised at the trivial character. 




\begin{lemma}\label{lem:coinv}
Assume $E(K)[p]=0$ and $\alpha:\Gamma_K^-\rightarrow R^\times$ is a crystalline character such that:
\begin{itemize}
\item[(a)] ${\rm corank}_{R}\rH^1_{\Fcal_{\rm BK}}(K,W_{\alpha^{-1}})=1$,
\item[(b)] The restriction map
\[
\rH^1_{\Fcal_{\rm ord}}(K,W_{\alpha^{-1}})\xrightarrow{{\rm loc}_v}\rH^1_{\rm ord}(K_v,W_{\alpha^{-1}})
\]
is nonzero. 
\end{itemize}
Then the following hold:
\begin{itemize}
\item[(i)] $\rH^1_{\Fcal_{\rm Gr}}(K,W_\alpha)$ is finite, and  $\rH^1_{\Fcal_{\rm Gr}}(K,T_\alpha)=0$.
\item[(ii)] $\rH^1(G_{K,\Sigma},M_\alpha^\pm)_{\Gamma_K^\pm}=0$.
\item[(iii)] For any $S'\subset\Sigma$ consisting of primes away from $p$, $\rH^1_{\Fcal_{\rm Gr}^{S'}}(K,M_\alpha^\pm)_{\Gamma_K^\pm}=0$.
\end{itemize}
\end{lemma}

\begin{proof}
%Since $W_\alpha=%T_pE(\alpha)\otimes(\Lambda_K^\pm)^\vee
%M_\alpha^\pm[\gamma^\pm-1]$, this is shown in the proof of
It follows from their definition by propagation that $\rH^1_{\Fcal_{\rm Gr}}(K,T_\alpha)$ is the $p$-adic Tate module of $\rH^1_{\Fcal_{\rm Gr}}(K,W_\alpha)$, and so part (i) is shown in Proposition\,3.2.1 of \cite{jsw}. For the proof of parts (ii) and (iii),  we shall adapt the arguments in the proof of [\emph{op.\,cit.}, Lem.\,3.3.5]. 
%We explain some of the details for the convenience of the reader. 
Let $\gamma^\pm\in\Gamma_K^{\pm}$ be any topological generator. From the relation $W_\alpha=%T_pE(\alpha)\otimes(\Lambda_K^\pm)^\vee
(M_\alpha^\pm)^{\Gamma_K^\pm}=M_\alpha^\pm[\gamma^\pm-1]$ we get an injection 
\begin{equation}\label{eq:inj-h2}
\rH^1(G_{K,\Sigma},M_\alpha^\pm)_{\Gamma_K^\pm}\hookrightarrow\rH^2(G_{K,\Sigma},W_\alpha).\nonumber
\end{equation}
Since the $p$-adic representation $V_{\alpha^{-1}}$ is pure of weight $-1$, we have $\rH^2(K_w,W_\alpha)=\rH^0(K_w,T_{\alpha^{-1}})^\vee=0$ for all $w\in\Sigma$, and so $\rH^2(G_{K,\Sigma},W_\alpha)=\sha_\Sigma^2(K,W_\alpha)$ %(with notations as in \cite[\S{4}]{milne-ADT}), 
which by Poitou--Tate duality is dual to $\sha_\Sigma^1(K,T_{\alpha^{-1}})$. %\simeq\sha_\Sigma^1(K,T_\alpha)$, using the action of complex conjugation for the last isomorphism. 
However, from $E(K)[p]=0$ the group $\sha_\Sigma^1(K,T_{\alpha^{-1}})$ is torsion-free, and from our assumption on $\alpha$ it is also of $\Z_p$-corank $0$. Hence $\sha_\Sigma^1(K,T_{\alpha^{-1}})=0$, so also $\rH^2(G_{K,\Sigma},W_\alpha)=0$, and the vanishing of $\rH^1(G_{K,\Sigma},M_\alpha^\pm)_{\Gamma_K^\pm}$ follows. This shows part (ii), and the vanishing of $\rH^1_{\Fcal_{\rm Gr}^{S'}}(K,M_\alpha^\pm)_{\Gamma_K^\pm}$ for any $S'\subset\Sigma$ as in part (iii) then follows from the exact sequence
\[
\rH^1(G_{K,\Sigma},W)=\rH^1(G_{K,\Sigma},M_\alpha^\pm)^{\Gamma_K^\pm}\xrightarrow{\lambda}\mathcal{P}_{\rm Gr}(M_\alpha^{\pm};S')^{\Gamma_K^\pm}\rightarrow\rH^1_{\Fcal_{\rm Gr}^{S'}}(K,M_\alpha^\pm)_{\Gamma_K^\pm}\rightarrow \rH^1(K^\Sigma/K,M_\alpha^\pm)_{\Gamma_K^\pm},
\]
in which the map $\lambda$ is surjective, being the composition of the restriction map $\rH^1(G_{K,\Sigma},W_\alpha)\rightarrow\mathcal{P}_{\Gr}(W_\alpha;S')$ (whose cokernel naturally injects into the dual of $\rH^1_{\Fcal_{\Gr}^{S'}}(K,T_{\alpha})=0$) and the natural map $\mathcal{P}_{\Gr}(W_\alpha;S')\rightarrow\mathcal{P}_{\Gr}(M_\alpha^{\pm};S')^{\Gamma_K^\pm}$, whose surjectivity follows from the fact that for $w\in\Sigma$, the local Galois group ${\rm Gal}(K_{\infty,\eta}^\pm/K_w)$ (for any $\eta\vert w$ in $K_\infty^\pm$) is either trivial or isomorphic to $\Z_p$, and so has $p$-cohomological dimension $\leq 1$.
\end{proof}

As standard, we shall often identify $\Lambda_K^\pm$ with the one variable power series ring $\Z_p[\![T^\pm]\!]$ via $\gamma^\pm=1+T^\pm$ upon the choice of a topological generator $\gamma^\pm\in\Lambda_K^\pm$. Thus denoting by $\mathcal{L}\mapsto\mathcal{L}^\pm$ the natural projection $\Lambda_K\rightarrow\Lambda_K^\pm$, the equality 
\[
\Lcal_p^{\rm PR}(E(\alpha)/K)^{+,S'}(0)=\Lcal_p^{\rm PR}(E(\alpha)/K)^{-,S'}(0)
\] 
is clear, reflecting in particular the fact that the trivial character is both cyclotomic and anticyclotomic. The next  result (which we shall only need for $S'=\emptyset$) is a parallel equality for characteristic power series. 

\begin{prop}\label{prop:euler-char}
Suppose $E(K)[p]=0$ and $\alpha:\Gamma_K^-\rightarrow R^\times$ is such that the conditions in Lemma~\ref{lem:coinv} hold.
Then the Selmer groups $\mathfrak{X}_{\rm Gr}^{S'}(E(\alpha)/K_\infty^+)$ and $\mathfrak{X}^{S'}_{\rm Gr}(E(\alpha)/K_\infty^-)$ are torsion over $\Lambda_K^+$ and $\Lambda_K^-$, respectively, where $S'\subset\Sigma$ is any subset consisting of primes away from $p$. 
%Let $\mathcak{F}_{\rm Gr}(E(\alpha)/K_\infty^\pm)\in\Lambda_{K,R}^\pm$ be a generator of the characteristic ideal of $\mathfrak{X}_{\rm Gr}(E(\alpha)/K_\infty^+)$, and assume that 
%
%Assume that $\phi\vert_{G_{p}}\neq 1,\omega$, and that the modules
%$\X_{\rm Gr}^S(E(\alpha)/K_\infty^\pm)$ are both $\Lambda^\pm_K$-torsion. 
Furthermore, we have the equality up to a $p$-adic unit:
\[
\mathcal{F}_{\rm Gr}^{S'}(E(\alpha)/K_\infty^+)(0)\;\sim_p\;\mathcal{F}_{\rm Gr}^{S'}(E(\alpha)/K_\infty^-)(0),
\]
where $\Fcal_{\rm Gr}^{S'}(E(\alpha)/K_\infty^\pm)\in\Lambda^\pm_{K,R}$ is any characteristic power series for $\X_{\rm Gr}^{S'}(E(\alpha)/K_\infty^\pm)$.
\end{prop}

\begin{proof} 
%Let $\gamma^\pm\in\Gamma_K^\pm$ be a topological generators.
By Poitou--Tate duality, the cokernel of the global-to-local map in the defining exact sequence
\[
0\rightarrow\rH^1_{\Fcal_{\Gr}}(K,W_\alpha)\rightarrow\rH^1(G_{K,\Sigma},W_\alpha)\rightarrow\mathcal{P}_{\Gr}(W_\alpha;\emptyset)
\]
injects into the Pontryagin dual of the Selmer group $\rH^1_{\Fcal_{\rm Gr}^*}(K,T_{\alpha^{-1}})$ dual to $\rH^1_{\Fcal_{\rm Gr}}(K,W_\alpha)$, defined by the local conditions given by the orthogonal complement of those in $\mathcal{P}_{\rm Gr}(W_\alpha;\emptyset)$ under local Tate duality. Since this dual Selmer group is torsion-free by the assumption $E(K)[p]=0$, and it follows from part (i) of Lemma~\ref{lem:coinv} and the finiteness of $\rH^1(K_w,T_{\alpha^{-1}})$ for finite primes $w\nmid p$  that it has finite order, we conclude that the above global-to-local map is surjective. It is then immediate that for any $S'$ as in the statement, we have an analogous exact sequence
\[
0\rightarrow\rH^1_{\Fcal_{\Gr}^{S'}}(K,W_\alpha)\rightarrow\rH^1(G_{K,\Sigma},W_\alpha)\rightarrow\mathcal{P}_{\Gr}(W_\alpha;S')\rightarrow 0.
\] 
A variant of Mazur's control theorem shows that the natural restriction map
\[
\rH^1_{\Fcal_{\Gr}^{S'}}(K,W_\alpha)\rightarrow
\rH^1_{\Fcal_{\rm Gr}^{S'}}(K,M_\alpha^\pm)^{\Gamma_K^\pm}
\]
has finite kernel and cokernel. By Lemma~\ref{lem:coinv} and the above remarks, it follows that $\mathfrak{X}_{\rm Gr}^{S'}(E(\alpha)/K_\infty^\pm)$ is $\Lambda_K^\pm$-torsion (for both choices of sign $\pm$), and together with the general result \cite[Lem.\,4.2]{greenberg-cetraro} for the $\Gamma_K^\pm$-Euler characteristic of $\mathcal{F}_{\rm Gr}^{S'}(E(\alpha)/K_\infty^\pm)$ we obtain the equality up to a $p$-adic unit:
\begin{equation}\label{eq:Euler-char}
\mathcal{F}_{\Gr}^{S'}(E(\alpha)/K_\infty^\pm)(0)\,\sim_p\;\#\rH^1_{\Fcal_{\Gr}^{S'}}(K,M_\alpha^\pm)^{\Gamma_K^\pm}.
\end{equation}
Now, from the Snake Lemma applied to the commutative diagram with exact rows
\[
\xymatrix{
0\ar[r]&\rH^1_{\Fcal_{\Gr}^{S'}}(K,W_\alpha)\ar[r]\ar[d]^{s^\pm}&\rH^1(K^\Sigma/K,W_\alpha)\ar[r]\ar[d]^{h^\pm}&\mathcal{P}_{\Gr}(W_\alpha;S')\ar[r]\ar[d]^{r^\pm}&0\\
0\ar[r]&\rH^1_{\Fcal_{\rm Gr}^{S'}}(K,M_\alpha^\pm)^{\Gamma_K^\pm}\ar[r]&\rH^1(K^\Sigma/K,M_\alpha^\pm)^{\Gamma_K^\pm}\ar[r]&\mathcal{P}_{\rm Gr}(M_\alpha^\pm;S')^{\Gamma_K^\pm},
}
\]
%where the exactness of the top row follows from the finiteness of $\rH^1_{\Fcal_{{\rm rel},{\rm str}}}(K,W_\alpha)$ in Lemma\,\ref{lem:coinv} as in \cite[Prop.\,3.3.3]{JSW}, 
we obtain
\begin{equation}\label{eq:snake}
\#\rH^1_{\Fcal_{\rm Gr}^{S'}}(K,M_\alpha^\pm)^{\Gamma_K^\pm}=\#\rH^1_{\Fcal_{\Gr}^{S'}}(K,W_\alpha)\cdot\frac{\#{\rm ker}(r^\pm)}{\#{\rm ker}(h^\pm)}.
\end{equation}

Clearly, 
\[
{\rm ker}(h^\pm)=\rH^1(\Gamma_K^\pm,(M_\alpha^\pm)^{G_K})=(M_\alpha^\pm)^{G_K}/(\gamma^\pm-1)(M_\alpha^\pm)^{G_K}=\rH^0(K_\infty^\pm,W_\alpha),
\]
and this vanishes by our assumptions. On the other hand, since we assume that every finite prime $w\in\Sigma$ splits in $K$, the argument in the proof of \cite[Prop.\,3.3.7]{jsw} (but noting that here the roles $v$ and $\bar{v}$ are reversed) shows that for both choices of sign $\pm$, the order of ${\rm ker}(r^\pm)$ is given 
\[
\#{\rm ker}(r^\pm)=\#\rH^0(K_v,W_{\alpha^{-1}})^2\cdot\prod_{w\in\Sigma\smallsetminus S',w\nmid p}c_w^{(p)}(W_\alpha),
\]
where $c_w^{(p)}=\#\rH^1_{\rm ur}(K_w,W_\alpha)$ is the $p$-part of the local Tamagawa number of $W_\alpha$ at $w$. Thus the value in (\ref{eq:snake}) is the same for both choices of sign $\pm$, and together with (\ref{eq:Euler-char}) this yields the result.
%
%Recall that we assume that all finite primes in $\Sigma$ split in $K$. Thus it follows from Proposition~\ref{prop:ps-null}, the anticyclotomic control theorem in \cite[Thm.~3.3.1]{JSW} and (an immediate adaptation of) the cyclotomic control theorem in \cite[Thm.~4.1]{greenberg-cetraro}, that both sides of the claimed equality are given by
%\[
%\#\rH^1_{\Fcal_{\rm rel,\rm str}}(K,W_\alpha)\cdot\#\rH^0(K_v,W_\alpha)\cdot\#\rH^0(K_{\bar{v}},W_\alpha)\cdot
%%\prod_{w\in\Sigma, w\nmid p}\rH^1_{\rm ur}(K_w,W_\alpha)
%\prod_{w\in S}\#\rH^1(K_w,W_\alpha)
%\]
%up to a $p$-adic unit.
\end{proof}



\subsubsection{Ancillary results for $\X_{\rm ord}(E(\alpha)/K_\infty^+)$}
Finally, in this section we prove a congruence relation modulo $\varpi^m$ for the (characteristic power series of the) Selmer groups $\X_{\rm ord}^S(E(\alpha)/K_\infty^+)$ and $\X_{\rm ord}^S(E/K_\infty^+)$, where $\alpha$ is an anticyclotomic character as above such that $\alpha\equiv 1\mod\varpi^m$. We start with the following preliminary lemma. 

%For our arguments later in the paper we shall also need the congruence in Proposition~\ref{prop:cong-Sel} below, whose proof builds on the following lemma.

\begin{lemma}\label{lem:almost-div-ord}
Assume that $E(K)[p]=0$ and that $\X_{\rm ord}^{S'}(E(\alpha)/K_\infty^+)$ is $\Lambda_{K,R}^+$-torsion, where $S'\subset\Sigma$ is a subset consisting of primes away from $p$. Then $\X_{\rm ord}^{S'}(E(\alpha)/K_\infty^+)$ has no nonzero finite $\Lambda_{K,R}^+$-submodules.
\end{lemma}

\begin{proof} 
This is shown in \cite[Prop.\,4.14]{greenberg-cetraro} when $\alpha=1$ and $S'=\emptyset$, and the general case follows from a slight variation of the same arguments (see e.g. \cite[Prop.\,2.3.3]{skinner-mult}). Alternatively, this can be seen as a special case of Greenberg's results \cite{greenberg-Sel}.
\end{proof}

Put
\[
S=\Sigma\smallsetminus\{p,\infty\},
\]
which as above we shall view as a set of primes of $\Q$ or of $K$ according to context. The next important result is an algebraic counterpart of Lemma~\ref{lem:cong-L}.

\begin{prop}\label{prop:cong-Sel}
Assume that $E(K)[p]=0$, that $\X_{\rm ord}^S(E(\alpha)/K_\infty^+)$ is $\Lambda^+_{K,R}$-torsion, and that
$\X_{\rm ord}^S(E/K_\infty^+)$ is $\Lambda_K^+$-torsion. 
%and that $\X_{\rm ord}^S(E(\alpha)/K_\infty^+)$ and $\X_{\rm ord}^S(E/K_\infty^+)$ are both $\Lambda^+$-torsion. 
If $\alpha\equiv 1\;({\rm mod}\,\varpi^m)$, then there are suitable characteristic power series $\Fcal_{\rm ord}^S(E(\alpha)/K_\infty^+)$ and $\Fcal_{\rm ord}^S(E/K_\infty^+)$ for the 
%$\Lambda_K^+$-modules 
modules $\X_{\rm ord}^S(E(\alpha)/K_\infty^+)$ and $\X_{\rm ord}^S(E/K_\infty^+)$, respectively, such that 
\[
\Fcal_{\rm ord}^S(E(\alpha)/K_\infty^+)\equiv\Fcal_{\rm ord}^S(E/K_\infty^+)\;\,({\rm mod}\,\varpi^m).
\]
%\ch_{\Lambda^+}(\X_{\rm ord}^S(E(\alpha)/K_\infty^+))\equiv\ch_{\Lambda^+}(\X^S_{\rm ord}(E/K_\infty^+))\pmod{\varpi^m}.
%\]
%where $\Fcal_{\rm ord}^S(E(\alpha)/K_\infty^+)$ and $\Fcal_{\rm ord}^S(E/K_\infty^+)$ are the
%characteristic power series for the 
%%$\Lambda_K^+$-modules 
%modules $\X_{\rm ord}^S(E(\alpha)/K_\infty^+)$ and $\X_{\rm ord}^S(E/K_\infty^+)$, respectively.  
\end{prop}

\begin{proof}
Since %under the hypotheses 
 $\X_{\rm ord}^S(E(\alpha)/K_\infty^+)$ and $\X_{\rm ord}^S(E/K_\infty^+)$ have no nonzero finite $\Lambda_K^+$-submodules by Lemma~\ref{lem:almost-div-ord}, 
%(taking $\alpha=1$ for the latter), 
their characteristic ideals are the same as their Fitting ideals (see \cite[Lem.~2.3.4(ii)]{skinner-mult}), so to prove the result it suffices to show that
\begin{equation}\label{eq:Sel-m}
\rH^1_{\Fcal_{\rm ord}^S}(K,M)[\varpi^m]\simeq\rH^1_{\Fcal_{\rm ord}^S}(K,M_\alpha)[\varpi^m],
\end{equation}
where $M=T_pE\widehat{\otimes}_{\Zp}(\Lambda_{K,R}^+)^\vee$ and $M_\alpha=T_pE(\alpha)\widehat\otimes_{\Z_p}(\Lambda_K^+)^\vee$. By the assumption  $E(K)[p]=0$, the natural maps  
\[
\rH^1(G_{K,\Sigma},M[\varpi^m])\rightarrow\rH^1(G_{K,\Sigma},M)[\varpi^m],\quad 
\rH^1(G_{K,\Sigma},M_\alpha[\varpi^m])\rightarrow\rH^1(G_{K,\Sigma},M_\alpha)[\varpi^m]
\]
are isomorphisms. Moreover, for any $w\vert p$ the restriction map $\rH^1(K_w,M^-)\rightarrow\rH^1(I_w,M^-)$ is an injection (this uses $p\nmid N$), where we put $M^-=(T_pE/{\rm Fil}_w^+T_pE)\otimes(\Lambda_K^+)^\vee$. Thus $\rH^1_{\Fcal_{\rm ord}^S}(K,M)[\varpi^m]$ is naturally identified with the kernel of the restriction map
\[
\rH^1(G_{K,\Sigma},M[\varpi^m])\rightarrow\prod_{w\vert p}\rH^1(I_w,M^-),
\]
which factors through $\rH^1(G_{K,\Sigma},M[\varpi^m])\rightarrow\prod_{w\vert p}\rH^1(I_w,M^-[\varpi^m])$.  Since the kernel of the natural map $\rH^1(I_w,M^-[\varpi^m])\rightarrow\rH^1(I_w,M^-)$ is given by $(M^-)^{I_w}/p^m(M^-)^{I_w}$, and this is zero since $(M^-)^{I_w}\simeq{\rm Hom}_{\rm cts}(R,\Q_p/\Z_p)$ is $p$-divisible, we conclude that 
\[
\rH^1_{\Fcal_{\rm ord}^S}(K,M)[\varpi^m]={\rm ker}\biggl(\rH^1(G_{K,\Sigma},M[\varpi^m])\rightarrow\prod_{w\vert p}\rH^1(I_w,M^-[\varpi^m])\biggr).
\]
Letting $M_\alpha^-=(T_pE(\alpha)/{\rm Fil}_w^+T_pE(\alpha))\otimes(\Lambda_K^+)^\vee$ we similarly find $(M_\alpha^-)^{I_w}\simeq{\rm Hom}_{\rm cts}(R,\Q_p/\Z_p)$, noting that after restriction to $G_{K_{\infty,w}^+}$ the character $\alpha$ becomes unramified; thus  $(M_\alpha^-)^{I_w}/p^m(M_\alpha^-)^{I_w}=0$ and so $\rH^1_{\Fcal_{\rm ord}^S}(K,M_\alpha)[\varpi^m]$ is identified with the kernel of the restriction map
\[
\rH^1(G_{K,\Sigma},M_\alpha[\varpi^m])\rightarrow\prod_{w\vert p}\rH^1(I_w,M_\alpha^-[\varpi^m]).
\]
Since $M_\alpha[\varpi^m]=M[\varpi^m]$, this proves (\ref{eq:Sel-m}) and yields the result.
\end{proof}



\section{Beilinson--Flach classes}\label{BF}


Let $f=\sum_{n=1}^\infty a_nq^n\in S_2(\Gamma_0(N))$ be a newform with 
%rational
Fourier coefficients in $\Q$, and fix a prime $p\nmid 2N$. We denote by ${Y_1(N)}$ the modular curve of level $\Gamma_1(N)$. The $p$-adic Galois representation $V_f$ associated to $f$ can be geometrically realized as the maximal quotient of $\rH^1_{\et}(\overline{Y_1(N)},\Q_p)(1)$ on which the Hecke operators $T_n$ acts as $a_n$. (Note that this is denoted $V_{\Q_p}(f)^*$ in \cite[Def.~6.3]{RSCMF}, and corresponds to  $V_{\Q_p}(f)(1)$ in the notations of \cite{kato-euler-systems}.) Let $T_f$ be the $\Z_p$-submodule $V_f$ generated by the image of $\rH^1_{\et}(\overline{Y_1(N)},\Z_p)(1)$, which 
%gives 
is a $G_\Q$-stable $\Z_p$-lattice 
%$T_f\subset 
in $V_f$. 

We assume that $f$ is ordinary at $p$, i.e., $a_p\in\Z_p^\times$, so there is a $G_{p}$-stable filtration 
\[
0\rightarrow T_f^+\rightarrow T_f\rightarrow T_f^-\rightarrow 0
\] 
with $T_f^{\pm}$ free rank one $\Z_p$-modules with the $G_{p}$-action on $T_f^-$ given by the unramified character sending an arithmetic Frobenius to the $p$-adic unit root of $x^2-a_px+p$. Replacing ${\rm Fil}_w^+(T_pE)$ with $T_f^+$ we can define the ordinary local condition for $T_f$ as in Definition \ref{defilocalconds}.

Let $K$ be an imaginary quadratic field satisfying (\ref{eq:spl}). 
In this section we introduce the special case of the Beilinson--Flach classes of \cite{explicit} that will be needed for our arguments, and deduce some applications.


\subsection{Reciprocity laws}\label{subsec:ERL}

Let $\boldsymbol{g}\in\Lambda_{\boldsymbol{g}}\llbracket{q}\rrbracket$ be the CM Hida family from $\S\ref{subsec:L-2var-II}$, and fix a character $\alpha:\Gamma_K\rightarrow R^\times$ with values in the ring of integers of a finite extension $\Phi/\Q_p$. With a slight abuse of notation, we continue to denote by $\Lambda_K$ the Iwasawa algebra $R\llbracket{\Gamma_K}\rrbracket$ and by $\Lambda_{\bg}$ the extension of scalars $\Lambda_{\bg}\otimes_{\Z_p} R$. 

%Note that we have a decomposition
%\[
%\Gamma_K\simeq\Gamma_v\times\Gamma_K^+
%\]
%given by the direct sum of the two natural projections. Using this, we may uniquely write $\alpha=\alpha_v\cdot\alpha^+$, with $\alpha_v$ (resp. $\alpha^+$) a character of $\Gamma_v$ (resp. $\Gamma_K^+$).

Let $I_f \subset \Qp$ be the image of the unique Hecke-equivariant map $M_2(\Gamma_0(N),\Z_p)\rightarrow \Q_p$, with the Hecke action on $\Q_p$ 
such that $T_n$ acts as multiplication by $a_n$, that sends $f$ to $1$. Then $I_f$ is a free $\Z_p$-module of rank one. 
We similarly define $I_{\bg}^{\rm cusp} \subset \Frac(\Lambda_{\bg})$ as  the image of the $\Lambda_{\boldsymbol{g}}$-adic cuspforms in ${\rm Frac}(\Lambda_{\boldsymbol{g}})$ under the projection corresponding to ${\boldsymbol{g}}$. Note that $I_{\bg}^{\rm cusp}$ is a finitely-generated $\Lambda_{\bg}$-module. 


\begin{thm}%[Kings--Loeffler--Zerbes]
\label{thm:KLZ}
There exists a class
\[
{BF}_\alpha\in\rH^1_{\Fcal_{\rm ord,\rm rel}}(K,T_f(\alpha)\widehat\otimes\Lambda_K)
\]
and injective $\Lambda_K$-linear maps with pseudo-null cokernel
\begin{align*}
\widetilde{\rm Col}_{f}:\rH^1(K_{\bar{v}},T_f^-(\alpha)\widehat\otimes\Lambda_K) \to I_f\widehat\otimes\Lambda_K, \ \ 
\widetilde{\rm Col}_{\boldsymbol{g}}:\rH^1(K_{v},T_f^+(\alpha)\widehat\otimes\Lambda_K) \to I_{\bg}^{\rm cusp}\widehat\otimes\Lambda_K,
\end{align*}
satisfying the following explicit reciprocity laws:
\begin{itemize}
\item[(a)] 
$\widetilde{\rm Col}_f(p^-({\rm loc}_{\bar v}({BF}_\alpha)))=\Lcal_p(f(\alpha)/K,\Sigma^{(1)})$, where $p^-({\rm loc}_{\bar{v}}({BF}_\alpha))$ is the natural image of ${\rm loc}_{\bar{v}}({BF}_\alpha)$ in $\rH^1(K_{\bar{v}},T_f^-(\alpha)\widehat\otimes\Lambda_K)$;
\item[(b)] $\widetilde{\rm Col}_{\boldsymbol{g}}({\rm loc}_{v}({BF}_\alpha))=\Lcal_p(f(\alpha)/K,\Sigma^{(2')})$.
\end{itemize}
\end{thm}

\begin{proof} This is proved in \cite[\S5]{BST}, where it is deduced from results in \cite{explicit} 
in combination with results in \cite{Betina-Dimitrov-Pozzi}, though here we have reversed the roles of $v$ and $\bar v$. 
\end{proof}

%\begin{proof}
%Taking $m=1$ in \cite[Def.~8.1.1]{explicit} (and using \cite[Lem.~6.8.9]{RSCMF} to dispense with the auxiliary $c>1$ needed for the construction) we obtain a cohomology class
%\begin{equation}\label{eq:BF-v}
%BF^{f\bg_{\alpha_v}}\in\rH^1(\Q,T_f\widehat\otimes V(\bg_{\alpha_v})\widehat\otimes\Lambda_\Q).\nonumber
%\end{equation}
%where $V(\bg_{\alpha_v})$ is the big Galois representation  (as in \cite[Def.~7.2.5]{explicit}) associated to the CM Hida family
%\[
%\bg_{\alpha_v}=\sum_{(\mathfrak{a},\bar{v})=1}\alpha_v(\mathfrak{a})[\mathfrak{a}]q^{\mathbf{N}(\mathfrak{a})}\in\Lambda_{\bg}\llbracket
%{q}\rrbracket.
%\]
%(thus $\bg_{\alpha_v}=\bg$ when $\alpha=1$). As usual, here we let $G_\Q$ act on $\Lambda_\Q$ via $\Psi:G_\Q\rightarrow\Lambda_\Q^\times$. 
%
%As shown in \cite{BST} building on the results of \cite{Betina-Dimitrov-Pozzi} (especially for $\alpha=1$), there is a $G_\Q$-module isomorphism
%\[
%V({\bg_{\alpha_v}})\simeq{\rm Ind}_K^\Q\alpha_v\Psi_v^{-\tau},
%\]
%where $\Psi_v:G_K\twoheadrightarrow\Gamma_v\hookrightarrow\Lambda_{\bg}^\times$ is the universal character. Therefore twisting by $\alpha^+$ the class $BF^{f\bg_{\alpha_v}}$ and applying Shapiro's lemma, we obtain the class
%\[
%BF_\alpha\in\rH^1(K,T_f(\alpha)\widehat\otimes\Lambda_K).
%\]
%
%Now by \cite[Prop.~8.1.7]{explicit} the restriction ${\rm loc}_p(BF^{f\bg_{\alpha_v}})$ maps to zero under the natural map
%\begin{equation}\label{eq:--}
%\rH^1(\Q_p,T_f\widehat\otimes V(\bg_{\alpha_v})\widehat\otimes\Lambda_\Q)\to\rH^1(\Q_p,T_f^-\widehat\otimes V(\bg_{\alpha_v})^-\widehat\otimes\Lambda_\Q);
%\end{equation}
%equivalently, if we put 
%\begin{equation}\label{eq:bal}
%\mathscr{F}_p^{\rm bal}(T_f\widehat\otimes V(\bg_{\alpha_v})):=T_f^+\widehat\otimes V(\bg_{\alpha_v})+T_f\widehat\otimes V(\bg_{\alpha_v})^+,
%\end{equation}
%the class ${\rm loc}_p(BF^{f\bg_{\alpha_v}})$ belongs to the image of the natural map
%\[
%\rH^1(\Q_p,\mathscr{F}_p^{\rm bal}(T_f\widehat\otimes V(\bg_{\alpha_v}))\widehat\otimes\Lambda_\Q)\to\rH^1(\Q_p,T_f\widehat\otimes V(\bg_{\alpha_v})\widehat\otimes\Lambda_\Q),
%\]
%which is easily seen to be an injection. From Shapiro's lemma we find
%\[
%\mathscr{F}_w^{\rm bal}(T_f\widehat\otimes V(\bg_{\alpha_v}))\widehat\otimes\Lambda_\Q\simeq\begin{cases}
%T_f^+(\alpha_v)\widehat\otimes\Lambda_K&\textrm{if $w=v$,}\\[0.2em]
%T_f(\alpha_v)\widehat\otimes\Lambda_K&\textrm{if $w=\bar{v}$}.
%\end{cases}
%\]
%Thus the fact that ${\rm res}_p(BF^{f\bg_{\alpha_v}})$ lands in the kernel of $(\ref{eq:--})$ implies the inclusion 
%\[
%{\rm loc}_{v}(BF_\alpha)\in\rH^1_{\Fcal_{\rm ord}}(K_{v},T_f(\alpha)\widehat\otimes\Lambda_K),
%\] 
%proving the first claim in the theorem. 
%
%Next, put $\mathscr{F}_p^{-+}(T_f\widehat\otimes V(\bg_{\alpha_v}))=T_f^-\otimes V(\bg_{\alpha_v})^+$, and define $\widetilde{\rm Col}_f$ to be the twist by $\alpha^+$ of the composition
%\begin{align*}
%\rH^1(K_{\bar v},T_f^-(\alpha_v)\widehat\otimes\Lambda_K)&\simeq\rH^1(\Q_p,\mathscr{F}_p^{-+}(T_f\widehat\otimes V(\bg_{\alpha_v}))\widehat\otimes\Lambda_\Q)\\
%&\xrightarrow{\Lcal}\mathbb{D}(\mathscr{F}_p^{-+}(T_f\widehat\otimes V(\bg_{\alpha_v})))\widehat\otimes\Lambda_\Q\\
%&\xrightarrow{\eta_f\otimes\omega_{\bg_{\alpha_v}}}I_f\widehat\otimes\Lambda_{\bg}\widehat\otimes\Lambda_\Q\simeq I_f\otimes\Lambda_K,
%\end{align*}
%where $\mathcal{L}$ is the Perrin-Riou big logarithm map of \cite[Thm.~8.2.8]{explicit}, and $\eta_f, \omega_{\bg_{\alpha_v}}$ are as constructed in \cite[Prop.~10.1.1]{explicit}. Similarly, define $\widetilde{\rm Col}_{\bg_\alpha}$ to be the twist by $\alpha^+$ of the composition
%\begin{align*}
%\rH^1(K_{{v}},T_f^+(\alpha_v)\widehat\otimes\Lambda_K)&\simeq\rH^1(\Q_p,\mathscr{F}_p^{+-}(T_f\widehat\otimes V(\bg_{\alpha_v}))\widehat\otimes\Lambda_\Q)\\
%&\xrightarrow{\Lcal}\mathbb{D}(\mathscr{F}_p^{+-}(T_f\widehat\otimes V(\bg_{\alpha_v})))\widehat\otimes\Lambda_\Q\\
%&\xrightarrow{\omega_f\otimes\eta_{\bg_{\alpha_v}}}I_{\bg}\widehat\otimes\Lambda_{\Q}\simeq I_{\bg}\otimes\Lambda_K.
%\end{align*}
%The second claim in the theorem now follows from the reciprocity law of \cite[Thm.~10.2.2]{explicit}.
%\end{proof}

%Similarly as in Definition~\ref{def:PR}, put
%\[
%BF_{E(\alpha)}:=\frac{\deg(\pi_E)}{c_E^2}\cdot BF_{\alpha},
%\]
%where $\pi_E:X_0(N)\rightarrow E$ is a modular parametrization of minimal degree, and $c_E$ is the associated Manin constant.


Let $E_1$ be the optimal curve in the isogeny class associated with $f$, in the sense of \cite{stevens-invmath}, with optimal parametrization $\pi_1:X_1(N)\rightarrow E_1$. 
The inclusion $Y_1(N)\hookrightarrow X_1(N)$ identifies the maximal quotient of $\rH^1_{\et}(\overline{X_1(N)},\Q_p(1))$ on which the Hecke operators $T_n$ acts as $a_n$
with the similar quotient of $\rH^1_{\et}(\overline{Y_1(N)},\Q_p(1))$, that is, with $V_f$. Via this identification and the parametrization $\pi_1$, the Tate module $T_pE_1$ is identified with the image of $\rH^1_{\et}(\overline{X_1(N)},\Z_p(1))$ in $V_f$; this is a sublattice of $T_f$. 
Let $E_\bullet/\Q$ be the elliptic curve in the isogeny class associated with $f$ constructed in \cite{wuthrich-int}. By construction, there is a cyclic isogeny
$\phi:E_1\rightarrow E_\bullet$ that is \'etale in characteristic $p$ and for which the resulting parametrization $\pi_\bullet = \phi\circ\pi_1:X_1(N)\rightarrow E_\bullet$ 
identifies $T_pE_\bullet$ with $T_f$ in $V_f$:
\[
T_pE_\bullet = T_f.
\] 
(cf. Proposition~8 and Theorem~4 of  \emph{op.\,cit.}).









%Let $E_0$ be the optimial curve in the isogeny class associated with $f$, in the sense of with minimal parametrization $\pi_0:X_0(N)\rightarrow E_0$. By construction (see \cite[Thm.~4]{wuthrich-int}), the curve $E_\bullet$ is an \'{e}tale quotient of $E_0$. Let $\pi_\bullet:X_0(N)\rightarrow E_\bullet$ be the composition of $\pi_0$ with an \'{e}tale isogeny
%$\phi:E_0\rightarrow E_\bullet$ of minimal degree: $\pi_\bullet = \phi\circ\pi_0$.

\begin{lemma}\label{lem:cong-deg}
We have $\deg(\pi_\bullet)I_f  = \Zp$.
\end{lemma}

\begin{proof} 
Let $\omega_{E_\bullet}$ (resp.~$\omega_{E_1}$) be a N\'eron differential for $E_\bullet$ (resp.~$E_1$). Then 
\[
\pi_{\bullet,*}\pi_\bullet^*\omega_{E_\bullet} = \deg(\pi_\bullet)\omega_{E_\bullet}.
\]
As the isogeny $\phi$ is \'etale, $\phi^*\omega_{E_\bullet} = u_1 \omega_{E_1}$ for some $u_1\in \Z_p^\times$. Similarly,
$\pi_1^* \omega_{E_1} = c_1\omega_f$ for some $c_1\in \Z_p^\times$, as $p\nmid 2N$ (see \cite[Cor.\,4.1]{mazur} and \cite[Prop.\,~3.3]{greenvats}). Hence, $\pi_{\bullet,*}\omega_f = u_1c_1\deg(\pi_\bullet)\omega_{E_\bullet}$
and so it suffices to show that $\pi_{\bullet,*}\omega_f = a \omega_{E_\bullet}$ for some $a\in\Zp$ such that $a I_f = \Zp$.

As $E_\bullet$ has ordinary reduction at the prime $p$, the induced map $\rH^1_{\et}(\overline{Y_1(N)},\Zp(1))\twoheadrightarrow T_f = T_pE_\bullet$ factors through 
projection to the ordinary summand $\mathcal{H}^{\ord} :=e_{\ord} \rH^1_{\et}(\overline{Y_1(N)},\Zp(1))$. There is a corresponding commutative diagram
\[
\begin{tikzcd}
\mathcal{H}^+ \arrow[r,hook] \arrow[d,two heads] & \mathcal{H}^{\ord} \arrow[r,two heads] \arrow[d,two heads] & \mathcal{H}^- \arrow[d,two heads] \\
T_{f}^+ \arrow[r,hook] & T_{f} \arrow[r,two heads] & T_f^-
\end{tikzcd}
\]
of $G_{\Qp}$-modules,
where $\mathcal{H}^+$ (resp.~$\mathcal{H}^-$) is the maximal submodule (resp.~quotient) on which the inertia group $I_p$ acts non-trivially (resp.~trivially); this non-trivial
action is via the cyclotomic character. The middle arrow is the defining projection, which induces the other two maps. 

Let $M(f)$ denote the maximal torsion-free quotient of $M_2(\Gamma_1(N),\Z_p)$ on which the Hecke operator $T_n$ acts as $a_n$ for all $n$. 
Since $p\nmid 2N$, the integral de Rham - \'{e}tale comparison isomorphisms (we need the log version for the open curve $Y_1(N)$)
induce compatible identifications 
\[
e_{\ord}\rH^0(\Omega_{X_1(N)/\Zp}(\log(\mathrm{cusps})) = e_{\ord} M_2(\Gamma_1(N))\simeq \mathcal{H}^-
\]
and
\[ 
M(f) \simeq T_f^-= T_pE_\bullet^- \simeq\rH^0(\Omega_{E_\bullet/\Zp}) = \Zp \omega_{E_\bullet}. 
\]
Via these, $\pi_{\bullet,*}\omega_f$ is identified with the image of $f$ in
$M(f)$. Since $\omega_{E_\bullet}$ is identified with a $\Zp$-generator of $M(f)$ and we have an isomorphism of rank one $\Zp$-modules $M(f) \simeq I_f$, $f\mapsto 1$,
it follows that $\pi_{\bullet,*}\omega_f = a \omega_{E_\bullet}$ for some $a\in \Zp$ such that $aI_f = \Zp$, as desired.
\end{proof} 



 

%\begin{defi}\label{def:BF-alpha}
%Fix an isogeny $\phi:E_0\rightarrow E_\bullet$, and let 
%\[
%{BF}_{0,\alpha}\in\rH^1(K,(T_pE_0)(\alpha)\widehat\otimes\Lambda_K)
%\]
%be the image of $BF_\alpha$ under the induced pullback map 
%\[
%\phi^*:\rH^1(K,(T_pE_{\bullet})(\alpha)\widehat\otimes\Lambda_K)\rightarrow\rH^1(K,(T_pE_0)(\alpha)\widehat\otimes\Lambda_K).
%\]
%\end{defi}
%

Recall the $\alpha$-twisted versions of the two-variable $p$-adic $L$-functions $\Lcal_p^{\rm PR}(E_\bullet/K)$ and $\Lcal_p^{\rm Gr}(f/K)$ introduced in Definitions~\ref{def:PR} and \ref{def:Gr}, respectively.

\begin{cor}\label{cor:KLZ}
%Let $BF_\alpha$ be as in Theorem~\ref{thm:KLZ}. 
There are injective $\Lambda_K$-linear maps with pseudo-null cokernel 
\[
{\rm Col}_{E_\bullet}:\rH^1(K_{\bar{v}},(T_pE_\bullet)^-(\alpha)\widehat\otimes\Lambda_K)\to\Lambda_K,\quad
{\rm Col}_{\boldsymbol{g}}:\rH^1(K_{v},(T_pE_\bullet)^+(\alpha)\widehat\otimes\Lambda_K)\to\Lambda_K^{\rm ur},
\]
such that 
\[
{\rm Col}_{E_\bullet}(p^-({\rm loc}_{\bar v}(BF_{\alpha})))=\Lcal_p^{\rm PR}(E_\bullet(\alpha)/K),\quad{\rm Col}_{\boldsymbol{g}}({\rm loc}_{v}(BF_{\alpha}))=\Lcal_p^{\rm Gr}(f(\alpha)/K).
\]
\end{cor}

\begin{proof}
By \cite[Cor.~4.1]{mazur} and \cite[Prop.~3.3]{greenvats}, the Manin constant associated to the modular parametrization $\pi_\bullet:X_1(N)\rightarrow E_\bullet$ %as defined in the proof of Lemma~\ref{lem:cong-L}
is a $p$-adic unit, so setting 
\[
{\rm Col}_{E_\bullet}:=\deg(\pi_{\bullet})\cdot H_p(f)\cdot\widetilde{\rm Col}_{f},
\]
the first part of the result follows from Theorem~\ref{thm:KLZ} and Lemma~\ref{lem:cong-deg}. On the other hand, as explained in the proof of Lemma~\ref{lem:integral}, the characteristic power series $H_{\bg}^{\rm cusp}$ of the CM family $\boldsymbol{g}$ divides $h_K\cdot\Lcal_v(K)^-$. Since on the other hand by \cite[Cor.~5.6]{hida-coates} $H(\boldsymbol{g})$ is divisible by $h_K\cdot\Lcal_v(K)^-$, it follows that the congruence ideal of $\boldsymbol{g}$ is generated by $h_K\cdot\Lcal_v(K)^-$. Thus setting 
\[
\quad{\rm Col}_{\boldsymbol{g}}:=h_K\cdot\Lcal_v(K)^-\cdot\widetilde{\rm Col}_{\boldsymbol{g}}
\] 
the second part of the result follows from %Definition~\ref{def:Gr} and 
Theorem~\ref{thm:KLZ}.
\end{proof}

\subsection{Iwasawa main conjectures}

One key application of Corollary~\ref{cor:KLZ} is in relating different instances of the main conjectures for the cyclotomic $\Z_p$-extension $K_\infty^+/K$. Similarly as in $\S\ref{subsec:Sel-str}$, for any elliptic curve $E/\Q$ we put 
\begin{align*}
\X_{\rm Gr}(E(\alpha)/K_\infty^+):=\X_{\rm rel,\rm str}(E(\alpha)/K_\infty^+),\quad
\mathfrak{S}_{\rm ord,\rm rel}(E(\alpha)/K_\infty^+)&=\rH^1_{\Fcal_{\rm ord,\rm rel}}(K,T_pE(\alpha)\otimes\Lambda_K^+),
\end{align*}
etc.. Write $BF_{\alpha}^+$ for the image of the class $BF_{\alpha}$ of Theorem~\ref{thm:KLZ} under the natural projection
\[
\rH^1_{\Fcal_{\rm ord,\rm rel}}(K,(T_pE_\bullet)(\alpha)\widehat\otimes\Lambda_K)\rightarrow
\rH^1_{\Fcal_{\rm ord,\rm rel}}(K,(T_pE_\bullet)(\alpha)\widehat\otimes\Lambda_K^+), 
\]
so we have $BF_{\alpha}^+\in\mathfrak{S}_{\rm ord,\rm rel}(E_\bullet(\alpha)/K_\infty^+)$.


\begin{prop}\label{prop:equiv}
Suppose $E_\bullet(K)[p]=0$ and that $\Lcal_p^{\rm PR}(E_\bullet(\alpha)/K)^+$ and $\Lcal_p^{\rm Gr}(f(\alpha)/K)^+$ are both nonzero. Then the following are equivalent:
\begin{itemize}
\item[(i)] $\mathfrak{S}_{\rm ord,\rm rel}(E_\bullet(\alpha)/K_\infty^+)$ has $\Lambda_K^+$-rank one, $\X_{\rm ord,\rm str}(E_\bullet(\alpha^{-1})/K_\infty^+)$ is $\Lambda_K^+$-torsion, and
\[
\ch_{\Lambda_K^+}\bigl(\X_{\rm ord,\rm str}(E_\bullet(\alpha^{-1})/K_\infty^+\bigr)\supset\ch_{\Lambda_K^+}\bigl(\mathfrak{S}_{\rm ord,\rm rel}(E_\bullet(\alpha)/K_\infty^+)/\Lambda_K^+BF_{\alpha}^+\bigr).
\]
\item[(ii)] $\mathfrak{S}_{\rm str,\rm rel}(E_\bullet(\alpha)/K_\infty^+)$ and $\X_{\rm Gr}(E_\bullet(\alpha)/K_\infty^+)$ are both $\Lambda_K^+$-torsion, and
\[
\ch_{\Lambda_K^+}\bigl(\X_{\rm Gr}(E_\bullet(\alpha)/K_\infty^+)\bigr)\Lambda_K^{+,\rm ur}\supset\bigl(\Lcal_p^{\rm Gr}(f(\alpha)/K)^+\bigr).
\]
\item[(iii)] $\mathfrak{S}_{\rm ord}(E_\bullet(\alpha)/K_\infty^+)$ and $\X_{\rm ord}(E_\bullet(\alpha)/K_\infty^+)$ are both $\Lambda_K^+$-torsion, and
\[
\ch_{\Lambda^+_K}\bigl(\X_{\rm ord}(E_\bullet(\alpha)/K_\infty^+)\bigr)\supset\bigl(\Lcal_p^{\rm PR}(E_\bullet(\alpha)/K)^+\bigr).
\]
\end{itemize}
The same result holds for the opposite divisibilities.
\end{prop}

\begin{proof}
This is a well-known consequence of Poitou--Tate duality and the reciprocity laws of Theorem~\ref{thm:KLZ}, but we provide the details for the convenience of the reader. Below we set $\mathfrak{S}_{\rm str,\rm rel}=\mathfrak{S}_{\rm str,\rm rel}(E_\bullet(\alpha)/K_\infty^+)$, $\X_{\rm Gr}=\X_{\Gr}(E_\bullet(\alpha)/K_\infty^+)$, etc. for the ease of notation; and similarly, $\X_{\rm ord}(\alpha^{-1})=\X_{\rm ord}(E_\bullet(\alpha^{-1})/K_\infty^+)$, $\X_{\rm Gr}(\alpha^{-1})=\X_{\Gr}(E_\bullet(\alpha^{-1})/K_\infty^+)$, etc.. Note that, since $\alpha^{\tau}=\alpha^{-1}$ and $(T_pE_{0}\otimes (\Lambda_K^+)^{\vee})^{\tau}\simeq T_pE_{0}\otimes (\Lambda_K^+)^{\vee}$, the action of complex conjugation gives rise to isomorphisms of $\Lambda_K^+$-modules:
\begin{equation}\label{eq:cc}
\X_{\Gr}(\alpha^{-1})\simeq \X_{\Gr}, \ \ \ \X_{\rm ord}(\alpha^{-1})\simeq \X_{\rm ord} \ \ \ \X_{\rm ord, str}(\alpha^{-1})\simeq \X_{\rm str, ord}.
\end{equation}

For the equivalence ${\rm (i)}\Leftrightarrow{\rm (ii)}$ consider the exact sequence
\begin{equation}\label{eq:PT1}
0\rightarrow\mathfrak{S}_{\rm str,\rm rel}\rightarrow\mathfrak{S}_{\rm ord,\rm rel}\rightarrow\rH^1_{\rm ord}(K_{v},T_pE_\bullet(\alpha)\otimes\Lambda_K^+)\rightarrow\X_{\rm Gr}(\alpha^{-1})\rightarrow\X_{\rm ord,\rm str}(\alpha^{-1})\rightarrow 0.
\end{equation}
In both cases, we see that $\mathfrak{S}_{\rm str,\rm rel}$ is $\Lambda_K^+$-torsion, hence trivial (by the assumption $E_\bullet(K)[p]=0$), and (\ref{eq:PT1}) yields the exact sequence
\[
0 \to\mathfrak{S}_{\rm ord,\rm rel}/\Lambda_K^+ BF^+_{\alpha} \to\rH^1_{\rm ord}(K_{v},T_pE_\bullet(\alpha)\otimes\Lambda_K^+)/\Lambda_K^+\loc_{v}(BF_{\alpha}^+)\to\X_{\rm Gr}(\alpha^{-1})\to \X_{\rm ord,\rm str}(\alpha^{-1}) \to 0.
\]
Since by Corollary~\ref{cor:KLZ} the second term in this exact sequence is pseudo-isomorphic---via the map ${\rm Col}_{\bg}$---to $\Lambda_K^{+,\rm ur}/(\Lcal_p^{\rm Gr}(f(\alpha)/K)^+)$, the equivalence ${\rm (i)}\Leftrightarrow{\rm (ii)}$ follows from this, the multiplicativity of characteristic ideals, and the isomorphisms \eqref{eq:cc}. On the other hand, for the equivalence ${\rm (i)}\Leftrightarrow{\rm (iii)}$ consider the exact sequence
\begin{equation}\label{eq:PT2}
0\rightarrow\mathfrak{S}_{\rm ord}\rightarrow\mathfrak{S}_{\rm ord,\rm rel}\rightarrow\rH_{/\rm ord}^1(K_{\bar v},T_pE_\bullet(\alpha)\otimes\Lambda_K^+)\rightarrow\X_{\rm ord}(\alpha^{-1})\rightarrow\X_{\rm ord,\rm str}(\alpha^{-1})\rightarrow 0,
\end{equation}
where 
\[
\rH^1_{/\rm ord }(K_{\bar v},T_pE_\bullet(\alpha)\otimes\Lambda_K^+):=\frac{\rH^1(K_{\bar v},T_pE_\bullet(\alpha)\otimes\Lambda_K^+)}{\rH^1_{\rm ord }(K_{\bar v},T_pE_\bullet(\alpha)\otimes\Lambda_K^+)}\simeq\rH^1(K_{\bar v},T_f^-(\alpha)\widehat\otimes\Lambda_K).
\]
Similarly as before, in both cases we find $\mathfrak{S}_{\rm ord}$ is $\Lambda_K^+$-torsion, so from (\ref{eq:PT2}) we obtain the exact sequence
\[
0\to \mathfrak{S}_{\rm ord,\rm rel}/\Lambda_K^+ BF_{\alpha}^+ \to \rH^1_{/\rm ord}(K_{\bar{v}},T_pE_\bullet\otimes\Lambda_K^+)/\Lambda_K^+p^-(\loc_{\bar{v}}(BF_{\alpha}^+))\to \X_{\rm ord}(\alpha^{-1})\to \X_{\rm ord,\rm str}(\alpha^{-1}) \to 0.
\]
Since by Corollary~\ref{thm:KLZ} the second term in this exact sequence is pseudo-isomorphic---via the map ${\rm Col}_{E_\bullet}$---to $\Lambda_K^+/(\Lcal_p^{\rm PR}(E_\bullet(\alpha)/K)^+)$, taking characteristic ideals and applying \eqref{eq:cc} the result follows.
\end{proof}

%{\color{blue} 
%\begin{rem}
%As noted in \cite[\S{2.1}]{wuthrich-int}, if $N$ is square-free and $E_\bullet(\Q)[p]=0$, then $E_\bullet=E_0$ (and so $d=1$ in Proposition~\ref{prop:equiv}).
%\end{rem} }
 
\subsection{The Beilinson--Flach Euler system divisibility}\label{sec:BF}

In the case $\alpha=1$, the ``upper bound'' divisibility in Conjecture~\ref{conj:IMC-K} for $\X_{\rm ord}(E_\bullet/K_\infty^+)$ can be deduced from Kato's work together with Proposition~\ref{prop:comp-Selcyc}. 
For our later arguments (especially for elliptic curves $E/\Q$ of rank $>1$), we will need a similar divisibility for a twist by a non-trivial character $\alpha$ of $\Gamma_K^-$, a result that we shall deduce from 
%the equivalence in 
Proposition~\ref{prop:equiv} and the next result.

%Let us denote by $\X_{\rm str}(E_{\bullet}\otimes {\rm Ind}_K^\Q(\alpha)/\Q_{\infty})$ the Pontryagin dual of $\rH^1_{\rm str}(\Q,T_pE_{\bullet}\otimes {\rm Ind}_K^\Q(\alpha)\otimes (\Lambda_\Q)^{\vee})$.

\begin{thm}\label{thm:BF-ES} 
Suppose $\alpha\neq 1$ is a non-trivial character of $\Gamma_K^-$ and 
$BF_{\alpha}^+\neq 0$. Then $\X_{\rm ord,\rm str}(E_\bullet(\alpha^{-1})/K_\infty^+)$ is $\Lambda_K^+$-torsion,  $\mathfrak{S}_{\rm ord,\rm rel}(E_\bullet(\alpha)/K_\infty^+)$ has $\Lambda_K^+$-rank one, and we have the divisibility
\[
\ch_{\Lambda_K^+}\bigl(\X_{\rm ord,\rm str}(E_\bullet(\alpha^{-1})/K_\infty^+\bigr)\supset\ch_{\Lambda_K^+}\bigl(\mathfrak{S}_{\rm ord,\rm rel}(E_\bullet(\alpha)/K_\infty^+)/\Lambda_K^+BF_{\alpha}^+\bigr)
\]
in $\Lambda_K^+\otimes\Q_p$.
\end{thm}

\begin{proof}
As a consequence of the cyclotomic Euler system constructed in \cite[Thm.\,8.1.3]{explicit} attached to the pair $(f,\boldsymbol{g})$, there exists an integer $r\ge 0$ such that $\varpi^rBF_\alpha^+$ 
%of the class $BF_\alpha$ in Theorem\,\ref{thm:KLZ} 
extends to a system of cohomology classes
\[
BF_{\alpha,m}^+\in\rH^1(K(\mu_m),(T_pE_\bullet)(\alpha)\widehat\otimes\Lambda_K^+)
\]
indexed by integers $m\geq 1$ coprime to $pcND_K$ for an auxiliary integer $c>1$ coprime to $6ND_K$ with $BF_{\alpha,1}^+=\varpi^rBF_\alpha^+$ and satisfying the Euler system norm relations. (The $\varpi^r$ appears because the specialization at the character $\alpha$ of the Galois module
associated to $\boldsymbol{g}$ in \cite{explicit} may not equal ${\rm Ind}_K^\Q(\alpha)$ but only contain this lattice with finite index.) 
%In particular, these classes form an Euler system for $(T_pE_0,\mathcal{K})$ with the local condition determined by (\ref{eq:bal}) in the sense of \cite[Def.~12.2.1]{explicit}, where $\mathcal{K}$ is the maximal abelian extension of $\Q$ unramified at the primes dividing $cND_K$.
%
Thus by the results of \cite[\S{12}]{explicit}, giving in particular a refinement of the results of \cite{rubin-ES} for Euler systems with a non-trivial local condition at $p$, it suffices to verify that the $G_\Q$-module $T:=(T_pE_\bullet)\otimes_{\Zp}{\rm Ind}_K^\Q(\alpha)$ satisfies the following hypotheses:
\begin{itemize}
%\item[(i)] $V=T\otimes\Q_p$ is irreducible as a $\Phi[{\rm Gal}(\overline{\Q}/\Q^{\rm ab})]$-module,
\item[(i)] { $V=T\otimes\Q_p$ is irreducible as a $\Phi[{\rm Gal}(\overline{\Q}/\Q(\mu_{p^\infty}))]$-module,}
\item[(ii)] There exists an element $\sigma\in{\rm Gal}(\overline{\Q}/\Q(\mu_{p^{\infty}}))$ such that ${\rm dim}_{\Phi}(V /(\sigma-1) V)=1$.
\end{itemize}
{Here, as before, $\Phi/\Q_p$ is a finite extension contaning the values of the character $\alpha$.}
(Note that \cite[Thm.~12.3.4]{explicit} gives a refinement of \cite[Thm.~II.3.3]{rubin-ES} under the big image hypothesis ``Hyp$(BI)$'' in \emph{op.\,cit.}; the same methods yield a corresponding refinement of \cite[Thm.~II.3.4]{rubin-ES} under the above hypotheses (i)-(ii).)
%But the verification of these hypotheses is standard.
{
It remains to verify hypotheses (i) and (ii).  

We first explain that $E_\bullet$ does not have CM. Suppose to the contrary that $E_\bullet$ has CM by an order
in an imaginary quadratic field $\mathcal{K}$. We note that $p$ must be unramified in $\mathcal{K}$, as $E_\bullet$ is assumed to have good reduction at $p$:
$V_\ell E_\bullet = T_\ell E_\bullet\otimes_{\Z_\ell}{\Q_\ell}$ is the induction from $\Gal(\overline{\Q}/\mathcal{K})$ of the $\ell$-adic Galois-representation associated
with a Hecke character of $\mathcal{K}$ and so is ramified at any prime $\neq\ell$ that is ramified in $\mathcal{K}/\Q$ while also being unramified
at any prime $\neq\ell$ at which $E_\bullet$ has good reduction. 
Next, recall that $E_\bullet[p]^{ss} \simeq \mathbb{F}_p(\phi)\oplus \mathbb{F}_p(\psi)$.  As $E_\bullet$ is assumed to have CM by $\mathcal{K}$, $a_\ell(E_\bullet)=0$ for any prime $\ell$ 
of good reduction for $E_\bullet$ that is inert in $\mathcal{K}/\Q$. It follows that $\psi = \phi\eta_{\mathcal{K}}$, where $\eta_\mathcal{K}$ is the quadratic character of the extension
$\mathcal{K}/\Q$. But then $\omega = \phi\psi = \phi^2\eta_{\mathcal{K}}$ implies that both $\phi$ and $\psi$ are ramified at $p$, contradicting the assumption
that $E_\bullet$ has good ordinary reduction at $p$ and hence one of $\phi$ and $\psi$ is unramified at $p$.

Let $\rho_{E_\bullet}:G_\Q\rightarrow\Aut_{\Zp}(T_pE_\bullet) \simeq \GL_2(\Zp)$ be the usual Galois action on the $p$-adic Tate module of $E_\bullet$.  
Let $K_\infty^{cyc} = K(\mu_{p^\infty})$. As $E_\bullet$ does not have CM, it follows from Serre's open image theorem 
that $\rho_{E_\bullet}(\Gal(\overline{\Q}/K_\infty^{cyc}))$ is an open subgroup of $\SL_2(\Zp)$. In particular, 
$V_pE_\bullet = T_pE_\bullet\otimes_{\Zp}\Qp$ is an irreducible $\Gal(\overline{\Q}/K_\infty^{cyc})$-representation.
It follows that the restriction of $V = V_pE_\bullet \otimes_{\Zp}{\mathrm{Ind}}_K^\Q(\alpha)$ to $\Gal(\overline{\Q}/K_\infty^{cyc})$
decomposes as $V = V_pE_\bullet(\alpha) \oplus V_pE_\bullet(\alpha^{-1})$, with the summands being irreducible.
Let $\ell$ be a prime that ramifies in $K$. Then $\ell\neq p$ as $p$ splits in $K$ by \eqref{eq:intro-spl}, and $E_\bullet$ has good reduction at $\ell$ 
as any prime dividing the conductor of $E_\bullet$ is unramified in $K$ by \eqref{eq:intro-Heeg}.  
Let $\tau_\ell\in I_\ell \subset \Gal(\overline{\Q}/\Q(\mu_{p^\infty}))$, where $I_\ell$ denotes the inertia subgroup at any prime above $\ell$, be an element with non-trivial image in $\Gal(K/\Q)$. 
Then $\tau_\ell$ acts trivially on $V_pE_\bullet$ but maps $V_pE_\bullet(\alpha)$ isomorphically onto $V_pE_\bullet(\alpha^{-1})$ 
(and vice versa). 

Suppose $0\neq W \subset V$ is an irreducible $\Phi[\Gal(\overline{\Q}/\Q(\mu_{p^\infty}))]$-subrepresentation.  
Then $W \cap V_pE_\bullet(\alpha)$ is either $0$ or all of $V_pE_\bullet(\alpha)$. As $W$ is $\tau_\ell$-stable and $V_pE_\bullet(\alpha)$ is not, it must
be that $W\cap V_pE_\bullet(\alpha) = 0$. It follows that $W$ projects isomorphically onto $V_pE_\bullet(\alpha^{-1})$ as 
a $\Gal(\overline{\Q}/K_\infty^{cyc})$-representation. The same argument also shows that $W$ projects isomorphically onto $V_pE_\bullet(\alpha)$ as 
a $\Gal(\overline{\Q}/K_\infty^{cyc})$-representation.  But this means that $V_pE_\bullet(\alpha)\simeq V_pE_\bullet(\alpha^{-1})$ as 
$\Gal(\overline{\Q}/K_\infty^{cyc})$-representations. In turn this implies that $V_pE_\bullet \simeq V_pE_\bullet(\alpha^2)$ as 
$\Gal(\overline{\Q}/K_\infty^{cyc})$-representations. However, the determinant of $V_pE_\bullet$ is trivial on $\Gal(\overline{\Q}/K_\infty^{cyc})$
while $\alpha^4$ is not (as $\alpha$ is non-trivial and pro-$p$), so $V_pE_\bullet \not\simeq V_pE_\bullet(\alpha^2)$. This contradiction proves (i).

To establish (ii), we first choose $\sigma_1\in \Gal(\overline{\Q}/K_\infty^{cyc})$ so that 
$\rho_{E_\bullet}(\sigma_1) \simeq \left(\smallmatrix 1 & * \\ 0 & 1 \endsmallmatrix\right)$ with $*\neq 0$. This is possible
as $\rho_{E_\bullet}(\Gal(\overline{\Q}/K_\infty^{cyc}))$ is an open subgroup of $\SL_2(\Zp)$. The same is then 
true of $\rho_{E_\bullet}(\tau_\ell\sigma_1)$ as $\rho_{E_\bullet}(\tau_\ell) = 1$. 
Let $V_\alpha$ be the $\Phi$-space underlying $\mathrm{Ind}_K^\Q(\alpha)$ and let $\rho_\alpha: G_\Q\rightarrow \Aut_\Phi(V_\alpha) \simeq \GL_2(\Phi)$ denote
the action of $G_\Q$ on $V_\alpha$. Then $\det\rho_\alpha = \eta_K$, the quadratic character associated with the extension $K/\Q$.
As $\tau_\ell\sigma_1 \not\in G_K$, it follows that $\rho_\alpha(\tau_\ell\sigma_1) \simeq \left(\smallmatrix 1 & 0 \\ 0 & -1 \endsmallmatrix\right)$.
Clearly, the element $\sigma = \tau_\ell\sigma_1\in \Gal(\overline{\Q}/\Q(\mu_{p^\infty}))$ satisfies (ii).
}
\end{proof}



%Next we note that the hypotheses that $\alpha$ is crystalline implies that it is the $p$-adic representation associated with some anticyclotomic Hecke character of
%$K$. Furthermore, this character is everywhere unramified as $\alpha$ is a character of $\Gamma_K^-$ and so unramified away from the primes above $p$.
%In particular, there is a CM newform $g$ of odd weight $2n+1$, level prime to $p$, and character $\chi_K$, such that the usual $p$-adic Galois
%representation associated to $g$ satisfies $\rho_g\otimes\varepsilon^{-n} \simeq {\rm Ind}_K^\Q(\alpha)$.  Note also that since $\alpha$ is non-trivial,
%it must have non-zero Hodge--Tate weights and so $n>0$; that is, the weight of $g$ is at least $3$.
%
%It follows that $V = V_f\otimes_{\Qp} V_g(-n)$ and so hypotheses (i) and (ii) are equivalent to $\mathrm{Hyp}(\Q(\mu_{p^\infty},V_f\otimes_\Qp V_g)$ of
%\cite[\S4.1]{loeffler-image}. As $f$ does not 
%








%Indeed, since our running assumptions imply that the newform $f$ is not of CM-type, hypothesis (i) is clear. On the other hand, hypothesis (ii) follows from \cite[Thm.~4.4.1]{loeffler-image}, noting that by \cite[\S{IV.2.2}]{serre-ladic} the quaternion algebra $B_f$ considered in \emph{loc.\,cit.} can be taken to be split.
%%
%%\textcolor{blue}{If $E$ had CM by $K$ and a rational $p$-isogeny, $p$ would have to be ramified in $K$ and therefore be a prime of bad reduction for $E$. Thus our assumptions imply $E$ does not have CM.}
%%
%%Then there exists a basis of $T_f$ for which $\operatorname{Im}({\rho_E})$ contains the subgroup of matrices whose reduction modulo $p^n$ for some $n>0$ is upper triangular and unipotent and, in particular, matrices of the form $\SmallMatrix{y &0 \\0&y^{-1}}$ for all $y\in 1+p^n\Z_p$. This follows from \cite[Theorem 1 and Remark 4.2.1]{greenberg-image} if $p>7$ or $p=5$ and from \cite[Corollary 1.11]{elladicimage} if $p=3$. The $p=7$ case is treated in \cite{greenbergrubin}, where the statement is shown to be true, with the exception of certain elliptic curves with CM by $\Q(\sqrt{-7})$. \textcolor{blue}{Maybe can just use Serre's open image to find the diagonal elements in the image?}
%%
%%We have (see for example \cite[Corollaire 2, p.300]{serre}) 
%%$\rho_E(\operatorname{Gal}(\bar{\Q}/\Q(\mu_{p^{\infty}})))=\rho_E(G_{\Q})\cap \SL(T_p(E))$.
%%
%%By definition $\alpha(g)=1$ for every $g\in \operatorname{Gal}(\bar{\Q}/K^{-}_{\infty}(\mu_{p^{\infty}}))$. Moreover $K^{-}_{\infty}\cap K(\mu_{p^{\infty}})=K$ and hence $\alpha(\operatorname{Gal}(K^{-}_{\infty}(\mu_{p^{\infty}})/K(\mu_{p^{\infty}})))=\alpha(\operatorname{Gal}(K^{-}_{\infty}/K))$. Thus there exists $g\in \operatorname{Gal}(K^{-}_{\infty}(\mu_{p^{\infty}})/K(\mu_{p^{\infty}}))$ such that $\alpha(g)= x \neq 1, x\in 1+p^N\Z_p$ and for any $\sigma\in \operatorname{Gal}(\bar{\Q}/K^{-}_{\infty}(\mu_{p^{\infty}}))$ we have $\alpha(\sigma\cdot g)= x$.
%%
%%Note that $K_{\infty}^{-}(\mu_{p^{\infty}})\cap \Q(E[p^{\infty}],\mu_{p^{\infty}})=\Q(\mu_{p^{\infty}})$, therefore $\rho_E(\operatorname{Gal}(\bar{\Q}/\Q(\mu_{p^{\infty}})))=\rho_E(\operatorname{Gal}(\bar{\Q}/K^{-}_{\infty}(\mu_{p^{\infty}})))$. 
%%
%%We deduce that there exists $\sigma \in \operatorname{Gal}(\bar{\Q}/K^{-}_{\infty}(\mu_{p^{\infty}}))$ such that
%%$\rho_E(\sigma)=\rho_E(h)^{-1}\SmallMatrix{x &0 \\0&x^{-1}},$
%%where $h\in \operatorname{Gal}(\bar{\Q}/K(\mu_{p^{\infty}}))$ is such that its image in $\operatorname{Gal}(K^{-}_{\infty}(\mu_{p^{\infty}})/K(\mu_{p^{\infty}}))$ is equal to $g$. Therefore we have $$\rho_E(h\cdot \sigma)=\SmallMatrix{x &0 \\0&x^{-1}}, \ \ \ \ \alpha(h\cdot \sigma) = \alpha(g)=x.$$
%%
%%Finally, since $\rm Ind(\alpha)_{| \operatorname{Gal}(\bar{\Q}/K(\mu_{p^{\infty}}))}=\alpha \oplus \alpha^c$, choosing $g$ so that it corresponds to a prime splitting in $K$, we obtain
%%$$\rho_E(h\cdot \sigma)=\SmallMatrix{x &0 \\0&x^{-1}}, \ \ \ \ \rho_{\rm Ind(\alpha)}(h\cdot \sigma) = \SmallMatrix{x &0 \\0&-x^{-1}}$$
%%and we have therefore found an element $h\sigma\in \operatorname{Gal}(\bar{\Q}/K(\mu_{p^{\infty}}))$ acting on $T$ with exactly one eigenvalue equal to one.
%\end{proof}



\section{Interlude: The rank one case and the general strategy}\label{sec:one}

In this section we give a proof of %Mazur's main conjecture, as in
Theorem~\ref{thm:CYC} %in the Introduction, 
under the following two additional hypotheses:
\begin{itemize}
\item[(a)] ${\rm corank}_{\Z_p}{\rm Sel}_{p^\infty}(E/K)=1$.
\item[(b)] The restriction map
\[
{\rm Sel}_{p^\infty}(E/K)\xrightarrow{{\rm loc}_v}E(K_v)\otimes\Q_p/\Z_p
\]
is nonzero.
\end{itemize}
%
%
%\begin{equation}\label{eq:cork1}
%\textrm{the $\Z_p$-corank of ${\rm Sel}_{p^\infty}(E/K)$ is $1$;}\tag{corank 1}
%\end{equation}
%\begin{equation}\label{eq:sur}
%\textrm{the restriction map ${\rm Sel}_{p^\infty}(E/K)\overset{{\rm loc}_w}\rightarrow{\rm im}(\delta_w)$ is surjective for all $w\mid p$},\tag{sur}
%\end{equation}
%where $\delta_w:E(K_w)\otimes\Q_p/\Z_p\rightarrow\rH^1(K_w,E[p^\infty])$ is the local Kummer map at $w$.
 
The short argument that follows, albeit independent from the discussion in the later sections, will allow us to motivate the more involved arguments needed for the proof of Theorem~\ref{thm:CYC} in general, and might provide some orientation to the reader.   
%Note that under the additional hypotheses (a) and (b), the results in the earlier sections of this paper and those in \cite{eisenstein} suffice for the proof.

Recall that $E_\bullet$ denotes the elliptic curve in the isogeny of $E$ constructed in \cite{wuthrich-int}. 
\sk

\noindent\emph{Step 1}. Under the above additional hypotheses, by  \cite[Thm.~C]{eisenstein} the module $\X_{\rm Gr}^S(E_\bullet/K_\infty^-)$ is $\Lambda^-_K$-torsion, with 
\begin{equation}\label{eq:ac-IMC}
\ch_{\Lambda^-_K}\bigl(\X_{\rm Gr}(E_\bullet/K_\infty^-)\bigr)\Lambda_K^{-,\rm ur}=\bigl(\Lcal_p^{{\rm BDP}}(f/K)\bigr).%\tag{ac}
\end{equation}

Conditions (a) and (b) above correspond to conditions (a) and (b) in Lemma~\ref{lem:coinv} with $\alpha=1$, and so from \eqref{eq:ac-IMC} we conclude that $\rH^1_{\Fcal_{\Gr}}(K,E_\bullet[p^\infty])$ is finite, and letting $\Fcal_{\rm Gr}(E_\bullet/K_\infty^-)\in\Z_p\llbracket{T}\rrbracket$  be a characteristic power series for $\X_{\rm Gr}(E_\bullet/K_\infty^-)$ we have
\begin{equation}\label{eq:ac-0}
\Fcal_{\rm Gr}(E_\bullet/K_\infty^-)(0)\;\sim_p\;\Lcal_p^{\rm BDP}(f/K)(0)\neq 0,%\tag{ac-0}
\end{equation}
where $\sim_p$ denotes equality up to a $p$-adic unit.
\sk

\noindent\emph{Step 2}. From Kato's result \cite[Thm.~17.4]{kato-euler-systems} (as refined in \cite[Thm.~16]{wuthrich-int} to an integral divisibility in the $p$-Eisenstein case) applied to $E_\bullet$ and $E_\bullet^K$, together with Proposition~\ref{prop:comp-Lcyc} and Proposition~\ref{prop:comp-Selcyc}, we deduce that $\X_{\rm ord}(E_\bullet/K_\infty^+)$ is $\Lambda_K^+$-torsion, and that we have the divisibility
\begin{equation}\label{eq:cyc-ord}
\ch_{\Lambda_K^+}\bigl(\X_{\rm ord}(E_\bullet/K_\infty^+)\bigr)\supset\bigl(\Lcal_p^{\rm PR}(E_\bullet/K)^{+}\bigr)%\tag{cyc-ord}
\end{equation}
in $\Lambda_K^+$. Moreover, the nonvanishing of $\Lcal_p^{\rm PR}(E_\bullet/K)^+$ follows from Rohrlich's nonvanishing results \cite{rohrlich-cyc}, while that of $\Lcal_p^{\rm Gr}(f/K)^+$ follows from (\ref{eq:ac-0}) and Proposition~\ref{prop:comp-Lac}, noting that 
\[
\Lcal_p^{\rm Gr}(f/K)^-(0)=\Lcal_p^{\rm Gr}(f/K)^+(0).
\] 
Therefore, by Proposition~\ref{prop:equiv} with $\alpha=1$ we obtain that $\X_{\rm Gr}(E_\bullet/K_\infty^+)$ is $\Lambda_K^+$-torsion, with the divisibility 
\begin{equation}\label{eq:cyc-Gr}
\ch_{\Lambda_K^+}\bigl(\X_{\rm Gr}(E_\bullet/K_\infty^+)\bigr)\supset\bigl(\Lcal_p^{\rm Gr}(f/K)^{+}\bigr)%\quad\textrm{in $\Lambda_K^{+,\rm ur}$}.%\tag{cyc}
\end{equation}
in $\Lambda_K^{+,\rm ur}$.
\sk

\noindent\emph{Step 3}. For $\Fcal_{\rm Gr}(E_\bullet/K_\infty^-)\in\Z_p\llbracket{T}\rrbracket$ a characteristic power series for $\X_{\rm Gr}^S(E_\bullet/K_\infty^-)$, we have the chain of relations 
\[
\Fcal_{\rm Gr}(E_\bullet/K_\infty^+)(0)\;\sim_p\;
\Fcal_{\rm Gr}(E_\bullet/K_\infty^-)(0)\;\sim_p\;\Lcal_p^{\rm BDP}(f/K)(0)\;\sim_p\;\Lcal_p^{\rm Gr}(f/K)^{-}(0)=\Lcal_p^{\rm Gr}(f/K)^+(0),
\]
using Proposition~\ref{prop:euler-char} with $\alpha=1$ (resp. Proposition~\ref{prop:factor-L-Gr}) for the first (resp. third) equality up to a $p$-adic unit. 
%
%with the equalities up to a $p$-adic unit following from Proposition~\ref{prop:euler-char} (with $\alpha=1$), relation (\ref{eq:ac-0}), and Proposition~\ref{prop:comp-Lcyc}, respectively. 
We thus conclude that
\[
\Fcal_{\rm Gr}(E_\bullet/K_\infty^+)(0)\;\sim_p\;\Lcal_p^{\rm Gr}(f/K)^{+}(0)\neq 0,
\]
which by easy commutative algebra (see \cite[Lem.~3.2]{skinner-urban}) implies that equality holds in (\ref{eq:cyc-Gr}):
\[
\ch_{\Lambda_K^+}\bigl(\X_{\rm Gr}(E_\bullet/K_\infty^+)\bigr)=\bigl(\Lcal_p^{\rm Gr}(f/K)^{+}\bigr).
\]
By Proposition~\ref{prop:equiv}, it follows that $\X_{\rm ord}(E_\bullet/K_\infty^+)$ is $\Lambda_K^+$-torsion, with
\[
\ch_{\Lambda_K^+}\bigl(\X_{\rm ord}(E_\bullet/K_\infty^+)\bigr)=\bigl(\Lcal_p^{\rm PR}(E_\bullet/K)^+\bigr),
\]
and by Proposition~\ref{prop:isog-inv} the same conclusion holds with $E$ in place of $E_\bullet$. By Proposition~\ref{prop:comp-Lcyc} and Proposition~\ref{prop:comp-Selcyc}, this equality of characteristic ideals together with Kato's divisibility for $E$ yields Theorem~\ref{thm:CYC} in this case. 
%(under the additional hypotheses (a) and (b) above).
\sk

In order to obtain Theorem~\ref{thm:CYC} in general:

\begin{itemize}
\item We shall prove Theorem~\ref{thm:AC} from the Introduction, removing the assumption ${\rm corank}_{\Z_p}{\rm Sel}_{p^\infty}(E/K)=1$ from  \cite[Thm.~C]{eisenstein}. 
%this is done in the next section. 
Then, %using the nonvanishing of $\Lcal_p(f/K)$, 
similarly as in \emph{Step 1} above, we shall obtain
\begin{equation}\label{eq:ac-alpha}
\Fcal_{\rm Gr}(E_\bullet(\alpha)/K_\infty^-)(0)\;\sim_p\;\Lcal_p^{\rm BDP}(f(\alpha)/K)(0)\neq 0,\nonumber%\tag{ac-$\alpha$}
\end{equation}
for any character $\alpha$ of $\Gamma_K^-$ %sufficiently close to $1$ 
away from the zeros of $\Lcal_p^{\rm BDP}(f/K)$ (so necessarily $\alpha\neq 1$ if the $\Z_p$-corank of ${\rm Sel}_{p^\infty}(E/K)$ is greater than $1$).

\item By arguments similar to those in \emph{Steps 2} and \emph{3} above, but complicated by the need to apply certain congruences and the use of $S$-imprimitive Selmer groups, we will show that for $\alpha$ sufficiently close to $1$, the module $\X_{\rm ord}(E_\bullet(\alpha)/K_\infty^+)$ is $\Lambda_K^+$-torsion, with
\[
\ch_{\Lambda_K^+}\bigl(\X_{\rm ord}(E_\bullet(\alpha)/K_\infty^+)\bigr)=\bigl(\Lcal_p^{\rm PR}(E_\bullet(\alpha)/K)^+\bigr).
\]
%Building on Theorem~\ref{thm:BF-ES} and results from $\S$\ref{sec:seltwist}, we shall deduce $\alpha$-twisted versions of the  divisibilities (\ref{eq:cyc-ord}) and (\ref{eq:cyc-Gr}) in \emph{Step 2} for an anticyclotomic character $\alpha$ as above sufficiently close to $1$.
%, getting in particular the divisibility
%(Note that $\alpha$ is anticyclotomic, whereas the divisibility
%
%that these divisibilities take place in the Iwasawa algebra of $\Gamma_K^+$, whereas $\alpha$ is a character of $\Gamma_K^-$, so the twisting theorems of \cite[Ch.\,{6}]{rubin-ES} don't apply.)
%
%\item Finally, building Theorem~\ref{thm:AC}, (\ref{eq:ac-alpha}), and the earlier results in this paper, we shall deduce a proof of the equality
%\[
%\ch_{\Lambda_K^+}\bigl(\X_{\rm ord}(E_\bullet(\alpha)/K_\infty^+)\bigr)=%\bigl(\Lcal_p^{\rm PR}(E_\bullet(\alpha)/K)^{+}\bigr).
%\]
%A congruence argument (taking $\alpha$ as above even closer to $1$ if needed) will then show that $\X_{\rm ord}(E_\bullet/K_\infty^+)$ and $\Lcal_p^{\rm PR}(E_\bullet/K)^+$ have the same Iwasawa invariants. This will yield the equality in (\ref{eq:cyc-ord}), as predicted by Conjecture\,\ref{conj:IMC-K}, from where the proof of  Theorem~\ref{thm:CYC} will follow as in \emph{Step 3}.
\end{itemize}

From this last equality (taking $\alpha$ to be sufficiently close to $1$), we deduce that the original $\mathfrak{X}_{\rm ord}(E_\bullet/K_\infty^+)$ and $\Lcal_p^{\rm PR}(E_\bullet/K)^+$ have the same Iwasawa invariants, which together with Kato's work will yield the proof of Theorem~\ref{thm:CYC}. 
%both Conjecture\,\ref{conj:IMC-K} for $\X_{\rm ord}(E/K)^+$ and $X_{\rm ord}(E/\Q_\infty)$.

%\section{Known divisibilities and rough strategy}

%Let $\alpha$ be an anticyclotomic character and $\mathbb{T}_{\alpha}:=T_p(E)(\alpha)\otimes \Lambda^+$. Consider the Selmer group $\X_{\rm str,\rm ord}:=(\rH^1_{\rm str, \rm ord}(K,T_p(E)(\alpha)\otimes (\Lambda^+)^{\vee}))^{\vee}$.


%\begin{equation}
%0 \to \rH^1_{\rm rel,\rm ord}(K,\mathbb{T}_{\alpha})/\Lambda^+\cdot BF \to \rH^1_{/\rm ord}(K_v,\mathbb{T}_{\alpha})/\loc_v(BF)\to \X(E(\alpha)/K)^+\to \X_{\rm str,\rm ord} \to 0
%\end{equation}
%\begin{equation}
%0 \to \rH^1_{\rm rel,\rm ord}(K,\mathbb{T}_{\alpha})/\Lambda^+\cdot BF \to \rH^1_{\rm ord}(K_{\bar{v}},\mathbb{T}_{\alpha})/\loc_{\bar{v}}(BF)\to X_{\rm Gr}(E(\alpha)/K)^+\to \X_{\rm str,\rm ord} \to 0
%\end{equation}

%\begin{itemize}
%\item non-zero class $BF \in \rH^1_{\rm rel,\rm ord}(K,\mathbb{T}_{\alpha})$
%\item explicit reciprocity law $\loc_v(BF) \leftrightsquigarrow \Lcal_p^{\rm PR}(E(\alpha)/K)^+$
%\item explicit reciprocity law $\loc_{\bar{v}}(BF) \leftrightsquigarrow \Lcal_p^{\rm Gr}(E(\alpha)/K)^+$
%\end{itemize}

%Use this and exact sequences above to prove:
%\begin{equation}\label{equivalenceMC}
%\tag{cyc-MCs}
%\begin{matrix}
%\Lcal_p^{\rm PR}(E(\alpha)/K)^+ \mid \ch(\X(E(\alpha)/K)^+) \Leftrightarrow \Lcal_p^{\rm PR}(E(\alpha)/K)^+ \mid \ch( \X_{\rm Gr}(E(\alpha)/K)^+),\\
%\ch(\X(E(\alpha)/K)^+) \mid \Lcal_p^{\rm Gr}(E(\alpha)/K)^+ \Leftrightarrow \ch( \X_{\rm Gr}(E(\alpha)/K)^+) \mid \Lcal_p^{\rm Gr}(E(\alpha)/K)^+\end{matrix}
%\end{equation}
%Kato's work \cite{kato-euler-systems} gives the divisibilities 
%\begin{equation}\label{katodivQ}
%\tag{Kato$/\Q$}
%\ch(\X(E/\Q) )  \mid \Lcal_p^{\rm MSD}(E/\Q) \ \ \text{ and } \ \ \ch(\X(E^K/\Q))\mid \Lcal_p^{\rm MSD}(E^K/\Q). 
%\end{equation}
%One also has a two-variable main conjecture for $E/K_{\infty}$, which predicts an equality $\Lcal(E/K_{\infty}) = \ch(X(E/K_{\infty}))$. 
%Using Proposition \ref{prop:comp-Lcyc} and Proposition \ref{prop:comp-Selcyc}, Kato divisibilities give the following 
%\begin{equation}\label{katodivK}
%\tag{Kato$/K$}
% \ch(\X(E/K)^+) \mid\Lcal_p^{\rm PR}(E/K)^+.
%\end{equation}



%One also has another two-variable main conjecture, that we will denote with the subscript $Gr$ and is instead related to the one-variable anticyclotomic BDP main conjecture. Our result \cite[Theorem  B]{eisenstein}, \textbf{under the corank 1 assumption}, combined with Proposition \ref{prop:comp-Lac}, gives an equality
%\begin{displaymath}
%\Lcal_p^{\rm Gr}(E/K)^- = \ch(\X_{\rm Gr}(E/K)^-).
%\end{displaymath} 

%In particular, evaluating at $T=0$, which is the intersection of the cyclotomic and anticyclotomic line, we get
%\[
%\Lcal_p^{\rm Gr}(E/K)^+(0) = \ch(\X_{\rm Gr}(E/K)^+)(0),
%\] 
%which is not zero if the rank over $K$ is one.

%This in particular implies, together with the claim \eqref{equivalenceMC} for $\alpha=1$ and Kato's divisibility \eqref{katodivK}, 
%$$\Lcal_p^{\rm PR}(E/K)^+=\ch(\X(E/K)^+).$$
%From this we can deduce the cyclotomic main conjecture for $E/\Q_{\infty}$ and $E^K/\Q_{\infty}$ \textbf{in the case of rank $\leq 1$} : indeed we can use again Proposition \ref{prop:comp-Lcyc} and Proposition \ref{prop:comp-Selcyc} to obtain
%\[
%\ch(\X(E/\Q) )\cdot  \ch(\X(E^K/\Q))  = \Lcal_p^{\rm MSD}(E/\Q) \cdot \Lcal_p^{\rm MSD}(E^K/\Q),
%\]
%which combined with the known divisibilities \eqref{katodivQ}, gives $\ch(\X(E/\Q) ) = \Lcal_p^{\rm MSD}(E/\Q) $ and $\ch(\X(E^K/\Q))  =  \Lcal_p^{\rm MSD}(E^K/\Q)$. (in particular we reprove the Greenberg--Vatsal result, since if an elliptic curve satisfies \eqref{assumption}, the quadratic twist does not).


%We can use the same strategy in the higher rank case. Steps
%\begin{itemize}
%\item remove the corank 1 assumption \cite[Theorem  B]{eisenstein}, $\S$\ref{sec:anticyc};
%\item instead of using it at the trivial character, we use it at anticyclotomic characters $\equiv 1 \mod p^t$ for $t\gg 0$;
%\item need one divisibility $\ch_{\Lambda^+}(\X_{\rm str,\rm ord})\mid \ch_{\Lambda^+}(\rH^1_{\rm rel,\rm ord}(K,\mathbb{T}_{\alpha})/\Lambda^+\cdot BF)$ coming from Beilinson--Flach classes, $\S$\ref{sec:BF}.
%\end{itemize}


\section{Anticyclotomic main conjecture}\label{sec:anticyc}

The key new result in this section is Theorem~\ref{thm:Zp-twisted}. The result is a Kolyvagin system bound complementing \cite[Thm.~3.2.1]{eisenstein} for characters $\alpha$ of $\Gamma_K^-$ that are close to $1$. We then use this result to obtain a version of \cite[Thm.~C]{eisenstein} and its corollaries removing the assumption that ${\rm corank}_{\Z_p}{\rm Sel}_{p^\infty}(E/K)=1$ (i.e., assumption (Sel) in \emph{loc.\,cit.}).  
%As explained in the preceding section, removing this rank one assumption is essential to our approach to Theorem\,\ref{thm:CYC}. The key new result in this section is Theorem~\ref{thm:Zp-twisted}.

Throughout this section, we let $E/\Q$ be an elliptic curve of conductor $N$, $p\nmid 2N$ be a prime of good ordinary reduction for $E$, and $K$ be an imaginary quadratic field of discriminant $D_K$ prime to $Np$. We assume that
\begin{equation}\label{eq:h1}
E(K)[p]=0.\tag{h1}
\end{equation}
%and denote by $\mathscr{L}=\mathscr{L}_E$ the set of primes $\ell\nmid N$ that are inert in $K$ and satisfy $a_\ell\equiv\ell+1\equiv 0\pmod{p}$, where $a_\ell=\ell+1-\vert\tilde{E}(\mathbb{F}_\ell)\vert$, and by $\mathscr{N}$ the set of square-free products of primes $\ell\in\mathscr{L}$. 

%Theorem 3.2.1 of \cite{eisenstein} was the key result to deduce, assuming \eqref{eq:h1}, one of the divisibilities in Conjecture~\ref{conj:PR}: 
%\begin{center}
%${\rm char}_{\Lambda^-_{K}}\bigl(\X_{\rm ord}(E/K_\infty^-)_{\rm tors}\bigr)$ divides ${\rm char}_{\Lambda^-_{K}}\bigl(\mathfrak{S}_{\rm ord}(E/K_\infty^-)/(\kappa_1^{\rm Hg})\bigr)^2 $ in 
%$\Lambda^-_K[\frac{1}{p},\frac{1}{(\gamma^- -1)}]$.
%\end{center}
%The proof of this result relies on a Kolyvagin system argument for twists by anticyclotomic characters sufficiently close $p$-adically to the ``zeroes of $f_{\Lambda^-_K}$'', where $f_{\Lambda^-_K}$ is a generator of ${\rm char}_{\Lambda^-_{K}}\bigl(\mathfrak{S}_{\rm ord}(E/K_\infty^-)/(\kappa_1^{\rm Hg})\bigr)$. This allows to prove that for $m\gg 0$ and for every height one prime ideal $\mathfrak{P}=(g)$ of $\Lambda^-_K$ dividing $f_{\Lambda^-_K}$ we have
%\begin{equation}\label{eq:oldbound}
%m\cdot{\rm ord}_{\mathfrak{P}}\bigl({\rm char}_{\Lambda^-_{K}}\bigl(\X_{\rm ord}(E/K_\infty^-)_{\rm tors}\bigr)\bigr)\leqslant 2m\cdot{\rm ord}_{\mathfrak{P}}(f_{\Lambda^-_K})+ \mathcal{E}_{\mathfrak{P},m}
%\end{equation}
%for an `error term' $\mathcal{E}_{\mathfrak{P},m}$ depending on $T_p(E)$, ${\rm rk}_{\Z_p}(\Lambda^-/\mathfrak{P})$ and linearly on the $p$-adic valuation of $g(0)+p^m$. In particular, if $\mathfrak{P}=\mathfrak{P}_0=(\gamma^- -1)$ divides $f_{\Lambda^-_K}$, since $\mathcal{E}_{\mathfrak{P}_0,m}$ depends linearly on $m$, we cannot let $m$ go to infinity and deduce ${\rm ord}_{\mathfrak{P}_0}\bigl({\rm char}_{\Lambda^-_{K}}\bigl(\X_{\rm ord}(E/K_\infty^-)_{\rm tors}\bigr)\bigr)\leqslant 2\;{\rm ord}_{\mathfrak{P}_0}(f_{\Lambda^-_K})$. Therefore, in this section we prove a variation of \cite[Theorem 3.2.1]{eisenstein} which applies only to characters $\alpha$ $p$-adically close to the trivial character (say $\alpha \equiv 1 \mod \varpi^m$) and involves an error term not depending on $m$. The strategy is similar to the one employed in \cite[$\S$3.3.3]{eisenstein} but technically more delicate. In order to remove the dependence on $m$, in $\S$\ref{seccheb} we tweak the \v{C}ebotarev argument to produce the Kolyvagin primes at the cost of having to work in $\S$\ref{secproofkoly} with the $\varpi^m$-torsion of the Selmer groups involved (see Remarks \ref{rmkcalpha}, \ref{rmkkolyvaginargument} for more details). We then obtain \eqref{eq:oldbound} for $\mathfrak{P}=\mathfrak{P}_0$ and an error term $\mathcal{E}$ such that $\mathcal{E}/m$ goes to zero as $m$ goes to infinity. This allows us to deduce Theorem \ref{thm:AC}, which is proved in $\S$\ref{subsec:ac-IMC}.


\subsection{A Kolyvagin-style bound for $\alpha\equiv 1\;({\rm mod}\,\varpi^m)$}
Let $R$ be the ring of integers of some finite extension $\Phi/\Qp$ and let $\varpi\in R$ be a uniformizer.
Let $\fm = (\varpi)$ be the maximal ideal of $R$.
Let $\alpha:\Gamma_K^- \to R^{\times}$ be an anticyclotomic character and $m$ a positive integer such that 
\begin{equation}\label{eq:pm}
\alpha\equiv 1\;({\rm mod}\,\varpi^m).
\end{equation}

Denote by $\rho_E:G_\Q\rightarrow{\rm Aut}_{\Z_p}(T_pE)$ the representation on the $p$-adic Tate module of $E$, and consider the $G_K$-modules 
\begin{equation}
T:=T_pE\otimes_{\Z_p}R(\alpha),\quad
V:=T\otimes_R \Phi,\quad A:=T\otimes_R\Phi/R \simeq V/T,\nonumber
\end{equation}
where $R(\alpha)$ is the free $R$-module of rank one on which $G_K$ acts via the composition of the projection $G_K\twoheadrightarrow \Gamma_K^-$ with
$\alpha$,  and the $G_K$-action on $T$ is via $\rho = \rho_E\otimes \alpha$.
%To ease notation, let $(T,\mathcal{F},\cL)$ denote the Selmer triple $(T,\Fcal_{\rm ord},\cL_E)$. 

Let $\mathscr{L}=\mathscr{L}_E$ be the set of primes $\ell\nmid N$ that are inert in $K$ and satisfy $a_\ell\equiv\ell+1\equiv 0\pmod{p}$, where $a_\ell=\ell+1-\vert\tilde{E}(\mathbb{F}_\ell)\vert$, and let $\mathscr{N}$ denote the set of square-free products of primes $\ell\in\mathscr{L}$. For each $\ell\in\mathscr{L}$, let $I_{\ell}\subset\Z_p$ be the smallest ideal containing $\ell+1$ for which the Frobenius element $\frob_{\lambda} \in G_{K_{\lambda}}$ acts trivially on $T / I_{\ell} T$, where $\lambda\mid \ell \in \mathscr{L}$. 
For any $k\geqslant 0$, let 
\[
T^{(k)}=T/\varpi^kT,\quad\cL^{(k)}=\{\ell\in\cL\colon I_\ell\subset p^k\Z_p\},
\]
and let $\mathscr{N}^{(k)}$ be the set of square-free products of primes $\ell\in\cL^{(k)}$. We refer the reader to \cite[\S{3.1}]{eisenstein} for the definition of the module of Kolyvagin systems $\mathbf{KS}(T,\Fcal_{\rm ord},\cL)$ associated to the triple $(T,\Fcal_{\rm ord},\cL)$ (here and until $\S\ref{subsec:ac-IMC}$, $\Fcal_{\rm ord}$ denotes the ordinary Selmer structure introduced in \cite[p.\,548]{eisenstein}, which is compatible with the discrete coefficients analogue defined in $\S$\ref{sec:seltwist} and Definition \ref{defilocalconds}). 

\begin{thm}\label{thm:Zp-twisted}  
There exist non-negative integers $\CM$ and $\CE$ depending only on $T_pE$ and $\operatorname{rank}_{\Z_p}(R)$ such that if
$m\geq \CM$ and if there is a Kolyvagin system $\kappa=\{\kappa_{n}\}_{n\in\cN}\in\mathbf{KS}(T,\Fcal_{\rm ord},\cL)$ with $\kappa_{1}\neq 0$,
then $\rH^1_{\Fcal_{\rm ord}}(K,T)$ has $R$-rank one and there is a finite $R$-module $M$ such that
\[
{\rm H}^1_{\Fcal_{\rm ord}}(K,A)\simeq \Phi/R \oplus M\oplus M
\]
with 
\[
{\rm length}_R(M)\leqslant {\rm length}_R\bigl({\rm H}^1_{\Fcal_{\rm ord}}(K,T)/R\cdot\kappa_{1}\bigr)+\CE.
\]
\end{thm}
\noindent The `error term' $\CE$ in this theorem is independent of $m$, but that comes at the expense of the result applying only to characters $\alpha$ that are sufficiently close to $1$ (as measured by $\CM$). The reader may wish to compare this theorem
with \cite[Thm.~3.2.1]{eisenstein} whose error term $E_\alpha$ is at least as large as $m$ but which also applies to $\alpha$ that are relatively far from $1$. Both results are crucial for the proof of Theorem \ref{thm:AC}.


\subsection{Structure of Selmer groups}

In the following we use $\Fcal$ to denote $\Fcal_{\rm ord}$ for simplicity. For any $k\geq 1$, let $R^{(k)}= R/\fm^k$.
We recall the following structure results:
% (see \cite[Lemma 3.3.1, Proposition 3.3.2]{eisenstein} :
\begin{lemma}%[{\cite[Lem.~3.3.1]{eisenstein}}]
\label{lemmamod} 
For every $n\in\cN^{(k)}$ and $0\leqslant i\leqslant k$ there are natural isomorphisms
\begin{equation}\label{eq2}
\rH^1_{\Fcal(n)}(K,T^{(k)}/\fm^iT^{(k)})\xrightarrow{\sim}\rH^1_{\Fcal(n)}(K,T^{(k)}[\fm^i])\xrightarrow{\sim}{\rm H}_{\Fcal(n)}^1(K,T^{(k)})[\fm^i]\nonumber
\end{equation}
induced by the maps $T^{(k)}/\fm^iT^{(k)}\xrightarrow{\pi^{k-i}}T^{(k)}[\fm^i]\rightarrow T^{(k)}$.
\end{lemma}

\begin{proof} See \cite[Lem.~3.3.1]{eisenstein}.
\end{proof}

\begin{prop}%[{\cite[Prop.~3.3.2]{eisenstein}}]
\label{propstructure} There is an integer $\epsilon\in\{0,1\}$ such that for all $k$ and every
every $n\in\cN^{(k)}$ there is an $R^{(k)}$-module $M^{(k)}(n)$ such that
\begin{equation}\label{eq1}
\rH^1_{\CF(n)}(K,T^{(k)}) \simeq (R/\fm^k)^\epsilon\oplus M^{(k)}(n)\oplus M^{(k)}(n).\nonumber
\end{equation}
%Moreover, $\epsilon$ can be taken to be $0$ or $1$ and is independent of $k$ and $n$.
\end{prop}

\begin{proof} See \cite[Prop.~3.3.2]{eisenstein}.
\end{proof}

%We want $k$ greater than that and since we also want $k$ such that $\ind(\kappa_{\alpha_m,1},\rH^1_{\Fcal}(K,T_m))=\ind(\kappa_{\alpha_m,1},\rH^1_{\Fcal}(K,T_m^{(k)}))$, i.e. $k> m\cdot s_{\Lambda}$, where $s_{\Lambda}=\ord_{(T)}(f_\Lambda)$, we take
%\begin{itemize}
%\item $m\gg 0$ so that $m\cdot s_{\Lambda} \gg 4(C_1+C_2)(\dim_{\mathbb{F}_p}\rH^1_{\Fcal}(K, E[p]))$;
%\item $k\gg m\cdot s_{\Lambda}$ or $k\gg  4(C_1+C_2)(\dim_{\mathbb{F}_p}\rH^1_{\Fcal}(K, E[p]))$ if $s_{\Lambda}=0$.
%\end{itemize}

By Lemma \ref{lemmamod} and \eqref{eq:pm}, if $k\geq m$ there is an isomorphism
\begin{equation}\label{eq:pmtorsion}
\rH^1_{\Fcal(n)}(K,T_E^{(m)})\simeq \rH^1_{\Fcal(n)}(K,T^{(k)})[\fm^{m}],
\end{equation}
where $T_E^{(m)}=T_p(E)\otimes_{\Z_p} R/\fm^m$.  We can then exploit the action of complex conjugation on the left hand side
(both $T_E^{(m)}$ and the Selmer structure $\Fcal(n)$ are stable under this action).   We make use of this in our subsequent analysis
of the structure of the $R$-modules $M^{(k)}(n)$ in terms of Kolyvagin classes. 


%we can use the fact that $\alpha\equiv 1 \mod \varpi^m$ and, by Lemma \ref{lemmamod}, the isomorphism 
%\begin{equation}\label{eq:pmtorsion}
%\rH^1_{\Fcal}(K,T_E^{(m)})\simeq \rH^1_{\Fcal}(K,T^{(k)})[\fm^{m}],
%\end{equation}
%where $T_E^{(m)}=T_p(E)\otimes R/\fm^m$.
%We can then exploit the action of the complex conjugation on the left hand side. 


\subsection{The \v{C}ebotarev argument}\label{seccheb}
We recall the definitions of the error terms $C_1,C_2$ of \cite[$\S$3.3.1]{eisenstein}. For $U = \Z_p^\times\cap \mathrm{im}(\rho_E)$ let
\[
C_1 := \min\{v_p(u-1)\colon u\in U\}.
\]
As $U$ is an open subgroup, $C_1<\infty$.
Recall also that $\End_{\Z_p}(T_pE)/\rho_E(\Z_p[G_{\Q}])$ is a torsion $\Z_p$-module and let
\[
C_2:=\min\bigl\{ n\geqslant 0 \colon p^n\End_{\Z_p}(T_pE)\subset\rho_E(\Z_p[G_{\Q}])\bigr\}.
\]
Let $r=\operatorname{rank}_{\Z_p}R$ and 
\[
e:=r(C_1+C_2).
\]

For any finitely-generated torsion $R$-module $M$ and $x\in M$, let
$$
\ord(x):=\min\{m\geqslant 0: \varpi^m\cdot x =0\} \ \ \text{and } \ \ \ \exp(M):=\max\{\ord(x):x\in M\}.
$$
The following result %is the analogue of \cite[Prop. 3.3.6]{eisenstein} and 
is one of the main tools for our proof of Theorem \ref{thm:Zp-twisted}. 
%The difference with the result in \emph{op. cit.} is that here we are working on the $\fm^m$-torsion of the Selmer groups of $T^{(k)}$ and, by \eqref{eq:pmtorsion}, the character $\alpha$ becomes trivial and we can directly work with the eigenspaces of the complex conjugation (whereas in \cite[Prop. 3.3.6]{eisenstein} we had to go over an extension trivialising the character and use ``some quadratic forms'' to achieve linear independence of the images). 
\begin{prop}\label{prop:prime2} 
Let $c^\pm \in \rH^1(K,T_E^{(m)})^\pm$.  Let $k\geq m$. 
Then there exist infinitely many primes $\ell\in \cL^{(k)}$ such that 
$$
\ord(\loc_\ell(c^\pm)) \geqslant \ord(c^\pm) - e.
$$
In particular, $R\cdot \loc_\ell(c^+)+ R\cdot \loc_\ell(c^-)$ has an $R$-submodule isomorphic to  
$$
R/\fm^{\max\{0,\ord(c^+)-e\}}\oplus  R/\fm^{\max\{0,\ord(c^-)-e\}}.
$$
\end{prop}

%The difference between this proposition and \cite[Prop. 3.3.6]{eisenstein} is primarily that here we have restricted ourselves to the the $\fm^m$-torsion of the Selmer groups of $T^{(k)}$  (see \cite{eq:pmtorsion}) and we can directly work with the eigenspaces of the complex conjugation.
%For the proof of {\em loc.~cit.} we worked over an extension trivialising the character $\alpha \mod \varpi^k$, and then used ``some quadratic forms'' to 
%estimate the linear independence of the images of the localisations of the classes. 

\begin{proof}
The proof of this proposition follows along the lines of that of \cite[Prop. 3.3.6]{eisenstein}. 

%For any $n>0$, let $U_n\subset \Aut_{\Z_p}(T_E^{(n)})$ be the image of $G_K$ under $\rho_E \mod \varpi^n$.
Let $u\in  \Z_p^\times\cap \mathrm{im}(\rho_E|_{G_{K_\infty}})= \Z_p^\times\cap \mathrm{im}(\rho_E)\subset \Z_p^\times\cap \mathrm{im}(\rho_E \otimes \alpha)$ (see \cite[Lemma 3.3.3]{eisenstein}) be such that $v_p(u-1) = C_1$.   Let $L$ be 
%the fixed field of the action of $G_K$ on $T/p^kT$, so $L$ is 
the composite of the fixed field of the action
of $G_K$ on $T_E/p^kT_E$ and the field $K_\alpha$ trivialising $\alpha \mod p^k$.  Then there is some $h$ in the center of $\Gal(L/K)$ 
such that $h$ acts on $T/p^kT$, and hence on $T^{(m)}\simeq T_E^{(m)}$, as multiplication by $u$. The kernel of the 
restriction map $\rH^1(K,T^{(m)})\rightarrow \rH^1(L,T^{(m)})$ is $\rH^1(L/K,T^{(m)})$ and it follows from the existence
of $h$ that the latter is annihilated by $u-1$ (this is essentially Sah's Lemma: if $c:\Gal(L/K)\rightarrow T^{(m)}$ is 
a $1$-cocycle, then $c(gh) = c(hg)$ and so $(h-1)c(g) = (g-1)c(h)$ and hence $p^{C_1}c$ is a coboundary). It follows that 
\begin{equation}\label{eq:kerm}
p^{C_1}\cdot \ker \bigl( \rH^1(K,T_E^{(m)})\to \rH^1(L, T_E^{(m)})\bigr) = 0.
\end{equation}




%For any $n>0$, 
%it follows from Sah's Lemma, as in the proof of \cite[Prop. 3.3.6]{eisenstein}, that $u-1$ annihilates $\ker \bigl( \rH^1(K,T^{(n)})\to \rH^1(L_nF, T^{(n)})\bigr)$ for any $F\subset K_\infty$ trivialising $\alpha$ modulo $\fm^k$,
%where $L_n$ is the fixed field of the kernel of the action of $G_K$ on $T_E^{(n)}$. In particular,
%\textcolor{red}{
%\begin{equation}\label{eq:kern}
%p^{C_1}\cdot \ker \bigl( \rH^1(K,T^{(n)})\to \rH^1(L_nF, T^{(n)})\bigr) = 0 \ \ \text{for all $n$}.
%\end{equation}
%Let $L$ be the fixed field of the action of $G_K$ on $T/p^kT$, so $L$ is the composite of $L_{k'}$ for $k' = k\cdot \ord_\varpi(p)$ and of the field $K_\alpha$ trivialising $\alpha \mod \varpi^k$. 
%As $E[p]^{G_K} =0$ by \eqref{eq:h1}, $\rH^1(K,T_E^{(m)}) \subset \rH^1(K,T^{(k)})$, and so 
%$$
%\ker \bigl( \rH^1(K,T_E^{(m)})\to \rH^1(L, T_E^{(m)})\bigr)\subset \ker \bigl( \rH^1(K,T^{(k)})\to \rH^1(L, T^{(k)})\bigr).
%$$
%It then follows from \eqref{eq:kern} that 
%\begin{equation}\label{eq:kerm}
%p^{C_1}\cdot \ker \bigl( \rH^1(K,T_E^{(m)})\to \rH^1(L, T_E^{(m)})\bigr) = 0.
%\end{equation}
%}

Let $d^\pm:= \ord(c^{\pm}) - r(C_1 + C_2)$. If $d^+=d^-\leq 0$, then there is nothing to prove. So assume at least one of $d^\pm$ is positive. 
By (\ref{eq:kerm}), the kernel of the restriction map
$$
\rH^1(K,T_E^{(m)}) \stackrel{res}{\rightarrow} \rH^1(L,T_E^{(m)}) = \Hom_{G_K}(G_L,T_E^{(m)})
$$
is annihilated by $p^{C_1}=\varpi^{rC_1}$. Let $f^\pm\in \Hom_{G_K}(G_L,T_E^{(m)})$ be the image of $c^\pm$. We then have 
$$
\ord(f^\pm) \geq \ord(c^\pm) - rC_1.
$$
As $f^\pm(G_L)$ is a $G_K$-submodule, $f^\pm(G_L) = \im(\rho_E)\cdot f^\pm(G_L)$ and so, by the definition of $C_2$,
the image of $f^\pm$ contains $p^{C_2}\operatorname{End}(T_p(E))\cdot f^\pm(G_L)$. 
Since $\ord(f^\pm) \geq \ord(c^\pm) - rC_1$, it follows that the $R$-span of the image of $f^\pm$ contains $\varpi^{m-\ord(c^\pm) + r(C_1+C_2)} T_E^{(m)}$. 
Since at least one of $d^+$ and $d^-$ is positive, 
it follows that at least one of $f^+$ and $f^-$ is non-trivial. 

Let $H\subset G_L$ be the intersection of the kernels of $f^+$ and of $f^-$, and let $Z = G_L/H$. Note that $H\neq G_L$ since 
some $f^\pm$ is non-trivial, so $Z$ is a non-trivial torsion $R$-module. Note also that $Z$ is stable under the action of complex conjugation since each $f^\pm$ is.  In particular, $Z$ decomposes into eigenspaces under the action of complex conjugation: $Z= Z^+\oplus Z^-$. 

Let  $g^\pm$ be the projection of $f^\pm$ to the summand $(T_E^{(m)})^{\pm} \cong R/\fm^m$.
Then the $R$-span of the image of $g^\pm$ contains an $R$-submodule isomorphic to $R/\fm^{\max\{0,d^\pm\}}$. 
We have $g^\pm(Z^-)=0$ since $f^\pm\in \Hom(G_L,E[p^m])^{\pm}$, which means $f^\pm(Z^-)\subset E[p^m]^{\mp}$. So we find $g^{\pm}(Z) = g^{\pm}(Z^+)$ and
that the $R$-span of $g^\pm(Z^+)$ contains a submodule isomorphic to $R/\fm^{\max\{0,d^\pm\}}$. It follows
that $Z^+$ is non-trivial.

If $d^\pm>0$, let $W_\pm \subset Z^+$ be the proper subgroup such that $g^\pm(W_\pm) = \varpi^{m-(d^\pm-1)}(T_E^{(m)})^\pm$. 
If $d^\pm \leq 0$, let $W_\pm  = 0$. Then both $W_+$ and $W_-$ are proper subgroups of $Z^+$ (since there exists some $z\in Z^+$ such that $g^\pm(z)\in \varpi^{m-d^\pm}(T_E^{(m)})^\pm$). It follows that
$W_+\cup W_-\neq Z^+$.  Let $z\in Z^+$, $z\not\in W_+\cup W_-$. By the definition of $W^\pm$, we have
\begin{equation}\label{eq:ordnek}
\ord(g^\pm(z)) \geq d^\pm.
\end{equation}


Let $M =  \overline{\Q}^{H}$, so $\operatorname{Gal}(M/L) = Z$. Let $g = \tau z \in G_\Q$, and let 
$\ell\nmid Np$ be any prime such that both $c^+$ and $c^-$ are unramified at $\ell$ and $\frob_\ell = g$ in $\operatorname{Gal}(M/\Q)$. 
The \v{C}ebotarev density theorem implies there are infinitely many such primes.
%\textcolor{red}{Since $Z$ fixes $E[p^k]\supset E[p^m]$ and $K$, $\frob_\ell$ acts as $\tau$ on both $E[p^k]$ and $K$.
%This means that $a_\ell(E) \equiv \ell +1 \equiv 0 \mod p^k$ and $\ell$ is inert in $K$. Since $\alpha_{|Z}=1$, we obtained $\ell\in \mathscr{L}^{(k)}$.}
 Since $Z$ fixes $L$ and since $L$ contains the fixed field of the $G_K$-action on $E[p^k]$,
$\frob_\ell$ acts as $\tau$ on both $E[p^k]$ and $K$.
This means that $a_\ell(E) \equiv \ell +1 \equiv 0 \mod p^k$ and $\ell$ is inert in $K$. 
Since $L$ also contain the fixed field of $\alpha\mod p^k$, for $\lambda\mid \ell$ a prime of $K$, it follows that  $\frob_\lambda$
acts trivially on $T/p^kT$ and hence that $\ell\in \mathscr{L}^{(k)}$. 

Since $\ell$ is inert in $K$, the Frobenius element of $\ell$ in $\operatorname{Gal}(\bar{\Q}/K)$ is $\frob_\ell^2$. Consider the restriction of
$c^\pm$ to $K_\ell$. Since $c^\pm$ is unramified at $\ell$, $\loc_\ell(c^\pm)$ is completely determined by the image $c^\pm(\frob_\ell^2)$ in $T_E^{(m)}/(\frob_\ell^2-1)T_E^{(m)}$. By the choice of $\ell$, $\frob_\ell^2$ acts trivially on
$T_E^{(m)}$, so  $T_E^{(m)}/(\frob_\ell^2-1)T_E^{(m)} = T_E^{(m)}$. Moreover, $\frob_\ell^2 = g^2  = z^2 \in \operatorname{Gal}(M/L)$,
so $c^\pm(\frob_\ell^2) = f^\pm(z^2) = 2g^\pm(z)$, where the second equality follows from the fact that the projection of $f^\pm$ to $(T_E^{(m)})^\mp$ maps $z\in Z^+$ to zero. Since $p$ is odd, (\ref{eq:ordnek}) yields $\ord(\loc_{\ell}(c^\pm)) = 
\ord(c^\pm(\frob_\ell^2)) = \ord(2g^\pm(z)) = \ord(g^\pm(z)) \geq d^\pm$.
\end{proof}

\begin{rmk}\label{rmkcalpha}
The primary difference between Proposition \ref{prop:prime2} and \cite[Prop.~3.3.6]{eisenstein} is  that here we have restricted ourselves to the $\fm^m$-torsion of the Selmer groups of $T^{(k)}$  (see (\ref{eq:pmtorsion})) and so we can directly work with the eigenspaces of complex conjugation.
For the proof of {\em loc.~cit.} we worked over an extension trivialising the character $\alpha \mod \varpi^k$ and then used ``some quadratic forms'' to 
estimate the linear independence of the images of the localisations of the classes.  The upshot is that our error term no longer involves
the $C_{\alpha}$ of \cite{eisenstein}. This is crucial for removing the corank one assumption in the proof of the anticyclotomic Iwasawa main conjecture in {\em op.~cit.}. However it causes some complications in the proof of Theorem \ref{thm:Zp-twisted}: we use the $\fm^m$-Selmer groups to control the 
image of the localisation at Kolyvagin primes of classes in the $\fm^k$-Selmer groups, and the resulting control is not as tight as in \cite{eisenstein}.
 \end{rmk}



\subsection{Proof of Theorem \ref{thm:Zp-twisted}}\label{secproofkoly}
The (co-)rank one claim in the theorem follows from \cite[Thm.~3.3.8]{eisenstein}: 
$$
\rH^1_\CF(K,T)\simeq R \ \ \text{and} \ \ \rH^1_\CF(K,A) \simeq \Phi/R \oplus M, \ \ M \simeq M_0\oplus M_0
$$
for some finitely-generated torsion $R$-module $M_0$ such that $M_0\simeq M^{(k)}(1)$ for all $k\gg 0$.
In fact, the proof of  \cite[Thm.~3.3.8]{eisenstein} shows that $M_0\simeq M^{(k)}(1)$ if $k>{\ind(\kappa_{1})+3r(C_1+C_2+m)}$.
In the current setting, the error term $C_{\alpha}$ of \emph{op. cit.} is equal to $m$; it is essentially this fact that prevents the arguments in {\em op.~cit.} from applying to prove the theorem in the current setting and is the reason we take a different approach below
to establish the bound 
\begin{equation}\label{eq:bound}
\tag{B}
s_1+\CE\geq {\rm length}_{R}(M^{(k)}(1)),
\end{equation}
where $s_1=\ind(\kappa_{1},\rH^1_{\Fcal}(K,T))$ and $\CE$ does not depend on $m$, provided $m$ is sufficiently large. 
As $\mathrm{\length}_R(M) = \mathrm{length}_R(M^{(k)}(1))$ for $k\gg 0$, the bound \eqref{eq:bound} implies 
the bound in Theorem \ref{thm:Zp-twisted}.



%First, we show that the rank one statement follows from results in \cite{eisenstein}. 
%Since $\rH^1_\CF(K,A) = \varinjlim_k \rH^1_\CF(K,T^{(k)})$, it follows that 
%\begin{displaymath}
%\rH^1_\CF(K,A)[\varpi^k] \simeq \rH^1_\CF(K,T^{(k)}).
%\end{displaymath}
%If $\epsilon =1$ for some $k$, then $\epsilon = 1$ for all $k$ and the corank of $\rH^1(K,A)$ is at least one. 
%If moreover, $M^{(k)}(1)$ has all exponents less than $k$ (that is, $\ord(x) < k$ for all $x\in M^{(k)}(1))$, then 
%it follows from the above displayed isomorphism that the corank equals one, and in fact, 
%$$
%\rH^1_\CF(K,A) \simeq \Phi/R \oplus M, \ \ M \simeq M_0\oplus M_0
%$$
%for some finitely-generated torsion $R$-module $M_0$ such that $M_0\simeq M^{(k')}(1) \simeq M^{(k)}(1)$ for all $k'\geq k$.
%That such a $k$ in fact exists follows from \cite[Theorem 3.3.8]{eisenstein}, which tells us more precisely that for all $k\gg 0$,  
%$\varpi^{\ind(\kappa_{1})+3(C_1+C_2+m)} \rH^1_\CF(K,T^{(k)})$ is cyclic and generated by $\varpi^{3(C_1+C_2+m)}\kappa_{1}$.
%
%
%In particular, if for some $k$ big enough, $\epsilon=1$ (which than gives us ``rank 1'' for every $k$) and $M^{(k)}(1)$ has exponents strictly smaller than $k$, then
%$$
%\rH^1_\CF(K,A) \simeq \operatorname{Frac}(R)/R \oplus M, \ \ M \simeq M_0\oplus M_0.
%$$
%for some finitely-generated torsion $R$-module $M_0$ such that $M_0\simeq M^{(k)}(1) \simeq M^{(k_0)}(1)$ for all $k_0\geq k$. This result follows from \cite[Theorem 3.3.8]{eisenstein}, which tells us more precisely that for all $k\gg 0$, we have that $\varpi^{\ind(\kappa_{\alpha,1})+3(C_1+C_2+m)} \rH^1_\CF(K,T^{(k)})$ is cyclic and generated by $\kappa_{\alpha,1}$. In the current setting, the error term $C_{\alpha}$ of \emph{op. cit.} is equal to $m$. This is the reason why we cannot go beyond this annihilation result with our previous method and we need a different approach to prove the bound
%\begin{equation}\label{eq:bound}
%\tag{B}
%s_1+E\geq {\rm length}_{\Z_p}(M^{(k)}(1)),
%\end{equation}
%where $s_1=\ind(\kappa_{\alpha,1},\rH^1_{\Fcal}(K,T))=\ind(\kappa_{\alpha,1},\rH^1_{\Fcal}(K,T^{(k)}))$ and $E$ does not depend on $C_{\alpha_m}=m$.

We now focus on the proof of \eqref{eq:bound}. 
% for some choice of $k\gg 0$ as above, which, as explained above, will conclude the proof of Theorem \ref{thm:Zp-twisted}.
%
%{More precisely, we need the theorem for the image of the $\Lambda$-adic Kolyvagin system localised at the height one prime $(T+p^m)\subset \Lambda$ for $m\gg 0$. Letting $f_\Lambda$ be a characteristic power series for $\rH^1_{\Fcal_\Lambda}(K,\mathbf{T})/\Lambda\kappa_1$, we have (maybe up to $O(1)$ as $m$ grows) 
%$$\ind(\kappa_{\alpha_m,1},\rH^1_{\Fcal}(K,T_m)) = m \cdot \ord_{(T)}(f_\Lambda)$$
%where $\alpha_m$ is the character with values in $\Z_p^{\times}$ induced by the quotient map $\Lambda\to \Lambda/(T+p^m)$ and $T_m=T_p(E)\otimes \Z_p(\alpha_m)$.
A finite torsion $R$-module $X$ is isomorphic to a sum of cyclic $R$-modules:  $X \simeq \oplus_{i=1}^{s(X)} R/\fm^{d_i}$ for some uniquely-determined
integers $d_i\geq 0$. For an integer $t\geq 0$ we let $\rho_t(X)= \# \{i: d_i > t\}$.  In particular, for $n\in\cN^{(m)}$ we let
\[
\rho_t(n): = \rho_t(\rH^1_{\F(n)}(K, T^{(m)})^+) + \rho_t(\rH^1_{\F(n)}(K,T^{(m)})^-).
\]
Note that it $t<m$ then $\rho_t(n)\geq 1$ since the $\epsilon$ of Proposition \ref{propstructure} is $1$ (the latter fact is implicit in the above rank one result).
We also let
$$
\rho: = 2(\rho_e(1)-1).
$$
Note that $\rho< 2\dim_\mathbb{F}(\rH^1_{\Fcal}(K,T^{(1)}) = 2\dim_{\mathbb F_p}(\rH^1_\Fcal(K,E[p])$, where $\mathbb{F}= R/\fm$ is the residue field of $R$, and hence $\rho$ is bounded by a constant independent of $k$, $m$, and $\alpha$. 


%Let $X$ be a finite torsion $R$-module, we can decompose it as direct sum of cyclic groups as follows $X=\oplus_{i=1}^{r} R/\fm^{m_i}$. We define $\rho_t(X)= \# \{i: m_i \gneq t\}$ and let
%\[
%\rho_t(\rH^1_{\F'}):=\rho_t(\rH^1_{\F'}(K, T_E^{(m)})^+) + \rho_t(\rH^1_{\F'}(K,T_E^{(m)})^-).
%\]
%Let $T:= 2(\rho_e(\rH^1_{\F})-1)$; note $T\leq 2 \dim_{\mathbb{F}}(\rH^1_{\Fcal}(K,T_E^{(1)})$, where $\mathbb{F}$ is the residue field of $R$, and hence $T$ is bounded by a constant independent on $k,m$ and $\alpha$.

\begin{proof}[Proof of \eqref{eq:bound}]
Let $s(n)= \dim_{\mathbb{F}}\rH^1_{\F(n)}(K,T^{(1)})-1 = \dim_\mathbb{F}M(n)[\fm]$, and let
\[
\CE = (\rho + (s(1) +1+ 2\rho)(5\rho+1))e \ \ \text{and} \ \ \CM = (1+5\rho)e.
\]
Note that $\CE$ and $\CM$ are bounded by constants that are independent of $m$ (and depend on $\alpha$ only through the $\Zp$-rank $r$ of $R$). 
Let $k$ be a fixed integer such that 
\begin{equation}\label{eq:kbig}
k > \mathrm{length}_{R}(M(1)) + \ind(\kap_1)+m + (6 s(1)+ 2) e.
\end{equation}
We will show that \eqref{eq:bound} holds provided 
$$
m > \CM.
$$
%We assume $k$ is very large, maybe
%\begin{equation}\label{eq:kbig}
%k > 2\mathrm{length}_{R}(M_0) + \ind(\kap_1)+m + (6 \dim_{\mathbb{F}}(\rH^1_{\Fcal}(K,T_E^{(1)})+ 2) e.
%\end{equation}

%Recall that $\rH^1_{\F(n)}(K,T^{(k)})\simeq R/\fm^{k}\oplus M(n)$, where $M(n) \simeq M_0(n)\oplus M_0(n)$. 
%Then $M(n)$ is isomorphic to a direct sum of cyclic $R$-modules:  $M(n) = \oplus_{i=1}^{s(n)} R/\fm^{d_i(n)}$, and we order
%the exponents so that 
%\[
% d_1(n)=d_2(n)\geq d_3(n)=d_4(n)\geq \dots ...
% \]
%
%We write $H^1_{\Fcal}(K,T^{(k)})\simeq R/\fm^{k}\oplus M(1)$, where $M(1)\simeq M_0\oplus M_0$ as direct sum of cyclic groups and order them as follows:
%\begin{displaymath}
%\rH^1_{\Fcal}(K,T^{(k)})=  R/\fm^{k} \oplus \left(\bigoplus_{i=1}^s R/\fm^{d_i(M(1))} \right)
%\ \ \ \text{where } d_1(M(1))=d_2(M(1))\geq d_3(M(1))=d_4(M(1))\geq \dots
%\end{displaymath}
%Similarly, for any $n\in \mathcal{N}^{(k)}$, we write $H^1_{\Fcal(n)}(K,T^{(k)})\simeq R/\fm^{k}\oplus M(n)$, where $M(n)= M^{(k)}(n)\oplus M^{(k)}(n)$ as in Proposition \ref{propstructure}, since we know $\varepsilon =1$.

If $\rho = 0$, then the exponent of $M(1)$ is at most $e$ and therefore
\[
\tfrac{1}{2} \mathrm{length}_{R}(M(1)) \leq \tfrac{1}{2} (s(1)+1)e \leq \ind(\kappa_1) + \CE.
\]
So we may assume $\rho>0$. 
%
%
%
%If $\rho_e(\rH^1_{\F})=1$, then we have that the exponent of $M(1)$ must be bounded by $e$ and therefore
%\[
%\ind(\kappa_1)+ \tfrac{s}{2} e \geq \tfrac{s}{2} e \geq \tfrac{1}{2}\mathrm{length}_{R}(M(1)).
%\]
%Assume $\rho_e(\rH^1_{\F})>1$. 
%Let us consider $c_0$ a generator of $R/\fm^{k}$. 
%%and $c_1$ a generator of $R/\fm^{d_1(M(1))}$. 
%We look at its image in the $p^m$-torsion of $ \rH^1_{\Fcal}(K,T_m^{(k)})$, namely
%\[
%p^{k-m}c_0=(c_0^+,c_0^-), 
%\]
%%\ \ \ p^{\max(0,d_1(M(1))-m)}c_1=(c_1^+,c_1^-),
%where we denote by $c_i^\pm$ the component in $ \rH^1_{\Fcal}(K,E[p^m])^\pm$. We fix a choice of sign $\varepsilon$ such that $\ord(c_0^\varepsilon)=m$. 

We will find sequences of integers $1=n_0, n_1,..., n_\rho \in \cN^{(k)}$ and $1=x_0,x_1,...,x_\rho$ along with integers $b(n_i)\geq 0$ that satisfy
\begin{itemize}
\item[(a)] $s(n_{i+1})-2 \leq s(n_{i})\leq s(n_{i+1})+2$;
\item[(b)] $\tfrac{1}{2}\mathrm{length}_{R}(M(n_i))\geq  \tfrac{1}{2}\mathrm{length}_{R}(M(n_{i-1}))- b(n_i)$ ;
\item[(c)] $\tfrac{1}{2}\mathrm{length}_{R}(M(n_i)) \leq \tfrac{1}{2}\mathrm{length}_{R}(M(n_{i-1})) + e$;
\item[(d)] $\ord(\kappa_{n_i})\geq \ord(\kappa_{n_{i-1}})-e$;
\item[(e)] $\ind(\kappa_{n_{i-1}})+e\geq \ind(\kappa_{n_i})+b(n_i)$;
\item[(f)]  $x_{i-1}+2\leq x_i\leq x_{i-1} +5$, $\rho_{x_i e}(n_i) \leq \rho_{x_{i-1}e}(n_{i-1})$, and
$\rho_{x_i e}(n_i) = \rho_{x_{i-1}e}(n_{i-1})>1$ only if $\rho_{(x_i+1)e}(\rH^1_{\F(n_{i})}(K,T^{(m)})^\pm)\geq 1$;
\item[(g)] if $\rho_{x_{i-2}e}(n_{i-2})>1$ then 
$\rho_{x_{i} e}(n_{i}) < \rho_{x_{i-2}e}(n_{i-2})$.
\end{itemize}
Before explaining the existence of such sequences, we demonstrate that \eqref{eq:bound} follows from (a)--(g). 


%We want to find a sequence of integers $n_0=1, n_1, \dots, n_T \in \mathcal{N}^{(k)}$, with $T = 2(\rho_e(\rH^1_{\F})-1)$, together with integers $b_{M(n_{t-1})}(n_t)\geq 0$, satisfying the following conditions: 
%\begin{itemize}
%\item[(a)] $s(n_{t+1})-2 \leq s(n_{t})\leq s(n_{t+1})+2$ where $s(n):=\dim_{\mathbb{F}}(M(n)[\fm])$;
%\item[(b)] $\tfrac{1}{2}\mathrm{length}_{R}(M(n_t))\geq  \tfrac{1}{2}\mathrm{length}_{R}(M(n_{t-1}))- b_{M(n_{t-1})}(n_t)$ ;
%\item[(c)] $\mathrm{length}_{R}(M(n_t)) \leqslant \mathrm{length}_{R}(M(n_{t-1})) + 2e$;
%\item[(d)] $\ord(\kappa_{n_t})\geq \ord(\kappa_{n_{t-1}})-e$;
%\item[(e)] $\ind(\kappa_{n_{t-1}})+e\geq \ind(\kappa_{n_t})+b_{M(n_{t-1})}(n_t)$;
%\item[(f)] $\rho_{x_{t-1}e}(\rH^1_{\F(n_{t-1})})\geq \rho_{x_te}(\rH^1_{\F(n_{t})})$, with $x_t=x_{t-1}+2$ or $x_{t-1}+3$ or $x_{t-1}+5$, $x_0 = 1$. Moreover we have an equality if $\rho_{x_{t-1}e}(\rH^1_{\F(n_{t-1})})=1$ and whenever else we have an equality, we can obtain a strict inequality at the next step, namely
%\[
%1<\rho_{x_{t-1}e}(\rH^1_{\F(n_{t-1})})= \rho_{x_te}(\rH^1_{\F(n_{t})}) \Rightarrow \rho_{x_{t}e}(\rH^1_{\F(n_{t})})\gneq \rho_{x_{t+1}e}(\rH^1_{\F(n_{t+1})}).
%\]
%%\item[(g)] letting $c_0$ be a generator of $\Z/p^k$ in $\rH^1_{\Fcal(n_t)}(K,T_m^{(k)})$, then $\ord((1+\varepsilon(n_i)\tau )p^{k-m} c_0)\geq m-ie$.
%\end{itemize}

From repeated appeals to (a) we obtain 
$$
s(n_\rho) \leq s(1) + 2\rho. 
$$
From repeated appeals to (g) we see that 
either $\rho_{x_ie}(n_i) = 1$ for some $1\leq i<\rho$, in which case $1\leq \rho_{x_\rho e}(n_\rho) \leq \rho_{x_ie} (n_i) = 1$, 
or $1\leq \rho_{x_\rho e} (n_\rho) \leq \rho_{x_{\rho-2} e} (n_{\rho-2})-1 \leq \cdots \leq \rho_{x_0 e}(n_0) - \tfrac{1}{2}\rho = \rho_e(1) - \tfrac{1}{2}\rho = 1$.
In either case, we see
$$
\rho_{x_\rho e}(n_\rho) = 1.
$$
As $x_\rho \leq x_0 + 5\rho = 1 +5\rho$, it then follows that 
\begin{equation}\label{eq:length1}
\mathrm{length}_R(M(n_\rho)) \leq s(n_\rho)x_\rho e  \leq s(n_\rho)(5\rho+1) e \leq (s(1) + 2\rho)(5\rho+1) e.
\end{equation}
By repeatedly applying (b) and (e) we obtain
\begin{equation*}\begin{split}
\ind(\kap_1) + \rho e &  \geq \ind(\kap_{n_\rho}) + b(n_1) + b(n_2) + \cdots + b(n_\rho)  \\
&\geq \ind(\kap_{n_\rho}) + \tfrac{1}{2}\mathrm{length}_R(M(1)) - \tfrac{1}{2}\mathrm{length}_R(M(n_\rho)).
\end{split}
\end{equation*}
Combined with \eqref{eq:length1} this gives
\begin{equation*}\begin{split} 
\ind(\kap_1) + \CE & \geq \ind(\kap_{n_\rho}) + \tfrac{1}{2}\mathrm{length}_R(M(1)) - \tfrac{1}{2}\mathrm{length}_R(M(n_\rho)) + (s(1)+2\rho)(5\rho+1)e \\
& \geq \ind(\kap_{n_\rho}) + \tfrac{1}{2}\mathrm{length}_R(M(1)) + \tfrac{1}{2}\mathrm{length}_R(M(n_\rho))  \\
& \geq \tfrac{1}{2}\mathrm{length}_R(M(1)),
\end{split}
\end{equation*}
which is the bound \eqref{eq:bound}.



%From repeated use of (a) we obtain $s+2T\geq s(n_T)$ and from repeated use of (f) we obtain 
%\[
%\rho_{x_Te}(\rH^1_{\F(n_{T})})\leq \rho_{x_{T-2}e}(\rH^1_{\F(n_{T-2})}) -1 \leq \rho_{x_{T-4}e}(\rH^1_{\F(n_{T-4})}) - 2 \leq \dots \leq \rho_e(\rH^1_{\F})-T/2 = 1.
%\] 
%Since $x_T \leq 5T+1$, We deduce $s(n_T)(5T+1)e \geq  \mathrm{length}_{R}(M(n_T))$, therefore
%\begin{equation}\label{eq:finalbound}
%(s+2T)(5T+1)e\geq \mathrm{length}_{R}(M(n_T))
%\end{equation}
%Let $E:=Te+(s+2T)(5T+1)e$. By repeatedly applying (b) and (e), we obtain
%\begin{align*}
%\ind(\kappa_1)+ E&\geq \ind(\kappa_{n_1})+ (T-1)e +(s+2T)(5T+1)e+ \tfrac{1}{2}(\mathrm{length}_{R}(M(1))-\mathrm{length}_{R}(M(n_1))) \\
%&\geq\ind(\kappa_{n_2}) +(T-2)e+(s+2T)(5T+1)e+ \tfrac{1}{2}(\mathrm{length}_{R}(M(1))-\mathrm{length}_{R}(M(n_2)))\geq \dots \\
%&\geq \ind(\kappa_{n_{T-1}}) +e+(s+2T)(5T+1)e + \tfrac{1}{2}(\mathrm{length}_{R}(M(1))-\mathrm{length}_{R}(M(n_{T-1})))\\
%&\geq \ind(\kappa_{n_{T}}) +(s+2T)(5T+1)e +b_{M(n_{T-1})}(n_T)  + \tfrac{1}{2}\mathrm{length}_{R}(M(1))-\tfrac{1}{2}\mathrm{length}_{R}(M(n_{T-1}))\\
%&\geq \ind(\kappa_{n_{T}}) +(s+2T)(5T+1)e +\tfrac{1}{2}\mathrm{length}_{R}(M(1)) -\tfrac{1}{2}\mathrm{length}_{R}(M(n_T))\\
%&\geq \tfrac{1}{2}\mathrm{length}_{R}(M(1))
%\end{align*}
%where we used \eqref{eq:finalbound} for the final inequality.

We will now define the sequence $n_0=1, n_1, \dots, n_T \in \mathscr{N}^{(k)}$ (and subsequently the $b(n_i)$ and $x_i$) by making repeated use of Proposition \ref{prop:prime2} to choose suitable primes in $\mathscr{L}^{(k)}$. 
%\begin{rmk} 
%The strategy to produce such $n_i$ is similar to the one used in \cite[$\S$3.3.3]{eisenstein}, but technically more delicate. In particular comparing the condition (b) here and in \emph{op. cit.}, we notice that the one here is weaker (and would follow from condition (b) in \emph{op. cit.}). The issue is exactly the one hinted at in Remark \ref{rmkcalpha}: we have removed the dependence of the error term on $C_\alpha$, but at the cost of having to work with the $\fm^k$-torsion and, as we will see below, this prevents us from proving that we can take $b_{M(n_{t-1})}(n_t)=d_1(M(n_{t-1}))-e$ as in \cite{eisenstein} and the stronger statement (b) in \emph{op. cit.} This causes the need of the additional condition (f), which can be thought as an induction step on ``the number of summands of the $\fm^m$-torsion of the Selmer groups unbounded by (controlled) multiples of $e$''. Philosophically (?), this is what is done also in the proof of \cite[Lemma 1.6.4]{howard}, where however, being $e=0$, one can work with the $\fm$-torsion and has an equality for the order in Proposition \ref{prop:prime2}.
%\end{rmk}

Suppose $1=n_0, n_1, \dots, n_j \in \mathscr{N}^{(k)}$, $1=x_0, x_1, \dots, x_j$, and $0=b(n_0), b(n_1), \dots, b(n_j)$, $j<\rho$, are such that (a)--(f) hold for all $1\leq i\leq j$ (note that if $j=0$ then (a)--(f) are vacuously true).  
We will explain how to choose a prime $\ell\in \mathscr{L}^{(k)}$ such that $n_0,\dots,n_j,n_{j+1}=n_j\ell$
satisfy (a)--(f) for all $1\leq i\leq j+1$. Repeating this process yields the desired sequence $n_0,\dots,n_\rho$.


Let $c_0\in \rH^1_{\F(n_j)}(K,T^{(k)})$ generate an $R/\fm^{k}$-summand complementary to $M(n_{j})$, so
$\rH^1_{\F(n_j)}(K,T^{(k)}) = R \cdot c_0 \oplus M(n_j) \simeq R/\fm^k\oplus M(n_j)$.
Let $\nu \in \{\pm \}$ such that $\ord((1+\nu\tau)\varpi^{k-m}c_0)=m$, that is, the order of
$c^{\nu}= (1+\nu\tau)(\varpi^{k-m}c_0) \in \rH^1_{\F(n_j)}(K,T^{(m)})^{\nu}$ is $m$. Note that the order of $\varpi^{k-m}c_0$ is $m$, so there exists at least one $\nu \in \{\pm \}$ satisfying the desired condition.
%
%Let $\varepsilon(n) \in \{\pm \}$ such that $\ord((1-\varepsilon(n)\tau)\varpi^{k-m}c_0)=m$.
%To ease notation, let us assume $\varepsilon(n_i)=-$. 
Let $N: = \operatorname{exp} ( \rH^1_{\Fcal(n_j)}(K,T^{(m)})^{-\nu})$ and 
let $c^{-\nu} \in \rH^1_{\Fcal(n_j)}(K,T^{(m)})^{-\nu}$ have order $N$. 
%
%Consider $c^-\in \rH^1_{\Fcal(n_i)}(K,T_E^{(m)})^-$ a class such that $\ord(c^-)=\operatorname{exp} ( \rH^1_{\Fcal(n_i)}(K,T_E^{(m)})^-)=: N$. By definition of $\varepsilon(n_i)$, we have $\ord((1+\tau)\varpi^{k-m}c_0)=m$, for $c_0$ a generator of the $R/\fm^k$ summand.
%%
%%\textcolor{blue}{Assume for the time being that $\ord(c_0^+)=m, \ord(c_1^-)=\min(d_1(M(1)),m)$. Later need to show how to deal with the case everything has highest order component in the same eigenspace}.
%
We apply Proposition \ref{prop:prime2} to the classes $c^{\nu}$ and  $c^{-\nu}$ and obtain a prime $\ell\in \mathscr{L}^{(k)}$ such that
\[
\ord(\loc_{\ell}(\varpi^{k-m}c_0))\geq  \ord(\loc_{\ell}(c^{\nu}))\geq m-e \ \ \text{and} \ \ \ord(\loc_{\ell}(c^{-\nu}))\geq N-e,
\]
and $\loc_{\ell}( \rH^1_{\Fcal(n_j)}(K,T^{(m)}))$ has an $R$-submodule isomorphic to $R/\fm^{m-e}\oplus R/\fm^{N-e}$.
%
%\[
%\loc_{\ell}( \rH^1_{\Fcal(n_i)}(K,T_E^{(m)}))\supset R/\fm^{m-e}\oplus R/\fm^{N-e}.
%\]
As $m>e$, it follows that $\ord(\loc_{\ell}(c_0))\geq k-e$, and
hence $\loc_\ell(\rH^1_{\F(n_j)}(K,T^{(k)})) \simeq R/\fm^{k-a}\oplus R/\fm^{b}$ for some $a\leq e$ and $b\geq N-e$.
In particular, there is a short exact sequence
\begin{equation}\label{eq:imageloc}
0\to H \to \rH^1_{\Fcal(n_j)}(K,T^{(k)}) \xrightarrow{\loc_{\ell}} R/\fm^{k-a}\oplus R/\fm^{b}\to 0, \ \ \ a\leq e, \ b\geq N-e,
\end{equation}
where $H :=\rH^1_{\Fcal(n_j)_{\ell}}(K,T^{(k)}))$ is the kernel of the localisation at $\ell$. 
Global duality then implies that there is an exact sequence
\begin{equation}\label{eq:dualexactseq}
0 \to H \to  \rH^1_{\Fcal(n_j\ell)}(K,T^{(k)})\simeq R/\fm^k \oplus M(n_j\ell)\xrightarrow{\loc_{\ell}} R/\fm^{k-b'}\oplus R/\fm^{a'}\to 0, \ \ \ a\geq a', b'\geq b.
\end{equation}
Here we have used that the arithmetic dual of $T^{(k)} = T^{(k)}_\alpha$ is $T^{(k)}_{\alpha^{-1}}$ and that the complex conjugation $\tau$ induces an isomorphism $\rH^1_{\CF(n)}(K,T^{(k)}_{\alpha^{-1}})\simeq \rH^1_{\CF(n)}(K,T^{(k)}_{\alpha})$.
%We find $d_i(M(n_i\ell))\geq d_{i+1}(K)$ for $i=1, \dots, s(K)$. Combining this with \eqref{eq:d_iM1}, we obtain
%\begin{equation}\label{eq:d_iMell}
%d_i(M(n_i\ell))\geq d_{i+1}(M(n_i)) \text{ for }i=1,\dots j-2, \ \ \ \ \  d_i(M(n_i\ell))\geq d_{i+2}(M(n_i)) \text{ for }i\geq j-1 
%\end{equation}
%This implies that for every $i\not\equiv 0\mod 2$, $i\neq j$ if $j$ is odd $i\neq j-1$ if $j$ even, $1\leq i \leq s-1$, there exists one $d_{t}(M(n_i\ell))$  greater or equal than $d_i(M(n_i))$. Indeed, if $j=1$ or $j=2$, 
%\[
%d_{i-2}(M(n_i\ell))\geq d_{i}(M(n_i)) \ \ \text{ for }i=3,5,\dots, s-1. 
%\]
%If $j\neq 1,2$, the second inequalities take care of all the $d_i(M(n_i))$ for $i=j+2,j+4,\dots,s-1$ if $j$ is odd (or for $j=j+1,j+3,\dots,s-1$ if $j$ is even) and the first inequalities give 
%\[
%d_1(M(n_i\ell))\geq d_{2}(M(n_i))=d_1(M(n_i)), d_2(M(n_i\ell))\geq d_{3}(M(n_i)) , \dots, d_{j-2}(M(n_i\ell))\geq d_{j-1}(M(n_i)).
%\]

We now show (a)-(e) hold for $i=j+1$ with $n_{j+1}=n_j\ell$ and $b(n_{j+1}) = b'$. 

Let $h: = \dim_{\mathbb{F}} H[\fm]$. From \eqref{eq:imageloc} it follows that $h \leq 1+s(n_j) \leq h+2$, and from \eqref{eq:dualexactseq} it follows that
$h\leq 1+s(n_{j+1})\leq h+2$. Hence $s(n_j) -2 \leq h-1 \leq s(n_{j+1}) \leq h+1 \leq s(n_j) +2$, which shows that (a) holds. 

From\eqref{eq:dualexactseq} and  \eqref{eq:imageloc}  we find
\[
\length_R(M(n_{j+1})) = \length_R(H) - b' + a' = \length_R(M(n_j)) - (b+b') + (a+a').
\]
As $(b+b') \leq 2b' = 2b(n_{j+1})$, it follows that (b) holds. And as $(a+a') \leq 2e$, it follows that (c) holds.
%
%
%From the above exact sequences, (a) follows easily and we obtain
%\begin{equation}\label{eq:lengthMell}
%\tag{c}
%\mathrm{length}_{R}(M(n_i\ell)) = \mathrm{length}_{R}(M(n_i)) - (b+b') + (a+a') 
%\leqslant \mathrm{length}_{R}(M(n_i)) + 2e.
%\end{equation}
%and also, letting $n_{i+1}=n_i\ell$ and $b_{M(n_{i})}(n_{i+1})=b'$, we obtain
%\begin{equation}\label{eq:ineqb}
%\tag{b}
%\tfrac{1}{2}\mathrm{length}_{R}(M(n_i\ell))\geq  \tfrac{1}{2}\mathrm{length}_{R}(M(n_i))- b_{M(n_{i})}(n_{i+1}).
%\end{equation}

%From the finite-singular isomorphism and our choice of $\ell$, we obtain $\ord(\loc_{\ell}(\kappa_{\ell}))\geq k-\ind(\kappa_1)-e$. From the conditions on $k$, we must have $\loc_{\ell}(\kappa_{\ell})$ has a non trivial component in $R/\fm^{k-b'}$ and therefore it has index bigger than $b'+\ind(\kappa_\ell)$. We find $\ind(\kappa_1)+e\geq \ind(\kappa_\ell)+b'$ and obtain
%\begin{equation}\label{eq:indexes}
%\tag{e}
%\ind(\kappa_1)+e\geq \ind(\kappa_\ell)+b'.
%\end{equation}

%Combining (c) for $1\leqslant j \leqslant i$ yields 
%$$
%\mathrm{length}_R(M(n_{i})) \leqslant \mathrm{length}_R(M(n_0)) + 2i e.
%$$
%From this, together with $i< T \leq 2(s+1)$, and the assumption (\ref{eq:kbig}), we find
%$$
%k> \mathrm{length}_R(M(n_0)) + (6(s+1)+2)e \geqslant \mathrm{length}_R(M(n_i))+2e.
%$$

To verify that (d) holds for $i=j+1$, we first observe that 
$$
\ord(\kap_{n_{j+1}}) = \ord(\kap_{n_{j}\ell}) \geqslant \ord(\loc_{\ell}(\kap_{n_{j}\ell}))= \ord(\loc_\ell(\kap_{n_{j}})),
$$
the last equality following from the finite-singular relations of the Kolyvagin system.
So (d) holds if $\ord(\loc_\ell(\kap_{n_{j}})) \geqslant \ord(\kap_{n_{j}})-e$.
To see that this last inequality holds, we note that
$\ord(\kap_{n_{j}}) \geq \ord(\kap_{n_0}) - je$ by (d) for $1\leqslant i \leqslant j$. 
But $\ord(\kap_{n_0}) =\ord(\kap_1)= k-\ind(\kap_1)$ by the choice of $k$ (and the fact that $\rH^1_{\CF}(K,T)$ is torsion-free by \eqref{eq:h1}), 
and so by \eqref{eq:kbig} and repeated application of (c) for $1\leq i\leq j$ we have
\begin{equation}\label{eq:ordnj}\begin{split}
\ord(\kap_{n_j}) \geq k-\ind(\kap_1) - je & >  \length_R(M(1))+m+(6s(1)+2-j)e \\
& \geq \length_R(M(n_j))+m+(6s(1)+2-3j)e   \\
& \geq  \length_R(M(n_j)) + m + 2e  \\
& \geq  \length_R(M(n_j)) + 2e.
\end{split}\end{equation}
Write $\kap_{n_{j}} = x c_0 + y$ with $x\in R$ and $y \in M(n_j)$. 
Since $\ord(\kap_{n_j}) > \exp(M(n_j))$ by \eqref{eq:ordnj}, it follows that 
$x = \varpi^t u$ for $t = k-\ord(\kap_{n_{j}})$ and some $u\in R^\times$. 
%Let $ n= \exp(M(n_{i}))$.
It follows that
$$
\varpi^{\exp(M(n_{j}))}\loc_\ell(\kap_{n_j}) = \varpi^{\exp(M(n_{j}))+t} u\loc_\ell(c_0).
$$
By the choice of $\ell$, $\ord(\loc_\ell(c_0)) \geqslant k - e = t+\ord(\kap_{n_j}) - e > t+{\exp(M(n_{j}))}+e$, where the last inequality follows by
\eqref{eq:ordnj}. We then deduce that
$$
\ord(\loc_\ell(\kap_{n_j})) = \ord(\loc_\ell(c_0)) - t \geq k - e -t = \ord(\kap_{n_j})-e,
$$
which -- as noted at the start of this paragraph -- implies that (d) holds.

Next we verify (e) for $i={j+1}$. Let $c_1\in \rH^1_{\CF(n_{j+1})}(K,T^{(k)})$ be a generator of an $R/\fm^k$-summand
complementary to $M(n_{j+1})$, so $\rH^1_{\CF(n_{j+1})}(K,T^{(k)}) = R\cdot c_1\oplus M(n_{j+1})\simeq R/\fm^k\oplus M(n_{j+1})$. 
Write $\kap_{n_j} = u\pi^g c_1 + y$ and $\kap_{n_{j+1}} = v\pi^h c + y'$, where $u,v\in R^\times$, $y\in M(n_{j})$ and $y'\in M(n_{j+1})$. 
Arguing as in the preceding proof that (d) holds for $i={j+1}$ shows that 
$\ord(\kap_{n_i})> \exp(M(n_i))+2e$ for $1\leqslant i\leqslant j+1$ and in particular for $i=j$ and $i=j+1$. 
Hence  $g = k - \ord(\kap_{n_j})$ and $h = k - \ord(\kap_{n_{j+1}})$.
Arguing further as in the proof that (d) holds also yields
$$
\ord(\loc_\ell(\kap_{n_j}))  = \ord(\loc_\ell(c_0))-g \ \ \text{and} \ \  \ord(\loc_\ell(\kap_{n_{j+1}})) = \ord(\loc_\ell(c_1))-h.
$$
From the finite-singular relations for the Kolyvagin system the left-hand sides of both equalities are equal and therefore
$$
h-g = \ord(\loc_\ell(c_1))-\ord(\loc_\ell(c_0)).
$$
We refer again to the short exact sequences \eqref{eq:imageloc} and \eqref{eq:dualexactseq}.
By the choice of $\ell$, 
$\ord(\loc_\ell(c_0))\geq k-e> \exp(M(n_j))\geq b$, the last inequality by \cite[Lem.~3.3.10(ii)]{eisenstein}. Hence we must have 
$\ord(\loc_\ell(c_0))=k-a$. Similarly,  we also must have $\ord(\loc_\ell(c_1))=k-b'$. Thus we find
$$
h-g =  (k-b')-(k-a) = a-b' \leq e -b'.
$$
Since $h-g = \ord(\kap_{n_j})-\ord(\kap_{n_{j+1}})$, this proves $\ord(\kap_{n_j})+b'\leq \ord(\kap_{n_{j+1}})+e$ and hence, since we have shown $\ord(\kap_{n_j})=k-\ind(\kap_{n_j})$ and $\ord(\kap_{n_{j+1}})=k-\ind(\kap_{n_{j+1}})$, (e) holds.

%\begin{rmk}
%Note that, so far, the proof has followed the same lines of \cite[$\S$3.3.3]{eisenstein}. However, we cannot prove, as in \emph{op. cit.}, that 
%$$d_t(M(n_{i+1})) \geqslant d_{t+2}(M(n_{i})), \ \ t=1,\dots,\dim_{\mathbb{F}}(M(n_{i})[\fm])-2$$
%due to the fact that knowing the generator of the localisation of $\rH^1_{\Fcal(n_i)}(K,T^{(k)})[\fm^m]^-$ does not tell us the localisation of which class in $\rH^1_{\Fcal(n_i)}(K,T^{(k)})$ generates the summand $R/\fm^b$ in \eqref{eq:imageloc}. Thus, as shown above, to conclude the proof of \eqref{eq:bound}, we need to prove (f), which roughly  tells us that, even without being able to control the single $d_t(M(n_{i}))$s, the number of unbounded ones decreases. Actually, in general, we can only show that this number does not increase and, if we are in the unfortunate situation where it is stable (which happens essentially if all the ``big summands'' are in the same eigenspace of the big-order multiple of the generator of $R/\fm^k$), then at the next step of induction it decreases. 
%\end{rmk}
%We now prove (f).

So far, our arguments have not wandered far from the lanes of \cite[\S 3.3.3]{eisenstein}.
However, at this point we cannot continue along the same path and deduce -- in the notation of {\em op.~cit.} (see also below) --
that $d_t(M(n_{j+1})) \geqslant d_{t+2}(M(n_{j}))$, $t=1,\dots,\dim_{\mathbb{F}}(M(n_{i})[\fm])-2$. This is because 
knowing the localisation of $\rH^1_{\Fcal(n_j)}(K,T^{(k)})[\fm^m]^{-\nu}$ is not sufficient to determine which class(es) in $\rH^1_{\Fcal(n_j)}(K,T^{(k)})$ 
have localistion generating the summand $R/\fm^b$ in \eqref{eq:imageloc}.  Instead, to conclude the proof of \eqref{eq:bound}, we establish (f) and (g). This roughly  tells us that -- even without being able to control the individual $d_t(M(n_{j}))$s -- the number of large exponents decreases. Actually, in general we can only show that this number does not increase, but if we are in the unfortunate situation where it is stable (which happens essentially if all the ``big summands'' are in the same eigenspace), then at the step $n_{j+2}$ this number decreases. 


It remains to define $x_{j+1}$ and verify (f) and (g).  To this end we introduce some more notation.  A finitely-generated torsion $R$-module $X$ can be written 
as sum of cyclic $R$-modules $X\simeq \oplus_{i=1}^{s(X)} R/\fm^{d_i(X)}$, with the exponents $d_i(X)$ uniquely determined.  We shall always suppose that
the $d_i(X)$ have been labeled so that $d_1(X)\geq d_2(X) \geq \cdots \geq d_{s(X)}(X)$.  Note that $s(X) = \dim_{\mathbb{F}} X[\fm]$.
We will adopt the convention that $d_i(X)= 0$ if $i>s(X)$, extending the $d_i(X)$ to all positive $i$.

%Let $\nu= \nu$. 
Taking the $\fm^m$-torsion of the exact sequence \eqref{eq:imageloc} we obtain two short exact sequences:
\begin{equation}\label{eq:se-1}
0 \to H[\fm^m]^{\nu} \to \rH^1_{\F(n_j)}(K,T^{(m)})^{\nu} \xrightarrow{\loc_\ell} R/\fm^{m-a_{\nu}}\to 0, \ \ 0\leq a_{\nu}\leq e, 
\end{equation}
and
\begin{equation}\label{eq:se-2}
0 \to H[\fm^m]^{-\nu} \to \rH^1_{\F(n_j)}(K,T^{(m)})^{-\nu} \xrightarrow{\loc_\ell} R/\fm^{b_{\nu}}\to 0, \ \ N-e\leq b_{\nu}\leq N,
\end{equation}
The bounds on the exponents for the modules on the right come from the choice of $\ell$ with respect to the classes
$c^{\nu} = (1+\nu\tau)\varpi^{k-m}c_0$ and $c^{-\nu}$.  Global duality then yields two additional short exact sequences
\begin{equation}\label{eq:se-3}
0 \to H[\fm^m]^{\nu} \to \rH^1_{\F(n_{j+1})}(K,T^{(m)})^{\nu} \xrightarrow{\loc_\ell} R/\fm^{a'_{\nu}}\to 0, \ \ 0\leq a_{\nu}'\leq e, 
\end{equation}
and
\begin{equation}\label{eq:se-4}
0 \to H[\fm^m]^{-\nu} \to \rH^1_{\F(n_{j+1})}(K,T^{(m)})^{-\nu} \xrightarrow{\loc_\ell} R/\fm^{m-b'_{\nu}}\to 0, \ \ b_{\nu}'\leq b_{\nu}.
\end{equation}

As $\ord(c^{\nu}) = m$ and $\ord(\loc_\ell(c^{\nu})) = m-a$ with $a\leq e$, it follows from \eqref{eq:se-1}
that $d_i(H[\fm^m]^{\nu}) \leq d_{i+1}(\rH^1_{\F(n_j)}(K,T^{(m)})^{\nu})+e$.  From \eqref{eq:se-3} we deduce
that $d_i(\rH^1_{\F(n_{j+1})}(K,T^{(m)})^{\nu})\leq d_i(H[\fm^m]^{\nu}) +e$, and so
$$
d_i(\rH^1_{\F(n_{j+1})}(K,T^{(m)})^{\nu}) \leq d_{i+1}(\rH^1_{\F(n_j)}(K,T^{(m)})^{\nu})+2e.
$$
%and 
%$$
%d_i(\rH^1_{\F(n_{j+1})}(K,T^{(m)})^{\nu}) \leq 2e, \ \ i> s(H[\fm^m]).
%$$
Let $i_0=\rho_{x_je}(\rH^1_{\F(n_j)}(K, T^{(m)})^{\nu})$. It follows that
$d_{i_0}(\rH^1_{\F(n_{j+1})}(K,T^{(m)})^{\nu}) \leq d_{i_0+1}(\rH^1_{\F(n_j)}(K,T^{(m)})^{\nu})+2e \leq x_j e + 2e$, 
so 
\begin{equation}\label{eq:rhopluspart}
\rho_{(x_j+2)e}(\rH^1_{\F(n_{j+1})}(K, T^{(m)})^{\nu}) \leq i_0-1  = \rho_{x_je}(\rH^1_{\F(n_j)}(K, T^{(m)})^{\nu})-1.
\end{equation}


%Firstly, from the exact sequences obtained taking the $\varpi^m$-torsion of the exact sequences \eqref{eq:imageloc},\eqref{eq:dualexactseq}, we obtain
%\[
%0 \to K[\varpi^m]^+ \to R/\fm^m\cdot c_0^+ \oplus M(n_i)[\varpi^m]^+ \to R/\fm^{m-a}\to 0;
%\]
%\[
%0 \to K[\varpi^m]^- \to \rH^1_{\Fcal(n_i)}(K,T_E^{(m)})^- \to R/\fm^{y}\to 0, \ \ \ y\geq N-e.
%\]
%Since we have shown that $R/\fm^m\cdot c_0^+$ surjects to $ R/\fm^{m-a}$, from the first sequence we obtain
%\[
%0 \to K[\varpi^m]^+\cap R/\fm^m \to  K[\varpi^m]^+ \to M(n_i)[\varpi^m]^+\to 0.
%\]
%Since $K[\varpi^m]^+\cap R/\fm^m \simeq R/\fm^a$ and $a\leq e$, we find $d_i(K[\varpi^m]^+)\leq d_i(M(n_i)[\varpi^m]^+)+e$ for $i=1, \dots, s^+$, where $s^+= \dim_{\mathbb{F}_p}(M(n_i)[\varpi^m]^+)$. 
%If $\dim_{\mathbb{F}_p}(K[\varpi^m]^+)=s^+ +1, d_{s^+ +1}(K[\varpi^m]^+)\leq e$. Using global duality, we find
%\[
%0 \to K[\varpi^m]^+ \to  \rH^1_{\Fcal(n_i\ell)}(K,T_E^{(m)})^+ \to R/\fm^{a'}\to 0, \ \ \ a'\leq e.
%\]
%and we obtain
%\[
%d_i( \rH^1_{\Fcal(n_i\ell)}(K,T_E^{(m)})^+)\leq d_i(K[\varpi^m]^+)+a'\leq d_i(M(n_i)[\varpi^m]^+)+2e, \ \ \ 1\leq i \leq s^+
%\]
%\[
%d_i( \rH^1_{\Fcal(n_i\ell)}(K,T_E^{(m)})^+)\leq 2e, \ \ \ i\geq s^+.
%\]
%Let $i_0$ be such that $d_{i_0}(M(n_i))\leq x_ie$ and $d_{i_0-1}(M(n_i))\gneq x_ie$, so that $\rho_{x_ie}(\rH^1_{\F(n_i)}(K, T_E^{(m)})^+)=i_0$. From the above inequalities, we find $d_{i_0}( \rH^1_{\Fcal(n_i\ell)}(K,T_E^{(m)})^+)\leq d_{i_0}(M(n_i)[\varpi^m]^+)+2e\leq 2e + x_ie$, therefore
%\begin{equation}\label{rhopluspart}
%\rho_{(x_{i}+2)e}(\rH^1_{\F(n_i\ell)}(K, T_E^{(m)})^+)\leq i_0 -1 = \rho_{x_ie}(\rH^1_{\F(n_i)}(K, T_E^{(m)})^+) -1.
%\end{equation}

Next we consider the exact sequence \eqref{eq:se-4}. 
One of the cyclic summands of $\rH^1_{\cF(n_{j+1})}(K,T^{(m)})^{-\nu}$, say one isomorphic to $R/\fm^{d_t}$ for $d_t = d_t(\rH^1_{\cF(n_{j+1})},T^{(m)})^{-\nu})$,
surjects under $\loc_\ell$ onto $R/\fm^{m-b_{\nu}'}$.  There is therefore an $R$-module surjection $H[\fm^m]^{-\nu}\twoheadrightarrow \rH^1_{\cF(n_{j+1})}(K,T^{(m)})^{-\nu}/(R/\fm^{d_t})$, and so -- upon taking Pontryagin duals -- an $R$-module injection
$\rH^1_{\cF(n_{j+1})}(K,T^{(m)})^{-\nu}/(R/\fm^{d_t})\hookrightarrow H[\fm^m]^{-\nu}$. From this together with the injection in \eqref{eq:se-2} we conclude
that
\begin{equation}\label{eq:missingd}
d_i(\rH^1_{\cF(n_{j+1})}(K,,T^{(m)})^{-\nu}) \leq \begin{cases} d_i(\rH^1_{\cF(n_{j})}(K,T^{(m)})^{-\nu}) & i< t \\ 
 d_{i-1}(\rH^1_{\cF(n_{j})}(K,T^{(m)})^{-\nu}) & i>t.\end{cases}
 \end{equation}
 It then follows that
 \begin{equation*}\label{eq:minusineq}
  \rho_{x_je}(\rH^1_{\cF(n_{j+1})}(K,T^{(m)})^{-\nu}) \leq \rho_{x_je}(\rH^1_{\cF(n_{j})}(K,T^{(m)})^{-\nu})+1
 \end{equation*}
% \[
% \rho_{x_je}(\rH^1_{\cF(n_{j+1})}(K,T^{(m)})^{-\nu}) \leq \begin{cases} \rho_{x_je}(\rH^1_{\cF(n_{j})}(K,T^{(m)})^{-\nu}) & \rho_{x_je}(\rH^1_{\cF(n_{j})}(K,T^{(m)})^{-\nu})<t-1\\ \rho_{x_je}(\rH^1_{\cF(n_{j})}(K,T^{(m)})^{-\nu})+1 & \rho_{x_je}(\rH^1_{\cF(n_{j})}(K,T^{(m)})^{-\nu})\geq t-1.
% \end{cases}
% \]
 Together with \eqref{eq:rhopluspart} this implies
 \begin{equation}\label{eq:rhofirstineq}
 \rho_{(x_j+2)e}(n_{j+1}) \leq \rho_{x_je}(n_j),
 \end{equation}
 which proves the first inequality in (f) for any choice of $x_{j+1}\geq x_j+2$.
 %with a strict inequality if $t-1 > \rho_{x_j e}(\rH^1_{\F(n_j)}(K,T^{(m)})^{-\nu})$.
 

%The exact sequence obtained by global duality for the $-$-eigenspaces is 
%\[
%0 \to K[\varpi^m]^- \to  \rH^1_{\Fcal(n_i\ell)}(K,T_E^{(m)})^- \to R/\fm^{m-y'}\to 0, \ \ \ y'\geq y.
%\]
%One of the direct summands, say isomorphic to $R/\fm^x$, of $  \rH^1_{\Fcal(n_i\ell)}(K,T_E^{(m)})^-$ surjects to $R/\fm^{m-y'}$. We hence obtain an inclusion $ \rH^1_{\Fcal(n_i\ell)}(K,T_E^{(m)})^-/R/\fm^x \hookrightarrow K[\varpi^m]^-$, which combined with the exact sequence for $ \rH^1_{\Fcal(n_i)}(K,T_E^{(m)})^-$ gives
%\begin{equation}\label{eq:missingx}
%d_i( \rH^1_{\Fcal(n_i\ell)}(K,T_E^{(m)})^-/R/\fm^x )\leq d_i(K[\varpi^m]^-)\leq d_i( \rH^1_{\Fcal(n_i)}(K,T_E^{(m)})^-).
%\end{equation}
%Note that $d_i( \rH^1_{\Fcal(n_i\ell)}(K,T_E^{(m)})^-/R/\fm^x )$ is either $d_i( \rH^1_{\Fcal(n_i\ell)}(K,T_E^{(m)})^-)$ or $d_{i+1}( \rH^1_{\Fcal(n_i\ell)}(K,T_E^{(m)})^-)$.
%Let $j_0$ be such that $d_{j_0}( \rH^1_{\Fcal(n_i)}(K,T_E^{(m)})^-)\leq x_ie$ and $d_{j_0-1}( \rH^1_{\Fcal(n_i)}(K,T_E^{(m)})^-)\gneq x_ie$, so that we have $\rho_{x_ie}(\rH^1_{\F(n_i)}(K, T_E^{(m)})^+)=j_0-1$. The above inequality implies either $d_{j_0}( \rH^1_{\Fcal(n_i\ell)}(K,T_E^{(m)})^-)\leq x_ie\leq x_{i+1}e$ or $d_{j_0+1}( \rH^1_{\Fcal(n_i\ell)}(K,T_E^{(m)})^-)\leq x_ie$, giving
%\begin{equation}\label{rhominuspartweak}
%\rho_{(x_{i}+2)e}(\rH^1_{\F(n_i\ell)}(K, T_E^{(m)})^-)\leq j_0 = \rho_{x_ie}(\rH^1_{\F(n_i)}(K, T_E^{(m)})^-) +1.
%\end{equation} 

Suppose 
\begin{equation}\label{eq:case1} 
\tag{$\spadesuit$}
x_j e +e < N = \operatorname{exp} (\rH^1_{\F(n_j)}(K,T^{(m)})^{-\nu}).
\end{equation}
Considering the exact sequence \eqref{eq:se-2}, we see that one of the cyclic summands of $\rH^1_{\F(n_j)}(K,T^{(m)})^{-\nu}$, say one isomorphic
to $R/\fm^{d_h}$ for $d_h = d_h(\rH^1_{\F(n_j)}(K,T^{(m)})^{-\nu})$, surjects onto $R/\fm^{b_{\nu}}$. As $b_{\nu}\geq N - e > x_je$, it follows that
$d_h> x_je$, and so $\rho_{x_je}(\rH^1_{\F(n_j)}(K,T^{(m)})^{-\nu}/(R/\fm^{d_h})) = \rho_{x_je}(\rH^1_{\F(n_j)}(K,T^{(m)})^{-\nu}) - 1$.
On the other hand, since $d_h \leq N$, the $R$-module surjection $H[\fm^m]^{-\nu}\twoheadrightarrow \rH^1_{\cF(n_{j})}(K,T^{(m)})^{-\nu}/(R/\fm^{d_h})$ 
has kernel annihilated by $\varpi^e$ and so $\rho_{x_je+e}(H[\fm^m]^{-\nu}) \leq \rho_{x_je}(\rH^1_{\F(n_j)}(K,T^{(m)})^{-\nu}/(R/\fm^{d_h}))$.
Combining these inequalities yields
$$
\rho_{x_je+2e}(H[\fm^m]^{-\nu})\leq \rho_{x_je+e}(H[\fm^m]^{-\nu}) \leq \rho_{x_je}(\rH^1_{\F(n_j)}(K,T^{(m)})^{-\nu}) - 1.
$$
From the previously noted injection $\rH^1_{\cF(n_{j+1})}(K,T^{(m)})^{-\nu}/(R/\fm^{d_t})\hookrightarrow H[\fm^m]^{-\nu}$ we deduce that
$$
\rho_{x_je+2e}(\rH^1_{\cF(n_{j+1})}(K,T^{(m)})^{-\nu}) - 1 \leq \rho_{x_je+2e}(\rH^1_{\cF(n_{j+1})}(K,T^{(m)})^{-\nu}/(R/\fm^{d_t})) \leq
\rho_{x_je+2e}(H[\fm^m]^{-\nu}).
$$
Together the two displayed equations yield
$$
\rho_{x_je+2e}(\rH^1_{\cF(n_{j+1})}(K,T^{(m)})^{-\nu})\leq \rho_{x_je}(\rH^1_{\F(n_j)}(K,T^{(m)})^{-\nu}).
$$
Combining this with \eqref{eq:rhopluspart} we find 
\begin{equation}\label{eq:rhocase1}
\text{\eqref{eq:case1}} \implies \rho_{(x_j+2)e}(n_{j+1}) < \rho_{x_je}(n_j).
\end{equation}



%However, if the following condition holds
%\begin{equation}\label{eq:case1} 
%\tag{$\blacksquare$}
%x_i e +e < \operatorname{exp} (\rH^1_{\Fcal(n_i)}(K,T_E^{(m)})^{\varepsilon(n_i)}),
%\end{equation}
%we can get a stronger result. Namely,
%if $R/\fm^z$ is the summand of $ \rH^1_{\Fcal(n_i)}(K,T_E^{(m)})^-$ surjejcting to $R/\fm^y$, we obtain
%$$z=d_h(\rH^1_{\Fcal(n_i)}(K,T_E^{(m)})^-) \geq N-e > x_ie$$
%and therefore $h\leq j_0 -1$ and $d_{j_0-1}( \rH^1_{\Fcal(n_i)}(K,T_E^{(m)})^-/R/\fm^z ) = d_{j_0}( \rH^1_{\Fcal(n_i)}(K,T_E^{(m)})^-)$.
%Arguing as above, we obtain
%\[
%d_i(K[\varpi^m]^-)\leq d_i( \rH^1_{\Fcal(n_i)}(K,T_E^{(m)})^-/R/\fm^z )+z-y\leq d_i( \rH^1_{\Fcal(n_i)}(K,T_E^{(m)})^-/R/\fm^z )+ e.
%\]
%Hence if $d_{j_0}( \rH^1_{\Fcal(n_i\ell)}(K,T_E^{(m)})^-/R/\fm^x )=d_{j_0}( \rH^1_{\Fcal(n_i\ell)}(K,T_E^{(m)})^-)$, then we have shown above that $$d_{j_0}( \rH^1_{\Fcal(n_i\ell)}(K,T_E^{(m)})^-)\leq x_ie\leq x_{i+1}e.$$ Otherwise we find
%\begin{align*}
%d_{j_0-1}( \rH^1_{\Fcal(n_i\ell)}(K,T_E^{(m)})^-/R/\fm^x )&=\begin{cases}
%d_{j_0}( \rH^1_{\Fcal(n_i\ell)}(K,T_E^{(m)})^-)\\
%d_{j_0-1}( \rH^1_{\Fcal(n_i\ell)}(K,T_E^{(m)})^-)
%\end{cases}\leq d_{j_0-1}(K[\varpi^m]^-)\\&\leq  d_{j_0-1}( \rH^1_{\Fcal(n_i)}(K,T_E^{(m)})^-/R/\fm^z)+e =d_{{j_0}}( \rH^1_{\Fcal(n_i)}(K,T_E^{(m)})^-)+e\\&\leq (x_{i}+2)e.
%\end{align*}
%We have proved that for $x_{i+1}=x_i+2$
%\begin{equation}\label{rhominuspartstrong}
%\rho_{x_{i+1}e}(\rH^1_{\F(\ell)}(K, T_E^{(m)})^-)\leq j_0 -1= \rho_{x_ie}(\rH^1_{\F}(K, T_E^{(m)})^-).
%\end{equation}
%Combining this with \eqref{rhopluspart}, we find that if \eqref{eq:case1} holds then
%\[
%\rho_{x_{i+1}e}(\rH^1_{\F(\ell)})\leq \rho_{x_ie}(\rH^1_{\F})-1, \ \ \ x_{i+1}=x_i+2.
%\]

If \eqref{eq:case1} does not hold, then $N=\operatorname{exp} ( \rH^1_{\Fcal(n_j)}(K,T^{(m)})^{-\nu})\leq (x_j+1)e$.
From \eqref{eq:missingd} we see that each $d_i(\rH^1_{\F(n_{j+1})}(K,T^{(m)})^{-\nu})\leq (x_j+1)e$ except possibly for $i=t$. 
So either $\rho_{(x_j+1)e}(\rH^1_{\F(n_{j+1})}(K,T^{(m)})^{-\nu})=0$ or $t=1$ and $\rho_{(x_j+1)e}(\rH^1_{\F(n_{j+1})}(K,T^{(m)})^{-\nu})=1$.
If $\rho_{(x_j+1)e}(\rH^1_{\F(n_{j+1})}(K,T^{(m)})^{-\nu})=0$, then it follows from \eqref{eq:rhopluspart} that 
$\rho_{(x_j+2)e}(n_{j+1}) < \rho_{x_je}(n_j)$. 
If $d_t \leq (x_j+3)e$, then we also have $\rho_{(x_j+3)e}(n_{j+1}) < \rho_{x_je}(n_j)$ by similar reasoning, where as above $d_t = d_t(\rH^1_{\cF(n_{j+1})},T^{(m)})^{-\nu})$ is the length of the summand
surjecting under $\loc_\ell$ onto $R/\fm^{m-b_{\nu}'}$ in \eqref{eq:se-4}.

Suppose $\rho_{(x_j+3)e}(n_j)=1$. The proof of \eqref{eq:rhofirstineq} also holds for $x_j$ replaced with $x_j+3$, which shows
$\rho_{(x_j+5)e}(n_{j+1}) \leq \rho_{(x_{j}+3)e}(n_{j})=1$.

To summarize, we have shown that 
\begin{equation}\label{eq:equalityconds}
\rho_{(x_j+5)e}(n_{j+1}) = \rho_{x_je}(n_j)>1 \implies
\left\{
\begin{array}{c} \text{\eqref{eq:case1} does not hold}, \\ \rho_{(x_j+1)e}(\rH^1_{\F(n_{j+1})}(K,T^{(m)})^{-\nu})=1, \\  
d_t > (x_j+3)e, \\ \rho_{(x_j+3)e}(n_j)>1.
\end{array} 
\right.
\end{equation}

Suppose $\rho_{(x_j+5)e}(n_{j+1}) = \rho_{x_je}(n_j)>1$.  Since $\rho_{(x_j+5)e}(n_{j+1}) \leq \rho_{(x_j+3)e}(n_{j+1}) \leq \rho_{(x_j+2)e}(n_{j+1}) \leq \rho_{x_je}(n_j)$,
it follows that we also have $\rho_{(x_j+3)e}(n_{j+1})=\rho_{(x_j+2)e}(n_{j+1}) = \rho_{x_je}(n_j)$. 
Furthermore, all the conditions on the right-hand side of
\eqref{eq:equalityconds} hold. As $d_t> (x_j+3)e$, $\mathrm{exp}(\rH^1_{\F(n_{j+1})}(K,T^{(m)})^{-\nu} )> (x_j+3)e$,
so $\rho_{(x_j+3)e}(\rH^1_{\F(n_{j+1})}(K,T^{(m)})^{-\nu} )\geq 1$. 
As $\rho_{(x_j+3)e}(\rH^1_{\F(n_{j+1})}(K,T^{(m)})^{-\nu}) \leq \rho_{(x_j+1)e}(\rH^1_{\F(n_{j+1})}(K,T^{(m)})^{-\nu})=1$, it follows
that $\rho_{(x_j+3)e}(\rH^1_{\F(n_{j+1})}(K,T^{(m)})^{-\nu}) =1$. 
Since $\rho_{x_je}(n_j) \geq 2$,
$\rho_{(x_j+3)e}(\rH^1_{\F(n_{j+1})}(K,T^{(m)})^{\nu}) =  \rho_{(x_j+3)e}(n_{j+1})  - \rho_{(x_j+3)e}(\rH^1_{\F(n_{j+1})}(K,T^{(m)})^{-\nu})  
= \rho_{x_je}(n_j) -1 \geq 1$. This shows
\begin{equation}\label{eq:equalityconds2}
\rho_{(x_j+5)e}(n_{j+1}) = \rho_{x_je}(n_j)>1 \implies
\left\{
\begin{array}{c} \rho_{(x_j+2)e}(n_{j+1}) = \rho_{x_je}(n_j), \\ \rho_{(x_j+3)e}(\rH^1_{\F(n_{j+1})}(K,T^{(m)})^{\pm})\geq 1.
\end{array}\right.
\end{equation}







%If \eqref{eq:case1} is not satisfied, namely $$N=\operatorname{exp} ( \rH^1_{\Fcal(n_i)}(K,T_E^{(m)})^-)\lneq (x_i+1)e$$
%we have \eqref{rhopluspart}. Looking at \eqref{eq:missingx}, we find that all the summands of $\rH^1_{\Fcal(n_i\ell)}(K,T_E^{(m)})^-$ are bounded by $x_{i}e$, but the summand $R/\fm^x$. If also $x\leq x_{i}e+3e$ then we have $\rho_{x_{i}e+3e}(\rH^1_{\F(\ell)}(K, T_E^{(m)})^-)= 0 =\rho_{x_ie}(\rH^1_{\F}(K, T_E^{(m)})^-)$. And therefore
%\[
%\rho_{x_{i}e+3e}(\rH^1_{\F(\ell)})\leq \rho_{x_ie}(\rH^1_{\F})-1.
%\]
%Otherwise, if we cannot prove \eqref{rhominuspartstrong}, but only the weaker \eqref{rhominuspartweak}, we have 
%\[
%\rho_{x_{i}e+2e}(\rH^1_{\F(n_i\ell)})\leq \rho_{x_ie}(\rH^1_{\F(n_i)}).
%\] 

%So it remains to deal with the case $t=1$, $d_t>(x_j+3)e$, and $\rho_{(x_j+1)e}(\rH^1_{\F(n_{j+1})}(K,T^{(m)})^{-\nu})=1$; suppose these hold.
%Suppose also that 
%\begin{equation}\label{eq:extracondition}
%\tag{$\blacktriangle$}
%\rho_{(x_{j}+3)e}(n_{j})>1.
%\end{equation}
%As $d_t> (x_j+3)e$, we have $\mathrm{exp}(\rH^1_{\F(n_{j+1})}(K,T^{(m)})^{-\nu} )> (x_j+3)e$.
%Since $\rho_{(x_j+3)e}(\rH^1_{\F(n_{j+1})}(K,T^{(m)})^{-\nu}) \leq \rho_{(x_j+1)e}(\rH^1_{\F(n_{j+1})}(K,T^{(m)})^{-\nu})=1$, it follows
%that $\rho_{(x_j+3)e}(\rH^1_{\F(n_{j+1})}(K,T^{(m)})^{-\nu}) =1$. 
%Suppose then that $\rho_{(x_j+3)e}(n_{j+1}) = \rho_{x_je}(n_j)$. By \eqref{eq:extracondition}, $\rho_{x_je}(n_j) \geq 2$, hence
%$\rho_{(x_j+3)e}(\rH^1_{\F(n_{j+1})}(K,T^{(m)})^{\nu}) =  \rho_{(x_j+3)e}(n_{j+1})  - \rho_{(x_j+3)e}(\rH^1_{\F(n_{j+1})}(K,T^{(m)})^{-\nu})  
%= \rho_{x_je}(n_j) -1 \geq 1$. This shows,
%\begin{equation}\label{eq:notsq-tri}
%\text{not \eqref{eq:case1} but \eqref{eq:extracondition}  and} \
%\rho_{(x_j+3)e}(n_{j+1}) = \rho_{x_je}(n_j) \implies \rho_{(x_j+3)e}(\rH^1_{\F(n_{j+1})}(K,T^{(m)})^{\pm})\geq 1.
%\end{equation}
%
%Suppose now that \eqref{eq:extracondition} does not hold, so $\rho_{(x_{j}+3)e}(n_{j})=1$.
%The proof of \eqref{eq:rhofirstineq} also holds for $x_j$ replaced with $x_j+3$, with shows
%$\rho_{(x_j+5)e}(n_{j+1}) \leq \rho_{(x_{j}+3)e}(n_{j})=1$.  That is,
%\begin{equation}\label{eq:notri}
%\text{if \eqref{eq:extracondition} does not hold, then} \ \rho_{(x_j+5)e}(n_{j+1})= 1.
%\end{equation}

We can now complete our definition of $x_{j+1}$ and the verification of (f):
\begin{itemize}
\item[(i)] If $\rho_{(x_j+5)e}(n_{j+1}) = 1$ or $\rho_{(x_j+5)e}(n_{j+1}) < \rho_{x_je}(n_j)$, then 
$x_{j+1}:= x_j+5$.
\item[(ii)] If $\rho_{(x_j+5)e}(n_{j+1}) = \rho_{x_je}(n_j)>1$, then $x_{j+1}: = x_j +2$ and 
\eqref{eq:equalityconds2} shows that 
$$\rho_{(x_{j+1}+1)e}(\rH^1_{\F(n_{j+1})}(K,T^{(m)})^{\pm}) \geq 1.$$
\end{itemize}
Hence (f) holds for $i=j+1$.


%
% If $\rho_{x_j e}(n_j) = 1$, then $x_{j+1}: = x_j+2$ and \eqref{eq:rhofirstineq} shows $\rho_{x_{j+1}e}(n_{j+1})=1$.
%\item[(ii)] If $\rho_{x_j e}(n_j)>1$ but \eqref{eq:case1} holds, then $x_{j+1} = x_j+2$ and \eqref{eq:rhocase1} shows
%that $\rho_{x_{j+1}e}(n_{j+1}) < \rho_{x_j e}(n_j)$.
%\item[(iii)] If $\rho_{x_j e}(n_j)>1$, \eqref{eq:case1} holds, but \eqref{eq:extracondition} does not hold, then
%$x_{j+1}: = x_j+3$ and \eqref{eq:notsq-tri} shows that 
%either $\rho_{x_{j+1}e}(n_{j+1}) < \rho_{x_{j+1}e}(n_j)$ or $\rho_{x_{j+1}e}(n_{j+1}) = \rho_{x_{j+1}e}(n_j)$ and
%$\rho_{(x_j+3)e}(\rH^1_{\F(n_{j+1})}(K,T^{(m)})^{\pm})\geq 1$.
%\item[(iv)] If $\rho_{x_j e}(n_j)>1$ and neither \eqref{eq:case1} or \eqref{eq:extracondition} holds, then $x_{j+1}: = x_j+5$
%and \eqref{eq:notri} shows $\rho_{x_{j+1}e}(n_{j+1}) = 1$.
%\end{itemize}
%Hence (f) holds for $i=j+1$.

Finally, we verify (g) for $i=j+1$. Suppose $\rho_{x_{j-1}e}(n_{j-1})>1$. 
If $\rho_{x_j e}(n_j) < \rho_{x_{j-1}e}(n_{j-1})$, then $\rho_{x_{j+1}e}(n_{j+1})\leq \rho_{x_j e}(n_j) < \rho_{x_{j-1}e}(n_{j-1})$. Suppose then that
$\rho_{x_j e}(n_j) = \rho_{x_{j-1}e}(n_{j-1})$. It follows from (f) in the case $i=j$ (which holds by induction) 
that $\rho_{(x_j+1)e}(\rH^1_{\F(n_j)}(K,T^{(m)})^{\pm}) \geq 1$. This implies
that \eqref{eq:case1} holds, and so, by \eqref{eq:equalityconds} above, $\rho_{(x_j+5)e}(n_{j+1}) < \rho_{x_j e}(n_j)= \rho_{x_{j-1}e}(n_{j-1})$.
As $x_{j+1} = x_j+5$ in this case (see (i) above), this shows that (g) holds.
\end{proof}

This completes the proof of Theorem \ref{thm:Zp-twisted}.

\begin{rmk}\label{rmkkolyvaginargument}
The strategy to produce the $n_i\in \mathscr{N}^{(k)}$ as in the preceding proof is similar to the one employed in \cite[$\S$3.3.3]{eisenstein} but technically more delicate. In particular, comparing condition (b) here and condition (b) in \emph{op.\,cit.}, one will notice that the one stated herein is weaker. 
%(and would follow from condition (b) in \emph{op.\,cit.}). 
The issue is exactly the one hinted at previously: we have removed the dependence of the error term on $\alpha$, but at the cost of having to work with the $\fm^m$-torsion. This prevents us from proving that we can take $b_{M(n_{t-1})}(n_t)=d_1(M(n_{t-1}))-e$ (notation as in \cite{eisenstein}) -- and so being able to work with the stronger condition (b). This results in the need for the additional conditions (f) and (g), which can be thought of as an induction step on ``the number of summands of the $\fm^m$-torsion of the Selmer groups that are not bounded by (controlled) multiples of $e$''. Philosophically, this is what is done in the proof of \cite[Lemma 1.6.4]{howard}, where however, as $e=0$, one can work with the $\fm$-torsion and has an equality for the order in Proposition \ref{prop:prime2}. 
\end{rmk}


%
%
%
%
%We first assume
%\begin{equation}\label{eq:extracondition}
%\tag{$\blacktriangle$}
%\rho_{(x_{i}+3)e}(\rH^1_{\F(n_{i})})>1.
%\end{equation}
%We want to show that \eqref{eq:case1} is satisfied for $n_{i}\ell$ and $x_{i+1}=x_i +2$. We write
%\[
% \rH^1_{\Fcal(n_i\ell)}(K,T_m^{(k)})\simeq R/\fm^k\cdot c_{\ell} \oplus M(n_i\ell).
%\]
%If $\ord((1-\tau)p^{k-m}c_{\ell})=m$, namely $\varepsilon(n_{i+1})=+1$, we know $$\operatorname{exp}(\rH^1_{\Fcal(n_i\ell)}(K,T_E^{(m)})^+)\geq \operatorname{exp}(M(n_i)[\varpi^m]^+) \gneq (x_i+3) e,$$ thanks to our assumption \eqref{eq:extracondition}).
%
%If $\ord((1-\tau)p^{k-m}c_{\ell})<m$ and $\ord((1+\tau)p^{k-m}c_{\ell})=m$, namely $\varepsilon(n_{i+1})=-1$ we have $$x=\operatorname{exp}(\rH^1_{\Fcal(n_i\ell)}(K,T_E^{(m)})^-)>x_{i}e+2e+e$$ as required.



%The proofs of 
%
%Finally, the only remaining case is the case where \eqref{eq:extracondition} does not hold, namely 
%$$\rho_{(x_{i}+3)e}(\rH^1_{\F(n_{i})})=1.$$
% The proof of \eqref{rhopluspart} and \eqref{rhominuspartweak} yields similarly 
%\[
%\rho_{x_{i}e+5e}(\rH^1_{\F(n_i\ell)})\leq \rho_{x_ie+3e}(\rH^1_{\F(n_i)}) =1,
%\]
%implying $\rho_{x_{i}e+5e}(\rH^1_{\F(n_i\ell)})=1$.

%To recap, we have proved: 
%\begin{itemize}
%\item if $\rho_{x_{j}e}(n_j)=1$, then for $x_{j+1}=x_j+2$ we also have $\rho_{x_{j+1}e}(n_{j+1})=1$;
%\item if \eqref{eq:case1} holds, then $x_{j+1}=x_j+2$ and $\rho_{x_{i+1}e}(H^1_{\F(n_i\ell)})\lneq \rho_{x_{i}e}(H^1_{\F(n_i)})$;
%\item if \eqref{eq:extracondition} holds and \eqref{eq:case1} doesn't, then either ``we are lucky'' and for $x_{i+1}:=x_i+3$, we find $\rho_{x_{i+1}e}(H^1_{\F(n_i\ell)})\lneq \rho_{x_{i}e}(H^1_{\F(n_i)})$ or $x_{i+1}:=x_i+2$ and $\rho_{x_{i+1}e}(H^1_{\F(n_i\ell)})\leq \rho_{x_{i}e}(H^1_{\F(n_i)})$ but \eqref{eq:case1} holds for $n_i\ell$ (so, by the case above, we are sure we get a strict inequality for $n_{i+2}$);
%\item if \eqref{eq:extracondition} does not hold, then for $x_{i+1}=x_i+5$ and $\rho_{x_{i+1}e}(H^1_{\F(n_i\ell)})=1$.
%\end{itemize}
%
%\end{proof}

\subsection{The anticyclotomic Iwasawa main conjectures}\label{subsec:ac-IMC}

With Theorem \ref{thm:Zp-twisted} in hand, we can state and prove a strengthening of \cite[Thm.~3.4,1]{eisenstein} and consequent strengthenings 
of \cite[Thm.~4.1.2, Thm.~4.2.2]{eisenstein} as well as \cite[Cor.~4.2.3]{eisenstein}. Let $\Lambda=\Lambda_K^-$, 
%
%Let $E$, $p$, $K$ be as in the start $\S\ref{sec:anticyc}$, 
and note that the modules 
\[
\mathcal{X}=\rH^1_{\Fcal_\Lambda}(K,M_E)^\vee,\quad\rH^1_{\Fcal_\Lambda}(K,\mathbf{T})
\] 
in \cite[\S{3.4}]{eisenstein} are the same as the modules $\X_{\rm ord}(E/K_\infty^-)$ and $\mathfrak{S}_{\rm ord}(E/K_\infty^-)$ in $\S\ref{subsec:Sel-str}$, respectively.

\begin{thm}\label{thm:howard}
Assume $E(K)[p]=0$ and suppose there is a Kolyvagin system $\kappa\in\mathbf{KS}(\mathbf{T},\Fcal_\Lambda,\cL_E)$ with $\kappa_1\neq 0$. Then $\mathfrak{S}_{\rm ord}(E/K_\infty^-)$ has $\Lambda$-rank one, and there is a finitely generated torsion $\Lambda$-module $M$ such that
\begin{itemize}
\item[(i)] $\X_{\rm ord}(E/K_\infty^-)\sim\Lambda\oplus M\oplus M$,
\item[(ii)] ${\rm char}_\Lambda(M)$ divides ${\rm char}_\Lambda\bigl(\mathfrak{S}_{\rm ord}(E/K_\infty^-)/\Lambda\kappa_1\bigr)$ in $\Lambda[1/p]$.
\end{itemize}
\end{thm}
\begin{proof}
The proof is the same as that of \cite[Theorem 3.4.1]{eisenstein}. However, the height one prime $(\gamma^- -1)\subset \Lambda$ was excluded from the analysis
in {\em loc.~cit.} because the error term in \cite[Thm.~3.2.1]{eisenstein} increases as the characters $\alpha$ get $p$-adically closer to $1$. 
Replacing the appeal to {\em op.~cit.} with one to Theorem~\ref{thm:Zp-twisted} %(whose error terms in this way) 
for the case of the prime $(\gamma^- -1)$,
yields the theorem.
\end{proof}

Applying Theorem~\ref{thm:howard} to the Kolyvagin system $\kappa^{\rm Hg}=\{\kappa_n^{\rm Hg}\}_{n\in\mathscr{N}}$ of \cite[Thm.~4.1.1]{eisenstein}, we thus obtain the following. 
%(For the ease of comparison, below we keep the notations from \emph{loc.\,cit.}, and note that $\rH^1_{\Fcal_\Lambda}(K,\mathbf{T})$ and $\rH^1(K,M_E)$ are the same as the  Selmer modules $\mathfrak{S}_{\rm ord}(E/K_\infty^-)$ and $\mathfrak{X}_{\rm ord}(E/K_\infty^-)$ introduced in $\S\ref{sec:compact}$.)

\begin{thm}\label{thm:howard-HP}
Assume $E(K)[p]=0$. Then $\mathfrak{S}_{\rm ord}(E/K_\infty^-)$ has $\Lambda$-rank one, and there is a finitely generated torsion $\Lambda$-module $M$ such that
\begin{itemize}
\item[(i)] $\X_{\rm ord}(E/K_\infty^-)\sim\Lambda\oplus M\oplus M$,
\item[(ii)] ${\rm char}_\Lambda(M)$ divides ${\rm char}_\Lambda\bigl(\mathfrak{S}_{\rm ord}(E/K_\infty^-)/\Lambda\kappa_1^{\rm Hg}\bigr)$ in $\Lambda[1/p]$.
\end{itemize}
\end{thm}

Using this, 
%and letting 
%\[
%\X_E=\X_{\rm Gr}(E/K_\infty^-),\quad\Lcal_E=\Lcal_p^{\rm BDP}(f/K),
%\] 
we conclude just as for \cite[Thm.~4.2.2]{eisenstein}:

\begin{thm}\label{thm:BDP-IMC}
Suppose $K$ satisfies hypotheses {\rm (\ref{eq:intro-Heeg})}, {\rm (\ref{eq:intro-spl})}, and {\rm (\ref{eq:intro-disc})}, and that
$E[p]^{ss}=\mathbb{F}_p(\phi)\oplus\mathbb{F}_p(\psi)$ as $G_\bQ$-modules, with $\phi\vert_{G_p}\neq\mathds{1},\omega$.
Then $\X_{\rm Gr}(E/K_\infty^-)$ is $\Lambda$-torsion, and
\[
{\rm char}_\Lambda(\X_{\rm Gr}(E/K_\infty^-))\Lambda^{\rm ur}=(\Lcal_p^{\rm BDP}(f/K))
\]
as ideals in $\Lambda^{\rm ur}$. Hence the anticyclotomic Iwasawa--Greenberg main conjecture in Conjecture~\ref{conj:IMC} holds.
\end{thm}

And just as for \cite[Cor.~4.2.3]{eisenstein} (noting that the ambiguity by powers of $p$ in \emph{loc.\,cit.} can be removed), we then have Theorem~\ref{thm:AC} in the Introduction:

\begin{cor}\label{cor:PR} 
Suppose $K$ satisfies hypotheses {\rm (\ref{eq:intro-Heeg})}, {\rm (\ref{eq:intro-spl})},  {\rm (\ref{eq:intro-disc})}, and that $E[p]^{ss}=\mathbb{F}_p(\phi)\oplus\mathbb{F}_p(\psi)$ as $G_\bQ$-modules, with $\phi\vert_{G_p}\neq\mathds{1},\omega$. 
Then both ${\rm H}^1_{\Fcal_\Lambda}(K,\mathbf{T})$ and 
$\rH^1_{\Fcal_\Lambda}(K,M_E)^\vee$ have $\Lambda$-rank one, and
\[
\mathrm{char}_\Lambda\bigl(\rH^1_{\Fcal_\Lambda}(K,M_E)^\vee_{\rm tors}\bigr)=\mathrm{char}_\Lambda\bigl({\rm H}^1_{\Fcal_\Lambda}(K,\mathbf{T})/\Lambda\kappa_\infty\bigr)^2.
\]
%as ideals in $\Lambda[1/p]$.
\end{cor}
 

\section{Mazur's main conjecture}\label{sec:MC}

In this section we put everything together to deduce the proof of Theorem~\ref{thm:CYC} in the Introduction.

\subsection{The main result}

The following is the main result of this paper.

\begin{thm}\label{thm:A}
Let $E/\Q$ be an elliptic curve, and $p>2$ a prime of good reduction for $E$ such that $E[p]^{ss}=\mathbb{F}_p(\phi)\oplus\mathbb{F}_p(\psi)$ 
with $\phi\vert_{G_{\Q_p}}\neq 1,\omega$. Then the module $\X_{\rm ord}(E/\Q_\infty)$ is $\Lambda_\Q$-torsion, with
\[
\ch_{\Lambda_\Q}\bigl(\X_{\rm ord}(E/\Q_\infty)\bigr)=\bigl(\Lcal_p^{\rm MSD}(E/\Q)\bigr).
\]
In other words, Mazur's main conjecture for $E$ holds.
\end{thm}

\subsection{Proof of Mazur's main conjecture}

We divide it into three steps, similarly as we did in $\S\ref{sec:one}$. Choose an imaginary quadratic field $K$ satisfying hypotheses (\ref{eq:Heeg}), (\ref{eq:spl}), and (\ref{eq:disc}). As usual, we let $E_\bullet$ denote the elliptic curve in the isogeny class of $E$ constructed in \cite{wuthrich-int}, and put $S=\Sigma\smallsetminus\{p,\infty\}$.
\sk

\noindent\emph{Step 1}. The $p$-adic $L$-functions $\Lcal_p^{\rm PR}(E_\bullet/K)^+\in\Lambda_K^+$ and $\Lcal_p^{\rm BDP}(f/K)\in\Lambda_K^{-,{\rm ur}}$ are nonzero: For $\Lcal_p^{\rm PR}(E_\bullet/K)^+$ this follows from  Proposition~\ref{prop:comp-Lcyc}, Theorem~\ref{thm:MSD}, and Rohrlich's nonvanishing result \cite{rohrlich-cyc}; and for $\Lcal_p^{\rm BDP}(f/K)$ this is part of Theorem~\ref{thm:BDP}. Fix an integer $m>0$ such that 
\begin{equation}\label{eq:p^m}
\Lcal_p^{\rm PR}(E_\bullet/K)^+\neq 0\in\Lambda_K^+/p^m\Lambda_K^+,
\end{equation}
and take a crystalline character $\alpha:\Gamma_K^-\rightarrow R^\times$ with $\alpha\equiv 1\pmod{\varpi^m}$ such that  $\Lcal_p^{\rm BDP}(f(\alpha)/K)(0)\neq 0$. 

By Theorem~\ref{thm:BDP-IMC} %and Proposition~\ref{prop:prim-imprim} 
we then have that $\X_{\rm Gr}(E_\bullet(\alpha)/K_\infty^-)$ is $\Lambda_K^-$-torsion, with
\begin{equation}\label{eq:ac-alpha}
\Fcal_{\rm Gr}(E_\bullet(\alpha)/K_\infty^-)(0)\;\sim_p\;\Lcal_p^{\rm BDP}(f(\alpha)/K)(0)\neq 0,%\tag{ac-$\alpha$}
\end{equation}
where $\Fcal_{\rm Gr}(E_\bullet/K_\infty^-)\in R\llbracket{T}\rrbracket$ is any characteristic power series for $\X_{\rm Gr}(E_\bullet(\alpha)/K_\infty^-)$.
\sk

\noindent\emph{Step 2}. Since $\alpha$ is anticyclotomic, for $\alpha\neq 1$ the nonvanishing of $\Lcal_p^{\rm PR}(E_\bullet(\alpha)/K)^+$ and $\Lcal_p^{\rm Gr}(f(\alpha)/K)^+$ is not automatic, but for our choice of $\alpha$ we can show this easily.

\begin{lemma}\label{lem:nonzeroLs}
With $\alpha$ chosen as above, the $p$-adic $L$-functions $\Lcal_p^{\rm PR}(E_\bullet(\alpha)/K)^+$ and $\Lcal_p^{\rm Gr}(f(\alpha)/K)^+$ are both nonzero.
\end{lemma}

\begin{proof}
For $\Lcal_p^{\rm PR}(E_\bullet(\alpha)/K)^+$, this is clear from (\ref{eq:p^m}) and the congruence of Lemma~\ref{lem:cong-L}; and for $\Lcal_p^{\rm Gr}(f(\alpha)/K)^+$ it follows from the relations
\begin{equation}\label{eq:sim-p}
\Lcal_p^{\rm Gr}(f(\alpha)/K)^+(0)=\Lcal_p^{\rm Gr}(f(\alpha)/K)^-(0)\,\sim_p\,\Lcal_p^{\rm BDP}(f(\alpha)/K)(0),
\end{equation}
using Proposition~\ref{prop:comp-Lac} for the last equality up to a $p$-adic unit.
\end{proof}

%In light of Lemma~\ref{lem:nonzeroLs}
In light of this nonvanishing, by Corollary~\ref{cor:KLZ} the class $BF_{\alpha}^+\in\mathfrak{S}_{\rm ord,\rm rel}(E_\bullet(\alpha)/K_\infty^+)$ is nonzero, and so by Theorem~\ref{thm:BF-ES} the Selmer group $\X_{\rm ord,\rm str}(E_\bullet(\alpha^{-1})/K_\infty^+)$ is $\Lambda_K^+$-torsion, with
\begin{equation}\label{eq:BF-div-alpha}
\ch_{\Lambda_K^+}\bigl(\X_{\rm ord,\rm str}(E_\bullet(\alpha^{-1})/K_\infty^+)\bigr)\supset\ch_{\Lambda_K^+}\bigl(\mathfrak{S}_{\rm ord,\rm rel}(E_\bullet(\alpha)/K_\infty^+)/\Lambda_K^+BF_{\alpha}^+\bigr)
\end{equation}
in $\Lambda_K^+\otimes\Q_p$. Note the need to invert $p$ in this divisibility, an ambiguity that we shall remove in the next result.

%Despite having to invert $p$ in this divisibility, the divisibility in the next result can be shown to be integral (a key point later in the argument).

\begin{lemma}\label{lem:int-div-PR}
The module $\X_{\rm ord}(E_\bullet(\alpha)/K_\infty^+)$ is $\Lambda_K^+$-torsion, with
\begin{equation}\label{eq:cyc-ord-alpha}
\ch_{\Lambda_K^+}\bigl(\X_{\rm ord}(E_\bullet(\alpha)/K_\infty^+)\bigr)\supset\bigl(\Lcal_p^{\rm PR}(E_\bullet(\alpha)/K)^+\bigr) %\tag{cyc-ord-$\alpha$}
\end{equation}
in $\Lambda_K^+$.
\end{lemma}

\begin{proof}
The combination of Proposition~\ref{prop:equiv}, Lemma~\ref{lem:nonzeroLs}, and (\ref{eq:BF-div-alpha}) shows that $\X_{\rm ord}(E_\bullet(\alpha)/K_\infty^+)$ is $\Lambda_K^+$-torsion, with the claimed divisibility holding in $\Lambda_K^+\otimes\Q_p$. Let $S=\Sigma\smallsetminus\{p,\infty\}$. By Corollary~\ref{cor:imp-IMC}, it follows that $\mathfrak{X}_{\rm ord}^S(E_\bullet(\alpha)/K_\infty^+)$ is also $\Lambda_K^+$-torsion, with
\begin{equation}\label{eq:ord-cyc-Qp}
\ch_{\Lambda_K^+}\bigl(\X_{\rm ord}^S(E_\bullet(\alpha)/K_\infty^+)\bigr)\supset\bigl(\Lcal_p^{\rm PR}(E_\bullet(\alpha)/K)^{+,S}\bigr) %\tag{cyc-ord-$\alpha$}
\end{equation}
in $\Lambda_K^+\otimes\Q_p$. Denote by $\Fcal_{\rm ord}^S(E_\bullet(\alpha)/K_\infty^+)\in\Lambda_K^+$ a characteristic power series for $\X_{\rm ord}^S(E_\bullet(\alpha)/K_\infty^+)$, so from the above we have
\begin{equation}\label{eq:div-h}
\Fcal_{\rm ord}^S(E_\bullet(\alpha)/K_\infty^+)\cdot h=\varpi^k\cdot\Lcal_p^{\rm PR}(E_\bullet(\alpha)/K)^{+,S}
\end{equation}
for some $h\in\Lambda_K^+$ and $k\in\Z$. If $k<0$ there is nothing to show, so assume $k\geq 0$. From Kato's divisibility \cite[Thm.~17.4]{kato-euler-systems} (refined to an integral statement as in \cite[Thm.~16]{wuthrich-int}), Proposition~\ref{prop:comp-Lcyc}, and Proposition~\ref{prop:comp-Selcyc} we have that the untwisted Selmer group $\X_{\rm ord}(E_\bullet/K_\infty^+)$ is $\Lambda_K^+$-torsion, with the integral divisibility
\begin{equation}\label{eq:final-cyc-ord-1}
\ch_{\Lambda_K^+}\bigl(\X_{\rm ord}(E_\bullet/K_\infty^+)\bigr)\supset\bigl(\Lcal_p^{\rm PR}(E_\bullet/K)^+\bigr)
%\quad\textrm{in $\Lambda_K^+$}
%\tag{cyc-ord-1}
\end{equation}
in $\Lambda_K^+$, and so from Proposition~\ref{prop:prim-imprim-ord} we get that $\X_{\rm ord}^S(E_\bullet/K_\infty^+)$ is also $\Lambda_K^+$-torsion, and we have the integral divisibility
\begin{equation}\label{eq:final-cyc-ord-1-imp}
\ch_{\Lambda_K^+}\bigl(\X_{\rm ord}^S(E_\bullet/K_\infty^+)\bigr)\supset\bigl(\Lcal_p^{\rm PR}(E_\bullet/K)^{+,S}\bigr)
%\quad\textrm{in $\Lambda_K^+$}
%\tag{cyc-ord-1}
\end{equation}
in $\Lambda_K^+$. By the congruences of Proposition~\ref{prop:cong-Sel} and Lemma~\ref{lem:cong-L}, it follows from (\ref{eq:final-cyc-ord-1-imp}) that for $\alpha$ sufficiently close to $1$ (i.e. taking $m>0$ in \eqref{eq:p^m} sufficiently large) the $\mu$-invariant of $\Fcal_{\rm ord}^S(E_\bullet(\alpha)/K_\infty^+)$ is at most that of $\Lcal_p^{\rm PR}(E_\bullet(\alpha)/K)^{+,S}$. Thus from (\ref{eq:div-h}) we see that $h$ is divisible by $\varpi^k$, and therefore the divisibility (\ref{eq:ord-cyc-Qp}) holds in $\Lambda_K^+$.  Together with Corollary~\ref{cor:imp-IMC}, this yields the result.
\end{proof}

From the preceding two lemmas and Proposition~\ref{prop:equiv}, we deduce that $\X_{\rm Gr}(E_\bullet(\alpha)/K_\infty^+)$ is $\Lambda_K^+$-torsion, with
\begin{equation}\label{eq:cyc-alpha-final}
\bigl(\Fcal_{\rm Gr}(E_\bullet(\alpha)/K_\infty^+)\bigr)\supset\bigl(\Lcal_p^{\rm Gr}(f(\alpha)/K)^+\bigr)
%\tag{cyc-$\alpha$}
\end{equation}
in $\Lambda_K^{+,\rm ur}$, where $\Fcal_{\rm Gr}(E_\bullet(\alpha)/K_\infty^+)\in\Lambda_K^+$ is any characteristic power series for $\X_{\rm Gr}(E_\bullet(\alpha)/K_\infty^+)$. Together with the relations
\[
\Fcal_{\Gr}(E_\bullet(\alpha)/K_\infty^+)(0)\,\sim_p\,
\Fcal_{\Gr}(E_\bullet(\alpha)/K_\infty^-)(0)\,\sim_p\,\Lcal_p^{\rm BDP}(f(\alpha)/K)(0)\,\sim_p\,\Lcal_p^{\rm Gr}(f(\alpha)/K)^-(0)\neq 0
\]
following from Proposition~\ref{prop:euler-char}, relation (\ref{eq:ac-alpha}), and Proposition~\ref{prop:comp-Lac}, and noting that 
\[
\Lcal_p^{\rm Gr}(f(\alpha)/K)^-(0)=\Lcal_p^{\rm Gr}(f(\alpha)/K)^+(0),
\]
by easy commutative algebra (see \cite[Lem.\,3.2]{skinner-urban}) it follows that equality holds in (\ref{eq:cyc-alpha-final}), so we have
\[
\bigl(\Fcal_{\rm Gr}(E_\bullet(\alpha)/K_\infty^+)\bigr)=\bigl(\Lcal_p^{\rm Gr}(f(\alpha)/K)^+\bigr).
\]
(Note that conditions (a) and (b) in Lemma~\ref{lem:coinv} needed for the above application of Proposition~\ref{prop:euler-char} follow from our choice of $\alpha$ with $\Lcal_p^{\rm BDP}(f(\alpha)/K)(0)\neq 0$ together with the reciprocity law of \cite[Thm.\,5.7]{cas-hsieh1} and the result of \cite[Thm.\,3.2.1]{eisenstein} applied to the Heegner point Kolyvagin system in \cite[\S{4.1}]{eisenstein}.) 

In particular, for $\alpha$ as above sufficiently close to $1$, we conclude by Proposition~\ref{prop:equiv} that $\X_{\rm ord}(E_\bullet(\alpha)/K_\infty^+)$ is $\Lambda_K^+$-torsion, with
\begin{equation}\label{eq:alpha-equality}
\ch_{\Lambda_K^+}\bigl(\X_{\rm ord}(E_\bullet(\alpha)/K_\infty^+)\bigr)=\bigl(\Lcal_p^{\rm PR}(E_\bullet(\alpha)/K)^+\bigr).
\end{equation}
%\sk
%On the other hand, from Kato's divisibility \cite[Thm.~17.4]{kato-euler-systems} (as refined in \cite[Thm.~16]{wuthrich-int}), Proposition~\ref{prop:comp-Lcyc}, and Corollary~\ref{cor:comp-Selcyc} we have that the untwisted Selmer group $\X_{\rm ord}(E_\bullet/K_\infty^+)$ is $\Lambda_K^+$-torsion, with
%\begin{equation}\label{eq:cyc-ord-1}
%\ch_{\Lambda_K^+}\bigl(\X_{\rm ord}(E_\bullet/K_\infty^+)\bigr)\supset\bigl(\Lcal_p^{\rm PR}(E_\bullet/K)^+\bigr)\quad\textrm{in $\Lambda_K^+$. 
%}
%\tag{cyc-ord-1}
%\end{equation}

\noindent\emph{Step 3}. 
%Now it remains to show that equality holds in (\ref{eq:cyc-ord-1}), which we shall do by first establishing equality in (\ref{eq:cyc-ord-alpha}). 
%(Note that these divisibilities take place in the cyclotomic Iwasawa algebra $\Lambda_K^+$, whereas $\alpha$ is anticyclotomic.)
%
%Note that the divisibilities (\ref{eq:cyc-ord-alpha}) and (\ref{eq:cyc-ord-1}) are not directly related, since they take place in the cyclotomic Iwasawa algebra $\Lambda_K^+$ where $\alpha$ is anticyclotomic. However, we shall be able between the two by a congruence argument.
%
%\begin{prop}\label{prop:Gr-alpha}
%The module $\X_{\rm ord}(E_\bullet(\alpha)/K_\infty^+)$ is $\Lambda_K^+$-torsion, with 
%\[
%\ch_{\Lambda_K^+}\bigl(\X_{\rm ord}(E_\bullet(\alpha)/K_\infty^+)\bigr)=%\bigl(\Lcal_p^{\rm PR}(E_\bullet(\alpha)/K)^+\bigr).%
%\]
%\end{prop}
%
%\begin{proof}
%That $\X_{\rm ord}(E_\bullet(\alpha)/K_\infty^+)$ (and also $\X_{\rm Gr}(E_\bullet(\alpha)/K_\infty^+)$) is $\Lambda_K^+$-torsion was already shown above. On the other hand, 
%from the relations (\ref{eq:sim-p}) and (\ref{eq:ac-alpha}), and Propositions~\ref{prop:prim-imprim} and \ref{prop:euler-char} we obtain
%\[
%\Fcal_{\rm Gr}(E_\bullet(\alpha)/K_\infty^+)(0)
%\,\sim_p\,\Lcal_p^{\rm Gr}(f(\alpha)/K)^+(0)\neq 0,
%\]
%By easy commutative algebra (see \cite[Lem.~3.2]{skinner-urban}), it follows that equality holds in (\ref{eq:cyc-alpha-final}),
%so we have
%\[
%\bigl(\Fcal_{\rm Gr}(E_\bullet(\alpha)/K_\infty^+)\bigr)=\bigl(\Lcal_p^{\rm Gr}(f(\alpha)/K)^+\bigr).
%\]
%The result now follows from Proposition~\ref{prop:equiv} and Lemma~\ref{lem:nonzeroLs}.
%\end{proof}
%
We are now in a position to prove Conjecture~\ref{conj:IMC-K} for $\X_{\rm ord}(E/K_\infty^+)$.

\begin{thm}\label{thm:PR-IMC}
Let $E/\Q$ be an elliptic curve, and $p>2$ a prime of good reduction for $E$ such that $E[p]^{ss}=\mathbb{F}_p(\phi)\oplus\mathbb{F}_p(\psi)$ 
with $\phi\vert_{G_{\Q_p}}\neq 1,\omega$. 
Then module $\X_{\rm ord}(E/K_\infty^+)$ is $\Lambda_K^+$-torsion, with
\[
\ch_{\Lambda_K^+}(\X_{\rm ord}(E/K_\infty^+))=(\Lcal_p^{\rm PR}(E/K)^+).
\]
%as ideals in $\Lambda_K^+$.
\end{thm}

\begin{proof}
By Proposition~\ref{prop:isog-inv}, it suffices to prove the result of $E_\bullet$. Let $S=\Sigma\smallsetminus\{p,\infty\}$. As shown above, from Kato's work we can deduce that $\X_{\rm ord}^S(E_\bullet/K_\infty^+)$ is $\Lambda_K^+$-torsion, and we have the integral divisibility
\[
\ch_{\Lambda_K^+}\bigl(\X_{\rm ord}^S(E_\bullet/K_\infty^+)\bigr)\supset\bigl(\Lcal_p^{\rm PR}(E_\bullet/K)^{+,S}\bigr)
\]
in $\Lambda_K^+$. Take a character $\alpha:\Gamma_K^-\rightarrow R^\times$ with 
\[
\alpha\equiv 1\;({\rm mod}\,\varpi^m)
\] 
for some $m\gg 0$ so that the equality (\ref{eq:alpha-equality}) holds. (Note that the argument in \emph{Step 1} and \emph{Step 2} leading to that equality only excludes finitely many $\alpha$.) By Corollary~\ref{cor:imp-IMC} it follows that $\X_{\rm ord}^S(E_\bullet/K_\infty^+)$ is also $\Lambda_K^+$-torsion, and denoting by $\Fcal_{\rm ord}^S(E_\bullet/K_\infty^+)$ and $\Fcal^S_{\rm ord}(E_\bullet(\alpha)/K_\infty^+)\in\Lambda_K^+$ characteristic power series for $\X^S_{\rm ord}(E_\bullet/K_\infty^+)$ and $\X^S_{\rm ord}(E_\bullet(\alpha)/K_\infty^+)$, respectively, we have 
\begin{equation}\label{eq:key-cong}
\Fcal^S_{\rm ord}(E_\bullet/K_\infty^+)\equiv\Fcal^S_{\rm ord}(E_\bullet(\alpha)/K_\infty^+)\equiv\Lcal_p^{\rm PR}(E_\bullet(\alpha)/K)^{+,S}\equiv\Lcal_p^{\rm PR}(E_\bullet/K)^{+,S}\;({\rm mod}\,\varpi^m),
\end{equation}
as a consequence of Proposition~\ref{prop:cong-Sel}, the combination of (\ref{eq:alpha-equality}) and Corollary~\ref{cor:imp-IMC}, and Lemma~\ref{lem:cong-L}, respectively. Taking $m\gg 0$, it follows from the congruence (\ref{eq:key-cong}) that $\Fcal_{\rm ord}^S(E_\bullet/K_\infty^+)$ and $\Lcal_p^{\rm PR}(E_\bullet/K)^{+,S}$ have \emph{the same} Iwasawa $\lambda$- and $\mu$-invariants, and so equality holds in (\ref{eq:final-cyc-ord-1-imp}). By Corollary~\ref{cor:imp-IMC}, this yields the proof of the theorem.
\end{proof}

The proof of Mazur's main conjecture for $E$ now follows easily.

\begin{proof}[Proof of Theorem~\ref{thm:A}] 
%With Theorem~\ref{thm:PR-IMC} at hand, the result is now immediate.  
Choose an imaginary quadratic field $K$ satisfying hypotheses (\ref{eq:disc}), (\ref{eq:Heeg}), and (\ref{eq:spl}). As before, from \cite{kato-euler-systems} and \cite{wuthrich-int} we have the divisibilities
\[
\ch_{\Lambda_\Q}\bigl(\X_{\rm ord}(E/\Q_\infty)\bigr)\supset\bigl(\Lcal_p^{\rm MSD}(E/\Q)\bigr),\quad\ch_{\Lambda_\Q}\bigl(\X_{\rm ord}(E^K/\Q_\infty)\bigr)\supset\bigl(\Lcal_p^{\rm MSD}(E^K/\Q)\bigr)
\]
in $\Lambda_\Q$. If $\ch_{\Lambda_\Q}(\X_{\rm ord}(E/\Q_\infty))\neq(\Lcal_p^{\rm MSD}(E/\Q))$ then from Proposition~\ref{prop:comp-Lcyc} and Proposition~\ref{prop:comp-Selcyc} we conclude that 
\[
\bigl(\Lcal_p^{\rm PR}(E/K)^+\bigr)\subsetneq
\ch_{\Lambda_K^+}(\X_{\rm ord}(E/K_\infty^+)\bigr),
\]
but this contradicts Theorem~\ref{thm:PR-IMC}. Thus $\ch_{\Lambda_\Q}(\X_{\rm ord}(E/\Q_\infty))=(\Lcal_p^{\rm MSD}(E/\Q))$, concluding the proof.
\end{proof}


\subsection*{Conflict of interest statement} On behalf of all authors, the corresponding author states that there is no conflict of interest and that the manuscript complies to the Ethical Rules applicable for this journal.

\bibliography{references}
\bibliographystyle{alpha}

\end{document}
